\input{preamble}

% OK, start here.
%
\begin{document}

\title{\'Etale Cohomology}


\maketitle

\phantomsection
\label{section-phantom}

\tableofcontents


\section{Introduction}
\label{section-introduction}

\noindent
This chapter is the first in a series of chapter on the \'etale cohomology
of schemes. In this chapter we discuss the very basics of the \'etale
topology and cohomology of abelian sheaves in this topology. Many of the
topics discussed may be safely skipped on a first reading; please see
the advice in the next section as to how to decide what to skip.

\medskip\noindent
The initial version of this chapter was formed by the notes
of the first part of a course on \'etale cohomology taught by Johan de Jong
at Columbia University in the Fall of 2009. The original note takers were
Thibaut Pugin, Zachary Maddock and Min Lee. The second part of the
course can be found in the chapter on the trace formula, see
The Trace Formula, Section \ref{trace-section-introduction}.




\section{Which sections to skip on a first reading?}
\label{section-skip}

\noindent
We want to use the material in this chapter for the development of
theory related to algebraic spaces, Deligne-Mumford stacks, algebraic stacks,
etc. Thus we have added some pretty technical material to the original
exposition of \'etale cohomology for schemes. The reader can recognize this
material by the frequency of the word ``topos'', or by discussions related
to set theory, or by proofs dealing with very general properties of morphisms
of schemes. Some of these discussions can be skipped on a first reading.

\medskip\noindent
In particular, we suggest that the reader skip the following sections:
\begin{enumerate}
\item Comparing big and small topoi,
Section \ref{section-compare}.
\item Recovering morphisms,
Section \ref{section-morphisms}.
\item Push and pull,
Section \ref{section-monomorphisms}.
\item Property (A),
Section \ref{section-A}.
\item Property (B),
Section \ref{section-B}.
\item Property (C),
Section \ref{section-C}.
\item Topological invariance of the small \'etale site,
Section \ref{section-topological-invariance}.
\item Integral universally injective morphisms,
Section \ref{section-integral-universally-injective}.
\item Big sites and pushforward,
Section \ref{section-big}.
\item Exactness of big lower shriek,
Section \ref{section-exactness-lower-shriek}.
\end{enumerate}
Besides these sections there are some sporadic results that may be skipped
that the reader can recognize by the keywords given above.



%9.08.09
\section{Prologue}
\label{section-prologue}

\noindent
These lectures are about another cohomology theory. The first thing to remark
is that the Zariski topology is not entirely satisfactory. One of the main
reasons that it fails to give the results that we would want is that if $X$ is
a complex variety and $\mathcal{F}$ is a constant sheaf then
$$
H^i(X, \mathcal{F}) = 0, \quad \text{ for all } i > 0.
$$
The reason for that is the following. In an irreducible scheme (a variety in
particular), any two nonempty open subsets meet, and so the restriction
mappings of a constant sheaf are surjective. We say that the sheaf is
{\it flasque}. In this case, all higher {\v C}ech cohomology groups vanish, and
so do all higher Zariski cohomology groups. In other words, there are ``not
enough'' open sets in the Zariski topology to detect this higher cohomology.

\medskip\noindent
On the other hand, if $X$ is a smooth projective complex variety, then
$$
H_{Betti}^{2 \dim X}(X (\mathbf{C}), \Lambda) = \Lambda \quad \text{ for }
\Lambda = \mathbf{Z}, \ \mathbf{Z}/n\mathbf{Z},
$$
where $X(\mathbf{C})$ means the set of complex points of $X$. This is a feature
that would be nice to replicate in algebraic geometry. In positive
characteristic in particular.




\section{The \'etale topology}
\label{section-etale-topology}

\noindent
It is very hard to simply ``add'' extra open sets to refine the Zariski
topology. One efficient way to define a topology is to consider not only open
sets, but also some schemes that lie over them. To define the \'etale topology,
one considers all morphisms $\varphi : U \to X$ which are \'etale. If
$X$ is a smooth projective variety over $\mathbf{C}$, then this means
\begin{enumerate}
\item $U$ is a disjoint union of smooth varieties, and
\item $\varphi$ is (analytically) locally an isomorphism.
\end{enumerate}
The word ``analytically'' refers to the usual (transcendental) topology over
$\mathbf{C}$. So the second condition means that the derivative of $\varphi$
has full rank everywhere (and in particular all the components of $U$
have the same dimension as $X$).

\medskip\noindent
A double cover -- loosely defined as a finite degree $2$ map between varieties
-- for example
$$
\Spec(\mathbf{C}[t])
\longrightarrow
\Spec(\mathbf{C}[t]),
\quad t \longmapsto t^2
$$
will not be an \'etale morphism if it has a fibre consisting of a single point.
In the example this happens when $t = 0$. For a finite map between varieties
over $\mathbf{C}$ to be \'etale all the fibers should have the same number of
points. Removing the point $t = 0$ from the source of the map in the example
will make the morphism \'etale. But we can remove other points from the source
of the morphism also, and the morphism will still be \'etale.  To consider the
\'etale topology, we have to look at all such morphisms. Unlike the Zariski
topology, these need not be merely open subsets of $X$, even though their
images always are.

\begin{definition}
\label{definition-etale-covering-initial}
A family of morphisms $\{ \varphi_i : U_i \to X\}_{i \in I}$ is
called an {\it \'etale covering} if each $\varphi_i$ is an \'etale morphism
and their images cover $X$, i.e.,
$X = \bigcup_{i \in I} \varphi_i(U_i)$.
\end{definition}

\noindent
This ``defines'' the \'etale topology. In other words, we can now say what the
sheaves are. An {\it \'etale sheaf} $\mathcal{F}$ of sets
(resp.\ abelian groups, vector spaces, etc) on $X$ is the data:
\begin{enumerate}
\item for each \'etale morphism $\varphi : U \to X$ a set
(resp.\ abelian group, vector space, etc) $\mathcal{F}(U)$,
\item for each pair $U, \ U'$ of \'etale schemes over $X$,
and each morphism $U \to U'$ over $X$ (which is
automatically \'etale) a restriction map
$\rho^{U'}_U : \mathcal{F}(U') \to \mathcal{F}(U)$
\end{enumerate}
These data have to satisfy the condition that $\rho^U_U = \text{id}$
in case of the identity morphism $U \to U$
and that $\rho^{U'}_U \circ \rho^{U''}_{U'} = \rho^{U''}_U$
when we have morphisms $U \to U' \to U''$ of schemes \'etale over $X$
as well as the following {\it sheaf axiom}:
\begin{itemize}
\item[(*)] for every \'etale covering $\{ \varphi_i : U_i \to U\}_{i \in
I}$, the diagram
$$
\xymatrix{
\mathcal{F} (U) \ar[r] &
\Pi_{i \in I} \mathcal{F} (U_i) \ar@<1ex>[r] \ar@<-1ex>[r] &
\Pi_{i, j \in I} \mathcal{F} (U_i \times_U U_j)
}
$$
is an equalizer diagram in the category of sets
(resp.\ abelian groups, vector spaces, etc).
\end{itemize}

\begin{remark}
\label{remark-i-is-j}
In the last statement, it is essential not to forget the case where $i = j$
which is in general a highly nontrivial condition (unlike in the Zariski
topology). In fact, frequently important coverings have only one element.
\end{remark}

\noindent
Since the identity is an \'etale morphism, we can compute the global sections
of an \'etale sheaf, and cohomology will simply be the corresponding
right-derived functors. In other words, once more theory has been developed and
statements have been made precise, there will be no obstacle to defining
cohomology.




\section{Feats of the \'etale topology}
\label{section-feats}

\noindent
For a natural number $n \in \mathbf{N} = \{1, 2, 3, 4, \ldots\}$ it is true that
$$
H_\etale^2 (\mathbf{P}^1_\mathbf{C}, \mathbf{Z}/n\mathbf{Z}) =
\mathbf{Z}/n\mathbf{Z}.
$$
More generally, if $X$ is a complex variety, then its \'etale Betti numbers
with coefficients in a finite field agree with the usual Betti numbers of
$X(\mathbf{C})$, i.e.,
$$
\dim_{\mathbf{F}_q} H_\etale^{2i} (X, \mathbf{F}_q) =
\dim_{\mathbf{F}_q} H_{Betti}^{2i} (X(\mathbf{C}), \mathbf{F}_q).
$$
This is extremely satisfactory. However, these equalities only hold for torsion
coefficients, not in general. For integer coefficients, one has
$$
H_\etale^2 (\mathbf{P}^1_\mathbf{C}, \mathbf{Z}) = 0.
$$
By contrast $H_{Betti}^2(\mathbf{P}^1(\mathbf{C}), \mathbf{Z}) = \mathbf{Z}$
as the topological space $\mathbf{P}^1(\mathbf{C})$ is homeomorphic to
a $2$-sphere.
There are ways to get back to nontorsion coefficients from torsion ones by a
limit procedure which we will come to shortly.




\section{A computation}
\label{section-computation}

\noindent
How do we compute the cohomology of $\mathbf{P}^1_\mathbf{C}$ with coefficients
$\Lambda = \mathbf{Z}/n\mathbf{Z}$?
We use {\v C}ech cohomology. A covering of $\mathbf{P}^1_\mathbf{C}$ is given
by the two standard opens $U_0, U_1$, which are both
isomorphic to $\mathbf{A}^1_\mathbf{C}$, and whose intersection is isomorphic
to $\mathbf{A}^1_\mathbf{C} \setminus \{0\} = \mathbf{G}_{m, \mathbf{C}}$.
It turns out that the Mayer-Vietoris sequence holds in \'etale cohomology.
This gives an exact sequence
$$
H_\etale^{i-1}(U_0\cap U_1, \Lambda) \to
H_\etale^i(\mathbf{P}^1_\mathbf{C}, \Lambda) \to
H_\etale^i(U_0, \Lambda) \oplus
H_\etale^i(U_1, \Lambda) \to H_\etale^i(U_0\cap U_1,
\Lambda).
$$
To get the answer we expect, we would need to show that the direct sum in the
third term vanishes. In fact, it is true that, as for the usual topology,
$$
H_\etale^q (\mathbf{A}^1_\mathbf{C}, \Lambda) = 0
\quad \text{ for } q \geq 1,
$$
and
$$
H_\etale^q (\mathbf{A}^1_\mathbf{C} \setminus \{0\}, \Lambda) = \left\{
\begin{matrix}
\Lambda & \text{ if }q = 1\text{, and} \\
0 & \text{ for }q \geq 2.
\end{matrix}
\right.
$$
These results are already quite hard (what is an elementary proof?). Let us
explain how we would compute this once the machinery of \'etale cohomology is
at our disposal.

\medskip\noindent
{\bf Higher cohomology.} This is taken care of by the following general
fact: if $X$ is an affine curve over $\mathbf{C}$, then
$$
H_\etale^q (X, \mathbf{Z}/n\mathbf{Z}) = 0 \quad \text{ for } q \geq 2.
$$
This is proved by considering the generic point of the curve and doing some
Galois cohomology. So we only have to worry about the cohomology in degree 1.

\medskip\noindent
{\bf Cohomology in degree 1.} We use the following identifications:
\begin{eqnarray*}
H_\etale^1 (X, \mathbf{Z}/n\mathbf{Z}) = \left\{
\begin{matrix}
\text{sheaves of sets }\mathcal{F}\text{ on the \'etale site }X_\etale
\text{ endowed with an} \\
\text{action }\mathbf{Z}/n\mathbf{Z} \times \mathcal{F} \to \mathcal{F}
\text{ such that }\mathcal{F}\text{ is a }\mathbf{Z}/n\mathbf{Z}\text{-torsor.}
\end{matrix}
\right\}
\Big/ \cong
\\
 = \left\{
\begin{matrix}
\text{morphisms }Y \to X\text{ which are finite \'etale together} \\
\text{ with a free }\mathbf{Z}/n\mathbf{Z}\text{ action such that }
X = Y/(\mathbf{Z}/n\mathbf{Z}).
\end{matrix}
\right\}
\Big/ \cong.
\end{eqnarray*}
The first identification is very general (it is true for any cohomology theory
on a site) and has nothing to do with the \'etale topology. The second
identification is a consequence of descent theory. The last set describes a
collection of geometric objects on which we can get our hands.

\medskip\noindent
The curve $\mathbf{A}^1_\mathbf{C}$ has no nontrivial finite \'etale covering
and hence
$H_\etale^1 (\mathbf{A}^1_\mathbf{C}, \mathbf{Z}/n\mathbf{Z}) = 0$.
This can be seen either topologically or by using the argument in the next
paragraph.

\medskip\noindent
Let us describe the finite \'etale coverings
$\varphi : Y \to \mathbf{A}^1_\mathbf{C} \setminus \{0\}$.
It suffices to consider the case where $Y$ is
connected, which we assume. We are going to find out what $Y$ can be
by applying the Riemann-Hurwitz formula (of course this is a bit silly, and
you can go ahead and skip the next section if you like).
Say that this morphism is $n$ to 1, and consider a
projective compactification
$$
\xymatrix{
{Y\ } \ar@{^{(}->}[r] \ar[d]^\varphi &
{\bar Y} \ar[d]^{\bar\varphi} \\
{\mathbf{A}^1_\mathbf{C} \setminus \{0\}} \ar@{^{(}->}[r] &
{\mathbf{P}^1_\mathbf{C}}
}
$$
Even though $\varphi$ is \'etale and does not ramify, $\bar{\varphi}$ may
ramify at 0 and $\infty$. Say that the preimages of 0 are the points $y_1,
\ldots, y_r$ with indices of ramification $e_1, \ldots e_r$, and that the
preimages of $\infty$ are the points $y_1', \ldots, y_s'$ with indices of
ramification $d_1, \ldots d_s$. In particular, $\sum e_i = n = \sum d_j$.
Applying the Riemann-Hurwitz formula, we get
$$
2 g_Y - 2 = -2n + \sum (e_i - 1) + \sum (d_j - 1)
$$
and therefore $g_Y = 0$, $r = s = 1$ and $e_1 = d_1 = n$.
Hence $Y \cong {\mathbf{A}^1_\mathbf{C} \setminus \{0\}}$, and it is easy to
see that $\varphi(z) = \lambda z^n$ for some $\lambda \in \mathbf{C}^*$.
After reparametrizing $Y$ we may assume $\lambda = 1$. Thus our
covering is given by taking the $n$th root of the coordinate on
$\mathbf{A}^1_{\mathbf{C}} \setminus \{0\}$.

\medskip\noindent
Remember that we need to classify the coverings of
${\mathbf{A}^1_\mathbf{C} \setminus \{0\}}$ together with free
$\mathbf{Z}/n\mathbf{Z}$-actions on them.
In our case any such action corresponds
to an automorphism of $Y$ sending $z$ to $\zeta_n z$, where $\zeta_n$ is a
primitive $n$th root of unity. There are $\phi(n)$ such actions
(here $\phi(n)$ means the Euler function). Thus there are exactly
$\phi(n)$ connected finite \'etale coverings with a given free
$\mathbf{Z}/n\mathbf{Z}$-action, each corresponding to a primitive
$n$th root of unity. We leave it to the reader to see that the
disconnected finite \'etale degree $n$ coverings of
$\mathbf{A}^1_{\mathbf{C}} \setminus \{0\}$ with a given free
$\mathbf{Z}/n\mathbf{Z}$-action correspond one-to-one with $n$th
roots of $1$ which are not primitive.
In other words, this computation shows that
$$
H_\etale^1 (\mathbf{A}^1_\mathbf{C} \setminus \{0\},
\mathbf{Z}/n\mathbf{Z}) =
\Hom(\mu_n(\mathbf{C}), \mathbf{Z}/n\mathbf{Z}) \cong \mathbf{Z}/n\mathbf{Z}.
$$
The first identification is canonical, the second isn't, see
Remark \ref{remark-normalize-H1-Gm}.
Since the proof of Riemann-Hurwitz does not use the computation of
cohomology, the above actually constitutes a proof (provided we
fill in the details on vanishing, etc).




\section{Nontorsion coefficients}
\label{section-nontorsion}

\noindent
To study nontorsion coefficients, one makes the following definition:
$$
H_\etale^i (X, \mathbf{Q}_\ell) :=
\left( \lim_n H_\etale^i(X, \mathbf{Z}/\ell^n\mathbf{Z}) \right)
\otimes_{\mathbf{Z}_\ell} \mathbf{Q}_\ell.
$$
The symbol $\lim_n$ denote the {\it limit} of the system of
cohomology groups $H_\etale^i(X, \mathbf{Z}/\ell^n\mathbf{Z})$ indexed
by $n$, see
Categories, Section \ref{categories-section-posets-limits}.
Thus we will need to study systems of sheaves satisfying some compatibility
conditions.




\section{Sheaf theory}
\label{section-sheaf-theory}
%9.10.09

\noindent
At this point we start talking about sites and sheaves in earnest.
There is an amazing amount of useful abstract material that could fit
in the next few sections. Some of this material is worked out in earlier
chapters, such as the chapter on sites, modules on sites, and cohomology
on sites. We try to refrain from adding too much material here, just
enough so the material later in this chapter makes sense.




\section{Presheaves}
\label{section-presheaves}

\noindent
A reference for this section is
Sites, Section \ref{sites-section-presheaves}.

\begin{definition}
\label{definition-presheaf}
Let $\mathcal{C}$ be a category. A {\it presheaf of sets} (respectively, an
{\it abelian presheaf}) on $\mathcal{C}$ is a functor $\mathcal{C}^{opp} \to
\textit{Sets}$ (resp.\ $\textit{Ab}$).
\end{definition}

\noindent
{\bf Terminology.} If $U \in \Ob(\mathcal{C})$, then elements of
$\mathcal{F}(U)$ are called {\it sections} of $\mathcal{F}$ over
$U$. For $\varphi : V \to U$ in $\mathcal{C}$, the
map $\mathcal{F}(\varphi) : \mathcal{F}(U) \to \mathcal{F}(V)$
is called the {\it restriction map} and is often denoted $s \mapsto s|_V$
or sometimes $s \mapsto \varphi^*s$. The notation $s|_V$ is ambiguous
since the restriction map depends on $\varphi$, but it is a standard
abuse of notation. We also use the notation
$\Gamma(U, \mathcal{F}) = \mathcal{F}(U)$.

\medskip\noindent
Saying that $\mathcal{F}$ is a functor means that if
$W \to V \to U$ are morphisms in $\mathcal{C}$ and
$s \in \Gamma(U, \mathcal{F})$ then
$(s|_V)|_W = s |_W$, with the abuse of
notation just seen. Moreover, the restriction mappings corresponding to
the identity morphisms $\text{id}_U : U \to U$ are the identity.

\medskip\noindent
The category of presheaves of sets (respectively of abelian presheaves) on
$\mathcal{C}$ is denoted $\textit{PSh} (\mathcal{C})$ (resp. $\textit{PAb}
(\mathcal{C})$). It is the category of functors from $\mathcal{C}^{opp}$ to
$\textit{Sets}$ (resp. $\textit{Ab}$), which is to say that the morphisms of
presheaves are natural transformations of functors. We only consider the
categories $\textit{PSh}(\mathcal{C})$ and $\textit{PAb}(\mathcal{C})$
when the category $\mathcal{C}$ is small. (Our convention is that a category
is small unless otherwise mentioned, and if it isn't small it should be
listed in Categories, Remark \ref{categories-remark-big-categories}.)

\begin{example}
\label{example-representable-presheaf}
Given an object $X \in \Ob(\mathcal{C})$, we consider the rule
$$
\begin{matrix}
U & \longmapsto & h_X(U) = \Mor_\mathcal{C}(U, X) \\
V \xrightarrow{\varphi} U & \longmapsto &
(\psi \mapsto \psi \circ \varphi) : h_X(U) \to h_X(V).
\end{matrix}
$$
This defines a functor $h_X : \mathcal{C}^{opp} \to \textit{Sets}$
and hence a presheaf. This is called the
{\it representable presheaf associated to $X$.}
It is not true that representable presheaves are sheaves in every topology on
every site.
\end{example}

\begin{lemma}[Yoneda]
\label{lemma-yoneda}
\begin{slogan}
Morphisms between objects are in bijection with natural transformations
between the functors they represent.
\end{slogan}
Let $\mathcal{C}$ be a category, and $X, Y \in
\Ob(\mathcal{C})$. There is a natural bijection
$$
\begin{matrix}
\Mor_\mathcal{C}(X, Y) &
\longrightarrow &
\Mor_{\textit{PSh}(\mathcal{C})} (h_X, h_Y) \\
\psi &
\longmapsto &
h_\psi = \psi \circ - : h_X \to h_Y.
\end{matrix}
$$
\end{lemma}

\begin{proof}
See
Categories, Lemma \ref{categories-lemma-yoneda}.
\end{proof}




\section{Sites}
\label{section-sites}


\begin{definition}
\label{definition-family-morphisms-fixed-target}
Let $\mathcal{C}$ be a category. A {\it family of morphisms with fixed target}
$\mathcal{U} = \{\varphi_i : U_i \to U\}_{i\in I}$ is the data of
\begin{enumerate}
\item an object $U \in \mathcal{C}$,
\item a set $I$ (possibly empty), and
\item for all $i\in I$, a morphism $\varphi_i : U_i \to U$ of $\mathcal{C}$
with target $U$.
\end{enumerate}
\end{definition}

\noindent
There is a notion of a {\it morphism of families of morphisms with fixed
target}. A special case of that is the notion of a {\it refinement}.
A reference for this material is
Sites, Section \ref{sites-section-refinements}.

\begin{definition}
\label{definition-site}
A {\it site}\footnote{What we call a site is a called a category endowed with
a pretopology in \cite[Expos\'e II, D\'efinition 1.3]{SGA4}.
In \cite{ArtinTopologies} it is called a category with a Grothendieck
topology.} consists of a category $\mathcal{C}$ and a set
$\text{Cov}(\mathcal{C})$ consisting of families of morphisms with fixed target
called {\it coverings}, such that
\begin{enumerate}
\item (isomorphism) if $\varphi : V \to U$ is an isomorphism in $\mathcal{C}$,
then $\{\varphi : V \to U\}$ is a covering,
\item (locality) if $\{\varphi_i : U_i \to U\}_{i\in I}$ is a covering and
for all $i \in I$ we are given a covering
$\{\psi_{ij} : U_{ij} \to U_i \}_{j\in I_i}$, then
$$
\{
\varphi_i \circ \psi_{ij} : U_{ij} \to U
\}_{(i, j)\in \prod_{i\in I} \{i\} \times I_i}
$$
is also a covering, and
\item (base change) if $\{U_i \to U\}_{i\in I}$
is a covering and $V \to U$ is a morphism in $\mathcal{C}$, then
\begin{enumerate}
\item for all $i \in I$ the fibre product
$U_i \times_U V$ exists in $\mathcal{C}$, and
\item $\{U_i \times_U V \to V\}_{i\in I}$ is a covering.
\end{enumerate}
\end{enumerate}
\end{definition}

\noindent
For us the category underlying a site is always ``small'', i.e., its
collection of objects form a set, and the collection of coverings of
a site is a set as well (as in the definition above). We will mostly,
in this chapter, leave out the arguments that cut down the collection
of objects and coverings to a set. For further discussion, see
Sites, Remark \ref{sites-remark-no-big-sites}.

\begin{example}
\label{example-site-topological-space}
If $X$ is a topological space, then it has an associated site $X_{Zar}$
defined as follows: the objects of $X_{Zar}$ are the open subsets of $X$,
the morphisms between these are the inclusion mappings, and the coverings are
the usual topological (surjective) coverings. Observe that if
$U, V \subset W \subset X$ are open subsets then $U \times_W V = U \cap V$
exists: this category has fiber products. All the verifications are trivial and
everything works as expected.
\end{example}




\section{Sheaves}
\label{section-sheaves}

\begin{definition}
\label{definition-sheaf}
A presheaf $\mathcal{F}$ of sets (resp. abelian presheaf) on a site
$\mathcal{C}$ is said to be a {\it separated presheaf} if for all coverings
$\{\varphi_i : U_i \to U\}_{i\in I} \in \text{Cov} (\mathcal{C})$
the map
$$
\mathcal{F}(U) \longrightarrow \prod\nolimits_{i\in I} \mathcal{F}(U_i)
$$
is injective. Here the map is $s \mapsto (s|_{U_i})_{i\in I}$.
The presheaf $\mathcal{F}$ is a {\it sheaf} if for all coverings
$\{\varphi_i : U_i \to U\}_{i\in I} \in \text{Cov} (\mathcal{C})$, the
diagram
\begin{equation}
\label{equation-sheaf-axiom}
\xymatrix{
\mathcal{F}(U) \ar[r] &
\prod_{i\in I} \mathcal{F}(U_i) \ar@<1ex>[r] \ar@<-1ex>[r] &
\prod_{i, j \in I} \mathcal{F}(U_i \times_U U_j),
}
\end{equation}
where the first map is $s \mapsto (s|_{U_i})_{i\in I}$ and the two
maps on the right are
$(s_i)_{i\in I} \mapsto (s_i |_{U_i \times_U U_j})$ and
$(s_i)_{i\in I} \mapsto (s_j |_{U_i \times_U U_j})$,
is an equalizer diagram in the category of sets (resp.\ abelian groups).
\end{definition}

\begin{remark}
\label{remark-empty-covering}
For the empty covering (where $I = \emptyset$), this implies that
$\mathcal{F}(\emptyset)$ is an empty product, which is a final object in the
corresponding category (a singleton, for both $\textit{Sets}$ and
$\textit{Ab}$).
\end{remark}

\begin{example}
\label{example-sheaf-site-space}
Working this out for the site $X_{Zar}$ associated to a topological
space, see Example \ref{example-site-topological-space}, gives the usual
notion of sheaves.
\end{example}

\begin{definition}
\label{definition-category-sheaves}
We denote $\Sh(\mathcal{C})$ (resp.\ $\textit{Ab}(\mathcal{C})$)
the full subcategory of $\textit{PSh}(\mathcal{C})$
(resp.\ $\textit{PAb}(\mathcal{C})$) whose objects are sheaves. This is the
{\it category of sheaves of sets} (resp.\ {\it abelian sheaves}) on
$\mathcal{C}$.
\end{definition}




\section{The example of G-sets}
\label{section-G-sets}

\noindent
Let $G$ be a group and define a site $\mathcal{T}_G$ as follows: the underlying
category is the category of $G$-sets, i.e., its objects are sets endowed
with a left $G$-action and the morphisms are equivariant maps; and the
coverings of $\mathcal{T}_G$ are the families
$\{\varphi_i : U_i \to U\}_{i\in I}$ satisfying
$U = \bigcup_{i\in I} \varphi_i(U_i)$.

\medskip\noindent
There is a special object in the site $\mathcal{T}_G$, namely the $G$-set $G$
endowed with its natural action by left translations. We denote it ${}_G G$.
Observe that there is a natural group isomorphism
$$
\begin{matrix}
\rho : & G^{opp} & \longrightarrow & \text{Aut}_{G\textit{-Sets}}({}_G G) \\
& g & \longmapsto & (h \mapsto hg).
\end{matrix}
$$
In particular, for any presheaf $\mathcal{F}$, the set $\mathcal{F}({}_G G)$
inherits a $G$-action via $\rho$. (Note that by contravariance of
$\mathcal{F}$, the set $\mathcal{F}({}_G G)$ is again a left $G$-set.) In fact,
the functor
$$
\begin{matrix}
\Sh(\mathcal{T}_G) & \longrightarrow & G\textit{-Sets} \\
\mathcal{F} & \longmapsto & \mathcal{F}({}_G G)
\end{matrix}
$$
is an equivalence of categories. Its quasi-inverse is the functor $X \mapsto
h_X$. Without giving the complete proof (which can be found in
Sites, Section \ref{sites-section-example-sheaf-G-sets})
let us try to explain why this is true.
\begin{enumerate}
\item
If $S$ is a $G$-set, we can decompose it into orbits $S = \coprod_{i\in I}
O_i$. The sheaf axiom for the covering $\{O_i \to S\}_{i\in I}$ says that
$$
\xymatrix{
\mathcal{F}(S) \ar[r] &
\prod_{i\in I} \mathcal{F}(O_i) \ar@<1ex>[r] \ar@<-1ex>[r] &
\prod_{i, j \in I} \mathcal{F}(O_i \times_S O_j)
}
$$
is an equalizer. Observing that fibered products in $G\textit{-Sets}$ are
induced from fibered products in $\textit{Sets}$, and using the fact that
$\mathcal{F}(\emptyset)$ is a $G$-singleton, we get that
$$
\prod_{i, j \in I} \mathcal{F}(O_i \times_S O_j) = \prod_{i \in I}
\mathcal{F}(O_i)
$$
and the two maps above are in fact the same. Therefore the sheaf axiom merely
says that $\mathcal{F}(S) = \prod_{i\in I} \mathcal{F}(O_i)$.
\item
If $S$ is the $G$-set $S= G/H$ and $\mathcal{F}$ is a sheaf on $\mathcal{T}_G$,
then we claim that
$$
\mathcal{F}(G/H) = \mathcal{F}({}_G G)^H
$$
and in particular $\mathcal{F}(\{*\}) = \mathcal{F}({}_G G)^G$. To see this,
let's use the sheaf axiom for the covering $\{ {}_G G \to G/H \}$ of $S$. We
have
\begin{eqnarray*}
{}_G G \times_{G/H} {}_G G & \cong & G \times H \\
(g_1, g_2) & \longmapsto & (g_1, g_1^{-1} g_2)
\end{eqnarray*}
is a disjoint union of copies of ${}_G G$ (as a $G$-set). Hence the sheaf axiom
reads
$$
\xymatrix{
\mathcal{F} (G/H) \ar[r] &
\mathcal{F}({}_G G) \ar@<1ex>[r] \ar@<-1ex>[r] &
\prod_{h\in H} \mathcal{F}({}_G G)
}
$$
where the two maps on the right are $s \mapsto (s)_{h \in H}$ and $s \mapsto
(hs)_{h \in H}$. Therefore $\mathcal{F}(G/H) = \mathcal{F}({}_G G)^H$ as
claimed.
\end{enumerate}
This doesn't quite prove the claimed equivalence of categories, but it shows at
least that a sheaf $\mathcal{F}$ is entirely determined by its sections over
${}_G G$. Details (and set theoretical remarks) can be found in
Sites, Section \ref{sites-section-example-sheaf-G-sets}.




\section{Sheafification}
\label{section-sheafification}

\begin{definition}
\label{definition-0-cech}
Let $\mathcal{F}$ be a presheaf on the site $\mathcal{C}$ and
$\mathcal{U} = \{U_i \to U\} \in \text{Cov} (\mathcal{C})$.
We define the {\it zeroth {\v C}ech cohomology group} of
$\mathcal{F}$ with respect to $\mathcal{U}$ by
$$
\check H^0 (\mathcal{U}, \mathcal{F}) =
\left\{
(s_i)_{i\in I} \in \prod\nolimits_{i\in I }\mathcal{F}(U_i)
\text{ such that }
s_i|_{U_i \times_U U_j} = s_j |_{U_i \times_U U_j}
\right\}.
$$
\end{definition}

\noindent
There is a canonical map
$\mathcal{F}(U) \to \check H^0 (\mathcal{U}, \mathcal{F})$,
$s \mapsto (s |_{U_i})_{i\in I}$.
We say that a {\it morphism of coverings} from a covering
$\mathcal{V} = \{V_j \to V\}_{j \in J}$ to $\mathcal{U}$ is a triple
$(\chi, \alpha, \chi_j)$, where
$\chi : V \to U$ is a morphism,
$\alpha : J \to I$ is a map of sets, and for all
$j \in J$ the morphism $\chi_j$ fits into a commutative diagram
$$
\xymatrix{
V_j \ar[rr]_{\chi_j} \ar[d] & & U_{\alpha(j)} \ar[d] \\
V \ar[rr]^\chi & & U.
}
$$
Given the data $\chi, \alpha, \{\chi_j\}_{j \in J}$ we define
\begin{eqnarray*}
\check H^0(\mathcal{U}, \mathcal{F}) & \longrightarrow &
\check H^0(\mathcal{V}, \mathcal{F}) \\
(s_i)_{i\in I} & \longmapsto &
\left(\chi_j^*\left(s_{\alpha(j)}\right)\right)_{j\in J}.
\end{eqnarray*}
We then claim that
\begin{enumerate}
\item the map is well-defined, and
\item depends only on $\chi$ and is independent of the choice of
$\alpha, \{\chi_j\}_{j \in J}$.
\end{enumerate}
We omit the proof of the first fact.
To see part (2), consider another triple $(\psi, \beta, \psi_j)$ with
$\chi = \psi$. Then we have the commutative diagram
$$
\xymatrix{
V_j \ar[rrr]_{(\chi_j, \psi_j)} \ar[dd] & & &
U_{\alpha(j)} \times_U U_{\beta(j)} \ar[dl] \ar[dr] \\
& & U_{\alpha(j)} \ar[dr] & &
U_{\beta(j)} \ar[dl] \\
V \ar[rrr]^{\chi = \psi} & & & U.
}
$$
Given a section $s \in \mathcal{F}(\mathcal{U})$, its image in
$\mathcal{F}(V_j)$ under the map given by
$(\chi, \alpha, \{\chi_j\}_{j \in J})$
is $\chi_j^*s_{\alpha(j)}$, and
its image under the map given by $(\psi, \beta, \{\psi_j\}_{j \in J})$
is $\psi_j^*s_{\beta(j)}$. These
two are equal since by assumption $s \in \check H^0(\mathcal{U}, \mathcal{F})$
and hence both are equal to the pullback of the common value
$$
s_{\alpha(j)}|_{U_{\alpha(j)} \times_U U_{\beta(j)}} =
s_{\beta(j)}|_{U_{\alpha(j)} \times_U U_{\beta(j)}}
$$
pulled back by the map $(\chi_j, \psi_j)$ in the diagram.

\begin{theorem}
\label{theorem-sheafification}
Let $\mathcal{C}$ be a site and $\mathcal{F}$ a presheaf on $\mathcal{C}$.
\begin{enumerate}
\item The rule
$$
U \mapsto \mathcal{F}^+(U) :=
\colim_{\mathcal{U} \text{ covering of }U}
\check H^0(\mathcal{U}, \mathcal{F})
$$
is a presheaf. And the colimit is a directed one.
\item There is a canonical map of presheaves $\mathcal{F} \to \mathcal{F}^+$.
\item If $\mathcal{F}$ is a separated presheaf then $\mathcal{F}^+$ is a sheaf
and the map in (2) is injective.
\item $\mathcal{F}^+$ is a separated presheaf.
\item $\mathcal{F}^\# = (\mathcal{F}^+)^+$ is a sheaf, and the canonical
map induces a functorial isomorphism
$$
\Hom_{\textit{PSh}(\mathcal{C})}(\mathcal{F}, \mathcal{G}) =
\Hom_{\Sh(\mathcal{C})}(\mathcal{F}^\#, \mathcal{G})
$$
for any $\mathcal{G} \in \Sh(\mathcal{C})$.
\end{enumerate}
\end{theorem}

\begin{proof}
See Sites, Theorem \ref{sites-theorem-plus}.
\end{proof}

\noindent
In other words, this means that the natural map
$\mathcal{F} \to \mathcal{F}^\#$ is a left adjoint to the forgetful functor
$\Sh(\mathcal{C}) \to \textit{PSh}(\mathcal{C})$.




\section{Cohomology}
\label{section-cohomology}

\noindent
The following is the basic result that makes it possible to define cohomology
for abelian sheaves on sites.

\begin{theorem}
\label{theorem-enough-injectives}
The category of abelian sheaves on a site is an abelian category
which has enough injectives.
\end{theorem}

\begin{proof}
See
Modules on Sites, Lemma \ref{sites-modules-lemma-abelian-abelian} and
Injectives, Theorem \ref{injectives-theorem-sheaves-injectives}.
\end{proof}

\noindent
So we can define cohomology as the right-derived functors of the
sections functor: if $U \in \Ob(\mathcal{C})$ and
$\mathcal{F} \in \textit{Ab}(\mathcal{C})$,
$$
H^p(U, \mathcal{F}) :=
R^p\Gamma(U, \mathcal{F}) =
H^p(\Gamma(U, \mathcal{I}^\bullet))
$$
where $\mathcal{F} \to \mathcal{I}^\bullet$ is an injective resolution. To do
this, we should check that the functor $\Gamma(U, -)$ is left exact. This is
true and is part of why the category $\textit{Ab}(\mathcal{C})$ is abelian,
see
Modules on Sites, Lemma \ref{sites-modules-lemma-abelian-abelian}.
For more general discussion of cohomology on sites (including the
global sections functor and its right derived functors), see
Cohomology on Sites, Section \ref{sites-cohomology-section-cohomology-sheaves}.



\section{The fpqc topology}
\label{section-fpqc}
%9.15.09

\noindent
Before doing \'etale cohomology we study a bit the fpqc topology, since
it works well for quasi-coherent sheaves.

\begin{definition}
\label{definition-fpqc-covering}
Let $T$ be a scheme. An {\it fpqc covering} of $T$ is a family
$\{ \varphi_i : T_i \to T\}_{i \in I}$ such that
\begin{enumerate}
\item each $\varphi_i$ is a flat morphism and
$\bigcup_{i\in I} \varphi_i(T_i) = T$, and
\item for each affine open $U \subset T$ there exists a finite
set $K$, a map $\mathbf{i} : K \to I$ and affine opens
$U_{\mathbf{i}(k)} \subset T_{\mathbf{i}(k)}$ such that
$U = \bigcup_{k \in K} \varphi_{\mathbf{i}(k)}(U_{\mathbf{i}(k)})$.
\end{enumerate}
\end{definition}

\begin{remark}
\label{remark-fpqc}
The first condition corresponds to fp, which stands for
{\it fid\`element plat}, faithfully flat in french, and
the second to qc, {\it quasi-compact}. The second part of
the first condition is unnecessary when the second condition holds.
\end{remark}

\begin{example}
\label{example-fpqc-coverings}
Examples of fpqc coverings.
\begin{enumerate}
\item Any Zariski open covering of $T$ is an fpqc covering.
\item A family $\{\Spec(B) \to \Spec(A)\}$ is an fpqc
covering if and only if $A \to B$ is a faithfully flat ring map.
\item If $f: X \to Y$ is flat, surjective and quasi-compact, then $\{ f: X\to
Y\}$ is an fpqc covering.
\item The morphism
$\varphi :
\coprod_{x \in \mathbf{A}^1_k} \Spec(\mathcal{O}_{\mathbf{A}^1_k, x})
\to \mathbf{A}^1_k$,
where $k$ is a field, is flat and surjective. It is not quasi-compact, and
in fact the family $\{\varphi\}$ is not an fpqc covering.
\item Write
$\mathbf{A}^2_k = \Spec(k[x, y])$. Denote $i_x : D(x) \to \mathbf{A}^2_k$
and $i_y : D(y) \to \mathbf{A}^2_k$ the standard opens.
Then the families
$\{i_x, i_y, \Spec(k[[x, y]]) \to \mathbf{A}^2_k\}$
and
$\{i_x, i_y, \Spec(\mathcal{O}_{\mathbf{A}^2_k, 0}) \to \mathbf{A}^2_k\}$
are fpqc coverings.
\end{enumerate}
\end{example}

\begin{lemma}
\label{lemma-site-fpqc}
The collection of fpqc coverings on the category of schemes
satisfies the axioms of site.
\end{lemma}

\begin{proof}
See Topologies, Lemma \ref{topologies-lemma-fpqc}.
\end{proof}

\noindent
It seems that this lemma allows us to define the fpqc site of the category
of schemes. However, there is a set theoretical problem that comes up when
considering the fpqc topology, see
Topologies, Section \ref{topologies-section-fpqc}.
It comes from our requirement that sites are ``small'', but that no small
category of schemes can contain a cofinal system of fpqc coverings of a
given nonempty scheme. Although this does not strictly speaking prevent
us from defining  ``partial'' fpqc
sites, it does not seem prudent to do so. The work-around is to allow
the notion of a sheaf for the fpqc topology (see below) but to prohibit
considering the category of all fpqc sheaves.

\begin{definition}
\label{definition-sheaf-property-fpqc}
Let $S$ be a scheme. The category of schemes over $S$ is denoted
$\Sch/S$. Consider a functor
$\mathcal{F} : (\Sch/S)^{opp} \to \textit{Sets}$, in other words
a presheaf of sets. We say $\mathcal{F}$
{\it satisfies the sheaf property for the fpqc topology}
if for every fpqc covering $\{U_i \to U\}_{i \in I}$ of schemes over $S$
the diagram (\ref{equation-sheaf-axiom}) is an equalizer diagram.
\end{definition}

\noindent
We similarly say that $\mathcal{F}$
{\it satisfies the sheaf property for the Zariski topology} if for
every open covering $U = \bigcup_{i \in I} U_i$ the diagram
(\ref{equation-sheaf-axiom}) is an equalizer diagram. See
Schemes, Definition \ref{schemes-definition-representable-by-open-immersions}.
Clearly, this is equivalent to saying that for every scheme $T$ over $S$ the
restriction of $\mathcal{F}$ to the opens of $T$ is a (usual) sheaf.

\begin{lemma}
\label{lemma-fpqc-sheaves}
Let $\mathcal{F}$ be a presheaf on $\Sch/S$. Then
$\mathcal{F}$ satisfies the sheaf property for the fpqc topology
if and only if
\begin{enumerate}
\item $\mathcal{F}$ satisfies the sheaf property with respect to the
Zariski topology, and
\item for every faithfully flat morphism $\Spec(B) \to \Spec(A)$
of affine schemes over $S$, the sheaf axiom holds for the covering
$\{\Spec(B) \to \Spec(A)\}$. Namely, this means that
$$
\xymatrix{
\mathcal{F}(\Spec(A)) \ar[r] &
\mathcal{F}(\Spec(B)) \ar@<1ex>[r] \ar@<-1ex>[r] &
\mathcal{F}(\Spec(B \otimes_A B))
}
$$
is an equalizer diagram.
\end{enumerate}
\end{lemma}

\begin{proof}
See Topologies, Lemma \ref{topologies-lemma-sheaf-property-fpqc}.
\end{proof}

\noindent
An alternative way to think of a presheaf $\mathcal{F}$ on
$\Sch/S$ which satisfies the sheaf condition for the
fpqc topology is as the following data:
\begin{enumerate}
\item for each $T/S$, a usual (i.e., Zariski) sheaf $\mathcal{F}_T$ on
$T_{Zar}$,
\item for every map $f : T' \to T$ over $S$, a restriction mapping
$f^{-1}\mathcal{F}_T \to \mathcal{F}_{T'}$
\end{enumerate}
such that
\begin{enumerate}
\item[(a)] the restriction mappings are functorial,
\item[(b)] if $f : T' \to T$ is an open immersion then the restriction
mapping $f^{-1}\mathcal{F}_T \to \mathcal{F}_{T'}$ is an isomorphism, and
\item[(c)] for every faithfully flat morphism
$\Spec(B) \to \Spec(A)$ over $S$, the diagram
$$
\xymatrix{
\mathcal{F}_{\Spec(A)}(\Spec(A)) \ar[r] &
\mathcal{F}_{\Spec(B)}(\Spec(B)) \ar@<1ex>[r] \ar@<-1ex>[r] &
\mathcal{F}_{\Spec(B \otimes_A B)}(\Spec(B \otimes_A B))
}
$$
is an equalizer.
\end{enumerate}
Data (1) and (2) and conditions (a), (b) give the data of a presheaf
on $\Sch/S$ satisfying the sheaf condition for the Zariski topology.
By Lemma \ref{lemma-fpqc-sheaves} condition (c) then suffices to get the
sheaf condition for the fpqc topology.

\begin{example}
\label{example-quasi-coherent}
Consider the presheaf
$$
\begin{matrix}
\mathcal{F} : & (\Sch/S)^{opp} & \longrightarrow & \textit{Ab} \\
& T/S & \longmapsto & \Gamma(T, \Omega_{T/S}).
\end{matrix}
$$
The compatibility of differentials with localization implies that
$\mathcal{F}$ is a sheaf on the Zariski site.
However, it does not satisfy the sheaf condition for the fpqc topology.
Namely, consider the case
$S = \Spec(\mathbf{F}_p)$ and the morphism
$$
\varphi :
V = \Spec(\mathbf{F}_p[v])
\to
U = \Spec(\mathbf{F}_p[u])
$$
given by mapping $u$ to $v^p$. The family $\{\varphi\}$ is an fpqc covering,
yet the restriction mapping
$\mathcal{F}(U) \to \mathcal{F}(V)$
sends the generator $\text{d}u$ to $\text{d}(v^p) = 0$, so
it is the zero map, and the diagram
$$
\xymatrix{
\mathcal{F}(U) \ar[r]^{0} &
\mathcal{F}(V) \ar@<1ex>[r] \ar@<-1ex>[r] &
\mathcal{F}(V \times_U V)
}
$$
is not an equalizer. We will see later that $\mathcal{F}$ does in fact
give rise to a sheaf on the \'etale and smooth sites.
\end{example}

\begin{lemma}
\label{lemma-representable-sheaf-fpqc}
Any representable presheaf on $\Sch/S$ satisfies the
sheaf condition for the fpqc topology.
\end{lemma}

\begin{proof}
See
Descent, Lemma \ref{descent-lemma-fpqc-universal-effective-epimorphisms}.
\end{proof}

\noindent
We will return to this later, since the proof of this fact uses
descent for quasi-coherent sheaves, which we will discuss in the next
section. A fancy way of expressing the lemma is to say that
{\it the fpqc topology is weaker than the canonical topology}, or
that the fpqc topology is {\it subcanonical}. In the setting of sites
this is discussed in
Sites, Section \ref{sites-section-representable-sheaves}.

\begin{remark}
\label{remark-fpqc-finest}
The fpqc is finer than the Zariski, \'etale, smooth, syntomic, and fppf
topologies. Hence any presheaf
satisfying the sheaf condition for the fpqc topology will be a
sheaf on the Zariski, \'etale, smooth, syntomic, and fppf
sites. In particular
representable presheaves will be sheaves on the \'etale site of a scheme
for example.
\end{remark}

\begin{example}
\label{example-additive-group-sheaf}
Let $S$ be a scheme.
Consider the additive group scheme $\mathbf{G}_{a, S} = \mathbf{A}^1_S$
over $S$, see
Groupoids, Example \ref{groupoids-example-additive-group}.
The associated representable presheaf is given by
$$
h_{\mathbf{G}_{a, S}}(T) =
\Mor_S(T, \mathbf{G}_{a, S}) =
\Gamma(T, \mathcal{O}_T).
$$
By the above we now know that this is a presheaf of sets which satisfies the
sheaf condition for the fpqc topology. On the other hand, it is clearly
a presheaf of rings as well. Hence we can think of this as a functor
$$
\begin{matrix}
\mathcal{O} : &
(\Sch/S)^{opp} &
\longrightarrow &
\textit{Rings} \\
&
T/S &
\longmapsto &
\Gamma(T, \mathcal{O}_T)
\end{matrix}
$$
which satisfies the sheaf condition for the fpqc topology.
Correspondingly there is a notion of $\mathcal{O}$-module, and so on and
so forth.
\end{example}




\section{Faithfully flat descent}
\label{section-fpqc-descent}

\noindent
In this section we discuss faithfully flat descent for quasi-coherent
modules. More precisely, we will prove quasi-coherent modules satisfy
effective descent with respect to fpqc coverings.

\begin{definition}
\label{definition-descent-datum}
Let $\mathcal{U} = \{ t_i : T_i \to T\}_{i \in I}$ be a family of
morphisms of schemes with fixed target. A {\it descent datum} for
quasi-coherent sheaves with respect to $\mathcal{U}$ is a collection
$((\mathcal{F}_i)_{i \in I}, (\varphi_{ij})_{i, j \in I})$ where
\begin{enumerate}
\item $\mathcal{F}_i$ is a quasi-coherent sheaf on $T_i$, and
\item $\varphi_{ij} : \text{pr}_0^* \mathcal{F}_i \to
\text{pr}_1^* \mathcal{F}_j$ is an isomorphism of modules
on $T_i \times_T T_j$,
\end{enumerate}
such that the {\it cocycle condition} holds: the diagrams
$$
\xymatrix{
\text{pr}_0^*\mathcal{F}_i \ar[dr]_{\text{pr}_{02}^*\varphi_{ik}}
\ar[rr]^{\text{pr}_{01}^*\varphi_{ij}} & &
\text{pr}_1^*\mathcal{F}_j \ar[dl]^{\text{pr}_{12}^*\varphi_{jk}} \\
& \text{pr}_2^*\mathcal{F}_k
}
$$
commute on $T_i \times_T T_j \times_T T_k$.
This descent datum is called {\it effective} if there exist a quasi-coherent
sheaf $\mathcal{F}$ over $T$ and $\mathcal{O}_{T_i}$-module isomorphisms
$\varphi_i : t_i^* \mathcal{F} \cong \mathcal{F}_i$ compatible with
the maps $\varphi_{ij}$, namely
$$
\varphi_{ij} = \text{pr}_1^* (\varphi_j) \circ \text{pr}_0^* (\varphi_i)^{-1}.
$$
\end{definition}

\noindent
In this and the next section we discuss some ingredients of the proof
of the following theorem, as well as some related material.

\begin{theorem}
\label{theorem-descent-quasi-coherent}
If $\mathcal{V} = \{T_i \to T\}_{i\in I}$ is an fpqc covering, then all
descent data for quasi-coherent sheaves with respect to $\mathcal{V}$
are effective.
\end{theorem}

\begin{proof}
See
Descent, Proposition \ref{descent-proposition-fpqc-descent-quasi-coherent}.
\end{proof}

\noindent
In other words, the fibered category of quasi-coherent sheaves is a stack on
the fpqc site.
The proof of the theorem is in two steps. The first one is to realize that for
Zariski coverings this is easy (or well-known) using standard glueing of
sheaves (see
Sheaves, Section \ref{sheaves-section-glueing-sheaves})
and the locality of quasi-coherence. The second step is the case of an
fpqc covering of the form $\{\Spec(B) \to \Spec(A)\}$
where $A \to B$ is a faithfully flat ring map.
This is a lemma in algebra, which we now present.

\medskip\noindent
{\bf Descent of modules.}
If $A \to B$ is a ring map, we consider the complex
$$
(B/A)_\bullet : B \to B \otimes_A B \to B \otimes_A B \otimes_A B \to \ldots
$$
where $B$ is in degree 0, $B \otimes_A B$ in degree 1, etc, and the maps are
given by
\begin{eqnarray*}
b & \mapsto & 1 \otimes b - b \otimes 1, \\
b_0 \otimes b_1 & \mapsto & 1 \otimes b_0 \otimes b_1 - b_0 \otimes 1 \otimes
b_1 + b_0 \otimes b_1 \otimes 1, \\
& \text{etc.}
\end{eqnarray*}

\begin{lemma}
\label{lemma-algebra-descent}
If $A \to B$ is faithfully flat, then the complex $(B/A)_\bullet$ is exact in
positive degrees, and $H^0((B/A)_\bullet) = A$.
\end{lemma}

\begin{proof}
See Descent, Lemma \ref{descent-lemma-ff-exact}.
\end{proof}

\noindent
Grothendieck proves this in three steps. Firstly, he assumes that the map $A
\to B$ has a section, and constructs an explicit homotopy to the complex where
$A$ is the only nonzero term, in degree 0. Secondly, he observes that to prove
the result, it suffices to do so after a faithfully flat base change $A \to
A'$, replacing $B$ with $B' = B \otimes_A A'$. Thirdly, he applies the
faithfully flat base change $A \to A' = B$ and remark that the map
$A' = B \to B' = B \otimes_A B$ has a natural section.

\medskip\noindent
The same strategy proves the following lemma.

\begin{lemma}
\label{lemma-descent-modules}
If $A \to B$ is faithfully flat and $M$ is an $A$-module, then the
complex $(B/A)_\bullet \otimes_A M$ is exact in positive degrees, and
$H^0((B/A)_\bullet \otimes_A M) = M$.
\end{lemma}

\begin{proof}
See Descent, Lemma \ref{descent-lemma-ff-exact}.
\end{proof}

\begin{definition}
\label{definition-descent-datum-modules}
Let $A \to B$ be a ring map and $N$ a $B$-module. A {\it descent datum} for
$N$ with respect to $A \to B$ is an isomorphism
$\varphi : N \otimes_A B \cong B \otimes_A N$ of $B \otimes_A B$-modules such
that the diagram of $B \otimes_A B \otimes_A B$-modules
$$
\xymatrix{
{N \otimes_A B \otimes_A B} \ar[dr]_{\varphi_{02}} \ar[rr]^{\varphi_{01}} & &
{B \otimes_A N \otimes_A B} \ar[dl]^{\varphi_{12}} \\
& {B \otimes_A B \otimes_A N}
}
$$
commutes where $\varphi_{01} = \varphi \otimes \text{id}_B$ and similarly
for $\varphi_{12}$ and $\varphi_{02}$.
\end{definition}

\noindent
If $N' = B \otimes_A M$ for some $A$-module M, then it has a canonical descent
datum given by the map
$$
\begin{matrix}
\varphi_\text{can}: & N' \otimes_A B & \to & B \otimes_A N' \\
& b_0 \otimes m \otimes b_1 & \mapsto & b_0 \otimes b_1 \otimes m.
\end{matrix}
$$

\begin{definition}
\label{definition-effective-modules}
A descent datum $(N, \varphi)$ is called {\it effective} if there exists an
$A$-module $M$ such that $(N, \varphi) \cong (B \otimes_A M,
\varphi_\text{can})$, with the obvious notion of isomorphism of descent data.
\end{definition}

\noindent
Theorem \ref{theorem-descent-quasi-coherent} is a consequence the
following result.

\begin{theorem}
\label{theorem-descent-modules}
If $A \to B$ is faithfully flat then descent data with respect to $A\to B$
are effective.
\end{theorem}

\begin{proof}
See
Descent, Proposition \ref{descent-proposition-descent-module}.
See also
Descent, Remark \ref{descent-remark-homotopy-equivalent-cosimplicial-algebras}
for an alternative view of the proof.
\end{proof}

\begin{remarks}
\label{remarks-theorem-modules-exactness}
The results on descent of modules have several applications:
\begin{enumerate}
\item The exactness of the {\v C}ech complex in positive degrees for
the covering $\{\Spec(B) \to \Spec(A)\}$ where $A \to B$ is
faithfully flat. This will give some vanishing of cohomology.
\item If $(N, \varphi)$ is a descent datum with respect to a faithfully
flat map $A \to B$, then the corresponding $A$-module is given by
$$
M = \Ker \left(
\begin{matrix}
N & \longrightarrow & B \otimes_A N \\
n & \longmapsto & 1 \otimes n - \varphi(n \otimes 1)
\end{matrix}
\right).
$$
See
Descent, Proposition \ref{descent-proposition-descent-module}.
\end{enumerate}
\end{remarks}




%9.17.09
\section{Quasi-coherent sheaves}
\label{section-quasi-coherent}

\noindent
We can apply the descent of modules to study quasi-coherent sheaves.

\begin{proposition}
\label{proposition-quasi-coherent-sheaf-fpqc}
For any quasi-coherent sheaf $\mathcal{F}$ on $S$ the presheaf
$$
\begin{matrix}
\mathcal{F}^a : & \Sch/S & \to & \textit{Ab}\\
& (f: T \to S) & \mapsto & \Gamma(T, f^*\mathcal{F})
\end{matrix}
$$
is an $\mathcal{O}$-module which satisfies the sheaf condition for the
fpqc topology.
\end{proposition}

\begin{proof}
This is proved in
Descent, Lemma \ref{descent-lemma-sheaf-condition-holds}.
We indicate the proof here. As established in
Lemma \ref{lemma-fpqc-sheaves},
it is enough to check the sheaf property
on Zariski coverings and faithfully flat morphisms of affine schemes. The
sheaf property for Zariski coverings is standard scheme theory, since
$\Gamma(U, i^\ast \mathcal{F}) = \mathcal{F}(U)$ when
$i : U \hookrightarrow S$ is an open immersion.

\medskip\noindent
For $\left\{\Spec(B)\to \Spec(A)\right\}$ with $A\to B$ faithfully
flat and
$\mathcal{F}|_{\Spec(A)} = \widetilde{M}$
this corresponds to the fact that
$M = H^0\left((B/A)_\bullet \otimes_A M \right)$, i.e., that
\begin{align*}
0 \to M \to B \otimes_A M \to B \otimes_A B \otimes_A M
\end{align*}
is exact by
Lemma \ref{lemma-descent-modules}.
\end{proof}

\noindent
There is an abstract notion of a quasi-coherent sheaf on a ringed site.
We briefly introduce this here. For more information please consult
Modules on Sites, Section \ref{sites-modules-section-local}.
Let $\mathcal{C}$ be a category, and let $U$ be an object of $\mathcal{C}$.
Then $\mathcal{C}/U$ indicates the category of objects over $U$, see
Categories, Example \ref{categories-example-category-over-X}.
If $\mathcal{C}$ is a site, then $\mathcal{C}/U$ is a site as well, namely
the coverings of $V/U$ are families $\{V_i/U \to V/U\}$ of morphisms
of $\mathcal{C}/U$ with fixed target such that
$\{V_i \to V\}$ is a covering of $\mathcal{C}$. Moreover, given any
sheaf $\mathcal{F}$ on $\mathcal{C}$ the {\it restriction}
$\mathcal{F}|_{\mathcal{C}/U}$ (defined in the obvious manner)
is a sheaf as well. See
Sites, Section \ref{sites-section-localize}
for details.

\begin{definition}
\label{definition-ringed-site}
Let $\mathcal{C}$ be a {\it ringed site}, i.e., a site endowed with a
sheaf of rings $\mathcal{O}$. A sheaf of $\mathcal{O}$-modules $\mathcal{F}$ on
$\mathcal{C}$ is called {\it quasi-coherent} if for all
$U \in \Ob(\mathcal{C})$ there exists a covering
$\{U_i \to U\}_{i\in I}$ of $\mathcal{C}$ such that the restriction
$\mathcal{F}|_{\mathcal{C}/U_i}$ is isomorphic to the cokernel of
an $\mathcal{O}$-linear map of free $\mathcal{O}$-modules
$$
\bigoplus\nolimits_{k \in K} \mathcal{O}|_{\mathcal{C}/U_i}
\longrightarrow
\bigoplus\nolimits_{l \in L} \mathcal{O}|_{\mathcal{C}/U_i}.
$$
The direct sum over $K$ is the sheaf associated to the presheaf
$V \mapsto \bigoplus_{k \in K} \mathcal{O}(V)$ and similarly for the other.
\end{definition}

\noindent
Although it is useful to be able to give a general definition as above
this notion is not well behaved in general.

\begin{remark}
\label{remark-final-object}
In the case where $\mathcal{C}$ has a final object, e.g.\ $S$, it
suffices to check the condition of the definition for
$U = S$ in the above statement. See
Modules on Sites, Lemma \ref{sites-modules-lemma-local-final-object}.
\end{remark}

\begin{theorem}[Meta theorem on quasi-coherent sheaves]
\label{theorem-quasi-coherent}
Let $S$ be a scheme.
Let $\mathcal{C}$ be a site. Assume that
\begin{enumerate}
\item the underlying category $\mathcal{C}$ is a
full subcategory of $\Sch/S$,
\item any Zariski covering of $T \in \Ob(\mathcal{C})$
can be refined by a covering of $\mathcal{C}$,
\item $S/S$ is an object of $\mathcal{C}$,
\item every covering of $\mathcal{C}$ is an fpqc covering of schemes.
\end{enumerate}
Then the presheaf $\mathcal{O}$ is a sheaf on $\mathcal{C}$ and
any quasi-coherent $\mathcal{O}$-module on $(\mathcal{C}, \mathcal{O})$
is of the form $\mathcal{F}^a$ for some quasi-coherent sheaf
$\mathcal{F}$ on $S$.
\end{theorem}

\begin{proof}
After some formal arguments this is exactly Theorem
\ref{theorem-descent-quasi-coherent}. Details omitted. In
Descent, Proposition \ref{descent-proposition-equivalence-quasi-coherent}
we prove a more precise version of the theorem for the
big Zariski, fppf, \'etale, smooth, and syntomic sites of $S$,
as well as the small Zariski and \'etale sites of $S$.
\end{proof}

\noindent
In other words, there is no difference between quasi-coherent
modules on the scheme $S$ and quasi-coherent $\mathcal{O}$-modules
on sites $\mathcal{C}$ as in the theorem. More precise statements
for the big and small sites $(\Sch/S)_{fppf}$, $S_\etale$, etc
can be found in Descent, Sections
\ref{descent-section-quasi-coherent-sheaves},
\ref{descent-section-quasi-coherent-cohomology}, and
\ref{descent-section-quasi-coherent-sheaves-bis}.
In this chapter we will sometimes refer to a
``site as in Theorem \ref{theorem-quasi-coherent}''
in order to conveniently state results which hold in any of those
situations.






\section{{\v C}ech cohomology}
\label{section-cech-cohomology}

\noindent
Our next goal is to use descent theory to show that
$H^i(\mathcal{C}, \mathcal{F}^a) = H_{Zar}^i(S, \mathcal{F})$
for all quasi-coherent sheaves $\mathcal{F}$ on $S$, and
any site $\mathcal{C}$ as in Theorem \ref{theorem-quasi-coherent}.
To this end, we introduce {\v C}ech cohomology on sites.
See \cite{ArtinTopologies} and
Cohomology on Sites, Sections \ref{sites-cohomology-section-cech},
\ref{sites-cohomology-section-cech-functor}
and \ref{sites-cohomology-section-cech-cohomology-cohomology}
for more details.

\begin{definition}
\label{definition-cech-complex}
Let $\mathcal{C}$ be a category. Let
$\mathcal{U} = \{U_i \to U\}_{i \in I}$ be a family of morphisms of
$\mathcal{C}$ with fixed target. Assume that all the fibre products
$U_{i_0} \times_U \ldots \times_U U_{i_p}$ exists in $\mathcal{C}$.
Let $\mathcal{F} \in \textit{PAb}(\mathcal{C})$ be an abelian
presheaf. We define the {\it {\v C}ech complex}
$\check{\mathcal{C}}^\bullet(\mathcal{U}, \mathcal{F})$ by
$$
\prod_{i_0 \in I} \mathcal{F}(U_{i_0}) \to
\prod_{i_0, i_1 \in I} \mathcal{F}(U_{i_0} \times_U U_{i_1}) \to
\prod_{i_0, i_1, i_2 \in I}
\mathcal{F}(U_{i_0} \times_U U_{i_1} \times_U U_{i_2}) \to \ldots
$$
where the first term is in degree 0 and the maps are the usual ones.
The {\it {\v C}ech cohomology groups} are defined by
$$
\check{H}^p(\mathcal{U}, \mathcal{F}) =
H^p(\check{\mathcal{C}}^\bullet(\mathcal{U}, \mathcal{F})).
$$
\end{definition}

\noindent
In the above definition, it is essential to allow the indices
$(i_0, \ldots, i_p)$ to have repetitions.

\begin{lemma}
\label{lemma-cech-presheaves}
Notation and assumptions as in Definition \ref{definition-cech-complex}.
The functor $\check{\mathcal{C}}^\bullet(\mathcal{U}, -)$
is exact on the category $\textit{PAb}(\mathcal{C})$.
\end{lemma}

\noindent
In other words, if $0\to \mathcal{F}_1\to \mathcal{F}_2\to \mathcal{F}_3\to 0$
is a short exact sequence of presheaves of abelian groups, then
$$
0 \to \check{\mathcal{C}}^\bullet\left(\mathcal{U}, \mathcal{F}_1\right)
\to\check{\mathcal{C}}^\bullet(\mathcal{U}, \mathcal{F}_2) \to
\check{\mathcal{C}}^\bullet(\mathcal{U}, \mathcal{F}_3)\to 0
$$
is a short exact sequence of complexes.

\begin{proof}
This follows at once from the definition of a short exact sequence of
presheaves. Namely, as the category of abelian presheaves is the category of
functors on some category with values in $\textit{Ab}$, it is automatically an
abelian category: a sequence $\mathcal{F}_1\to \mathcal{F}_2\to \mathcal{F}_3$
is exact in $\textit{PAb}$ if and only if for all
$U \in \Ob(\mathcal{C})$, the sequence
$\mathcal{F}_1(U) \to \mathcal{F}_2(U) \to \mathcal{F}_3(U)$ is exact in
$\textit{Ab}$. So the complex above is merely a product of short exact
sequences in each degree. See also
Cohomology on Sites, Lemma \ref{sites-cohomology-lemma-cech-exact-presheaves}.
\end{proof}

\noindent
This shows that $\check{H}^\bullet(\mathcal{U}, -)$ is a $\delta$-functor.
We now proceed to show that it is a universal $\delta$-functor. We thus need to
show that it is an {\it effaceable} functor. We start by recalling the Yoneda
lemma.

\begin{lemma}[Yoneda Lemma]
\label{lemma-yoneda-presheaf}
For any presheaf $\mathcal{F}$ on a category $\mathcal{C}$
and $U \in \Ob(\mathcal{C})$ there is a functorial isomorphism
$$
\Hom_{\textit{PSh}(\mathcal{C})}(h_U, \mathcal{F}) =
\mathcal{F}(U).
$$
\end{lemma}

\begin{proof}
See Categories, Lemma \ref{categories-lemma-yoneda}.
\end{proof}

\noindent
Given a set $E$ we denote (in this section)
$\mathbf{Z}[E]$ the free abelian group on $E$. In a formula
$\mathbf{Z}[E] = \bigoplus_{e \in E} \mathbf{Z}$, i.e., $\mathbf{Z}[E]$ is
a free $\mathbf{Z}$-module having a basis consisting of the elements of $E$.
Using this notation we introduce the free abelian presheaf on a
presheaf of sets.

\begin{definition}
\label{definition-free-abelian-presheaf}
Let $\mathcal{C}$ be a category.
Given a presheaf of sets $\mathcal{G}$, we define the
{\it free abelian presheaf on $\mathcal{G}$},
denoted $\mathbf{Z}_\mathcal{G}$, by the rule
$$
\mathbf{Z}_\mathcal{G}(U)
=
\mathbf{Z}[\mathcal{G}(U)]
$$
for $U \in \Ob(\mathcal{C})$
with restriction maps induced by the restriction maps of $\mathcal{G}$.
In the special case $\mathcal{G} = h_U$ we write simply
$\mathbf{Z}_U = \mathbf{Z}_{h_U}$.
\end{definition}

\noindent
The functor $\mathcal{G} \mapsto \mathbf{Z}_\mathcal{G}$ is left adjoint to the
forgetful functor $\textit{PAb}(\mathcal{C}) \to \textit{PSh}(\mathcal{C})$.
Thus, for any presheaf $\mathcal{F}$, there is a canonical isomorphism
$$
\Hom_{\textit{PAb}(\mathcal{C})}(\mathbf{Z}_U, \mathcal{F})
=
\Hom_{\textit{PSh}(\mathcal{C})}(h_U, \mathcal{F})
=
\mathcal{F}(U)
$$
the last equality by the Yoneda lemma. In particular, we have the following
result.

\begin{lemma}
\label{lemma-cech-complex-describe}
Notation and assumptions as in Definition \ref{definition-cech-complex}.
The {\v C}ech complex $\check{\mathcal{C}}^\bullet(\mathcal{U}, \mathcal{F})$
can be described explicitly as follows
\begin{eqnarray*}
\check{\mathcal{C}}^\bullet(\mathcal{U}, \mathcal{F})
& = &
\left(
\prod_{i_0 \in I}
\Hom_{\textit{PAb}(\mathcal{C})}(\mathbf{Z}_{U_{i_0}}, \mathcal{F}) \to
\prod_{i_0, i_1 \in I}
\Hom_{\textit{PAb}(\mathcal{C})}(
\mathbf{Z}_{U_{i_0} \times_U U_{i_1}}, \mathcal{F}) \to \ldots
\right) \\
& = &
\Hom_{\textit{PAb}(\mathcal{C})}\left(
\left(
\bigoplus_{i_0 \in I} \mathbf{Z}_{U_{i_0}} \leftarrow
\bigoplus_{i_0, i_1 \in I} \mathbf{Z}_{U_{i_0} \times_U U_{i_1}} \leftarrow
\ldots
\right), \mathcal{F}\right)
\end{eqnarray*}
\end{lemma}

\begin{proof}
This follows from the formula above. See
Cohomology on Sites, Lemma \ref{sites-cohomology-lemma-cech-map-into}.
\end{proof}

\noindent
This reduces us to studying only the complex in the first argument of the
last $\Hom$.

\begin{lemma}
\label{lemma-exact}
Notation and assumptions as in Definition \ref{definition-cech-complex}.
The complex of abelian presheaves
\begin{align*}
\mathbf{Z}_\mathcal{U}^\bullet \quad : \quad
\bigoplus_{i_0 \in I} \mathbf{Z}_{U_{i_0}} \leftarrow
\bigoplus_{i_0, i_1 \in I} \mathbf{Z}_{U_{i_0} \times_U U_{i_1}} \leftarrow
\bigoplus_{i_0, i_1, i_2 \in I}
\mathbf{Z}_{U_{i_0} \times_U U_{i_1} \times_U U_{i_2}} \leftarrow
\ldots
\end{align*}
is exact in all degrees except $0$ in $\textit{PAb}(\mathcal{C})$.
\end{lemma}

\begin{proof}
For any $V\in \Ob(\mathcal{C})$ the complex of abelian groups
$\mathbf{Z}_\mathcal{U}^\bullet(V)$ is
$$
\begin{matrix}
\mathbf{Z}\left[
\coprod_{i_0\in I} \Mor_\mathcal{C}(V, U_{i_0})\right]
\leftarrow
\mathbf{Z}\left[
\coprod_{i_0, i_1 \in I}
\Mor_\mathcal{C}(V, U_{i_0} \times_U U_{i_1})\right]
\leftarrow \ldots = \\
\bigoplus_{\varphi : V \to U}
\left(
\mathbf{Z}\left[\coprod_{i_0 \in I} \Mor_\varphi(V, U_{i_0})\right]
\leftarrow
\mathbf{Z}\left[\coprod_{i_0, i_1\in I} \Mor_\varphi(V, U_{i_0}) \times
\Mor_\varphi(V, U_{i_1})\right]
\leftarrow
\ldots
\right)
\end{matrix}
$$
where
$$
\Mor_{\varphi}(V, U_i)
=
\{ V \to U_i \text{ such that } V \to U_i \to U \text{ equals } \varphi \}.
$$
Set $S_\varphi = \coprod_{i\in I} \Mor_\varphi(V, U_i)$, so that
$$
\mathbf{Z}_\mathcal{U}^\bullet(V)
=
\bigoplus_{\varphi : V \to U}
\left(
\mathbf{Z}[S_\varphi] \leftarrow
\mathbf{Z}[S_\varphi \times S_\varphi] \leftarrow
\mathbf{Z}[S_\varphi \times S_\varphi \times S_\varphi] \leftarrow
\ldots \right).
$$
Thus it suffices to show that for each $S = S_\varphi$, the complex
\begin{align*}
\mathbf{Z}[S] \leftarrow
\mathbf{Z}[S \times S] \leftarrow
\mathbf{Z}[S \times S \times S] \leftarrow \ldots
\end{align*}
is exact in negative degrees. To see this, we can give an explicit homotopy.
Fix $s\in S$ and define $K: n_{(s_0, \ldots, s_p)} \mapsto n_{(s, s_0,
\ldots, s_p)}.$ One easily checks that $K$ is a nullhomotopy for the operator
$$
\delta :
\eta_{(s_0, \ldots, s_p)}
\mapsto
\sum\nolimits_{i = 0}^p (-1)^p \eta_{(s_0, \ldots, \hat s_i, \ldots, s_p)}.
$$
See
Cohomology on Sites, Lemma \ref{sites-cohomology-lemma-homology-complex}
for more details.
\end{proof}

\begin{lemma}
\label{lemma-hom-injective}
Notation and assumptions as in Definition \ref{definition-cech-complex}.
If $\mathcal{I}$ is an injective object of $\textit{PAb}(\mathcal{C})$,
then $\check H^p(\mathcal{U}, \mathcal{I}) = 0$ for all $p > 0$.
\end{lemma}

\begin{proof}
The {\v C}ech complex is the result of applying the functor
$\Hom_{\textit{PAb}(\mathcal{C})}(-, \mathcal{I}) $ to the complex $
\mathbf{Z}^\bullet_\mathcal{U} $, i.e.,
$$
\check H^p(\mathcal{U}, \mathcal{I}) = H^p
(\Hom_{\textit{PAb}(\mathcal{C})} (\mathbf{Z}^\bullet_\mathcal{U},
\mathcal{I})).
$$
But we have just seen that $\mathbf{Z}^\bullet_\mathcal{U}$ is exact in
negative degrees, and the functor $\Hom_{\textit{PAb}(\mathcal{C})}(-,
\mathcal{I})$ is exact, hence $\Hom_{\textit{PAb}(\mathcal{C})}
(\mathbf{Z}^\bullet_\mathcal{U}, \mathcal{I})$ is exact in positive degrees.
\end{proof}

\begin{theorem}
\label{theorem-cech-derived}
Notation and assumptions as in Definition \ref{definition-cech-complex}.
On $\textit{PAb}(\mathcal{C})$ the functors $\check{H}^p(\mathcal{U}, -)$ are
the right derived functors of $\check{H}^0(\mathcal{U}, -)$.
\end{theorem}

\begin{proof}
By the Lemma \ref{lemma-hom-injective}, the functors
$\check H^p(\mathcal{U}, -)$ are universal
$\delta$-functors since they are effaceable.
So are the right derived functors of $\check H^0(\mathcal{U}, -)$. Since they
agree in degree $0$, they agree by the universal property of universal
$\delta$-functors. For more details see
Cohomology on Sites,
Lemma \ref{sites-cohomology-lemma-cech-cohomology-derived-presheaves}.
\end{proof}

\begin{remark}
\label{remark-presheaves-no-topology}
Observe that all of the preceding statements are about presheaves so we haven't
made use of the topology yet.
\end{remark}




\section{The {\v C}ech-to-cohomology spectral sequence}
\label{section-cech-ss}

\noindent
This spectral sequence is fundamental in proving foundational results on
cohomology of sheaves.

\begin{lemma}
\label{lemma-forget-injectives}
The forgetful functor $\textit{Ab}(\mathcal{C})\to \textit{PAb}(\mathcal{C})$
transforms injectives into injectives.
\end{lemma}

\begin{proof}
This is formal using the fact that the forgetful functor has a left adjoint,
namely sheafification, which is an exact functor. For more details see
Cohomology on Sites,
Lemma \ref{sites-cohomology-lemma-injective-abelian-sheaf-injective-presheaf}.
\end{proof}

\begin{theorem}
\label{theorem-cech-ss}
Let $\mathcal{C}$ be a site. For any covering
$\mathcal{U} = \{U_i \to U\}_{i \in I}$ of $U \in \Ob(\mathcal{C})$
and any abelian sheaf $\mathcal{F}$ on $\mathcal{C}$
there is a spectral sequence
$$
E_2^{p, q}
=
\check H^p(\mathcal{U}, \underline{H}^q(\mathcal{F}))
\Rightarrow
H^{p+q}(U, \mathcal{F}),
$$
where $\underline{H}^q(\mathcal{F})$ is the abelian presheaf
$V \mapsto H^q(V, \mathcal{F})$.
\end{theorem}

\begin{proof}
Choose an injective resolution $\mathcal{F}\to \mathcal{I}^\bullet$ in
$\textit{Ab}(\mathcal{C})$, and consider the double complex
$\check{\mathcal{C}}^\bullet(\mathcal{U}, \mathcal{I}^\bullet)$
and the maps
$$
\xymatrix{
\Gamma(U, I^\bullet) \ar[r] &
\check{\mathcal{C}}^\bullet(\mathcal{U}, \mathcal{I}^\bullet) \\
& \check{\mathcal{C}}^\bullet(\mathcal{U}, \mathcal{F}) \ar[u]
}
$$
Here the horizontal map is the natural map
$\Gamma(U, I^\bullet) \to
\check{\mathcal{C}}^0(\mathcal{U}, \mathcal{I}^\bullet)$
to the left column, and the vertical map is induced by
$\mathcal{F}\to \mathcal{I}^0$ and lands in the bottom row.
By assumption, $\mathcal{I}^\bullet$ is a complex of injectives in
$\textit{Ab}(\mathcal{C})$, hence by
Lemma \ref{lemma-forget-injectives}, it is a complex of injectives in
$\textit{PAb}(\mathcal{C})$. Thus, the rows of the double complex are
exact in positive degrees (Lemma \ref{lemma-hom-injective}), and
the kernel of $\check{\mathcal{C}}^0(\mathcal{U}, \mathcal{I}^\bullet)
\to \check{\mathcal{C}}^1(\mathcal{U}, \mathcal{I}^\bullet)$
is equal to
$\Gamma(U, \mathcal{I}^\bullet)$, since $\mathcal{I}^\bullet$
is a complex of sheaves. In particular, the cohomology of the total complex
is the standard
cohomology of the global sections functor $H^0(U, \mathcal{F})$.

\medskip\noindent
For the vertical direction, the $q$th cohomology group of the $p$th column is
$$
\prod_{i_0, \ldots, i_p}
H^q(U_{i_0} \times_U \ldots \times_U U_{i_p}, \mathcal{F})
=
\prod_{i_0, \ldots, i_p}
\underline{H}^q(\mathcal{F})(U_{i_0} \times_U \ldots \times_U U_{i_p})
$$
in the entry $E_1^{p, q}$. So this is a standard double complex spectral
sequence, and the $E_2$-page is as prescribed. For more details see
Cohomology on Sites,
Lemma \ref{sites-cohomology-lemma-cech-spectral-sequence}.
\end{proof}

\begin{remark}
\label{remark-grothendieck-ss}
This is a Grothendieck spectral sequence for the composition of functors
$$
\textit{Ab}(\mathcal{C}) \longrightarrow
\textit{PAb}(\mathcal{C}) \xrightarrow{\check H^0} \textit{Ab}.
$$
\end{remark}








\section{Big and small sites of schemes}
\label{section-big-small}

\noindent
Let $S$ be a scheme.
Let $\tau$ be one of the topologies we will be discussing.
Thus $\tau \in \{fppf, syntomic, smooth, \etale, Zariski\}$.
Of course if you are only interested in the \'etale topology, then
you can simply assume $\tau = \etale$ throughout. Moreover, we will
discuss \'etale morphisms, \'etale coverings, and \'etale sites
in more detail starting in Section \ref{section-etale-site}.
In order to proceed with the discussion of cohomology of
quasi-coherent sheaves it is convenient to introduce the
big $\tau$-site and in case $\tau \in \{\etale, Zariski\}$, the
small $\tau$-site of $S$. In order to do this we first introduce
the notion of a $\tau$-covering.

\begin{definition}
\label{definition-tau-covering}
(See
Topologies, Definitions
\ref{topologies-definition-fppf-covering},
\ref{topologies-definition-syntomic-covering},
\ref{topologies-definition-smooth-covering},
\ref{topologies-definition-etale-covering}, and
\ref{topologies-definition-zariski-covering}.)
Let $\tau \in \{fppf, syntomic, smooth, \etale, Zariski\}$.
A family of morphisms of schemes $\{f_i : T_i \to T\}_{i \in I}$ with fixed
target is called a {\it $\tau$-covering} if and only if
each $f_i$ is flat of finite presentation, syntomic, smooth, \'etale,
resp.\ an open immersion, and we have $\bigcup f_i(T_i) = T$.
\end{definition}

\noindent
The class of all $\tau$-coverings satisfies the axioms
(1), (2) and (3) of
Definition \ref{definition-site} (our definition of a site), see
Topologies, Lemmas
\ref{topologies-lemma-fppf},
\ref{topologies-lemma-syntomic},
\ref{topologies-lemma-smooth},
\ref{topologies-lemma-etale}, and
\ref{topologies-lemma-zariski}.

\medskip\noindent
Let us introduce the sites we will be working with.
Contrary to what happens in
\cite{SGA4}, we do not want to choose a universe. Instead we pick a ``partial
universe'' (which is a suitably large set as in
Sets, Section \ref{sets-section-sets-hierarchy}), and consider all schemes
contained in this set. Of course we make sure that our favorite base scheme
$S$ is contained in the partial universe. Having picked the underlying category
we pick a suitably large set of $\tau$-coverings which turns this into a site.
The details are in the chapter on topologies on schemes; there is a lot of
freedom in the choices made, but in the end the actual choices made will not
affect the \'etale (or other) cohomology of $S$ (just as in \cite{SGA4} the
actual choice of universe doesn't matter at the end). Moreover, the way the
material is written the reader who is happy using strongly inaccessible
cardinals (i.e., universes) can do so as a substitute.

\begin{definition}
\label{definition-tau-site}
Let $S$ be a scheme.
Let $\tau \in \{fppf, syntomic, smooth, \etale, \linebreak[0] Zariski\}$.
\begin{enumerate}
\item A {\it big $\tau$-site of $S$} is any of the sites
$(\Sch/S)_\tau$ constructed as explained above and in more detail in
Topologies, Definitions
\ref{topologies-definition-big-small-fppf},
\ref{topologies-definition-big-small-syntomic},
\ref{topologies-definition-big-small-smooth},
\ref{topologies-definition-big-small-etale}, and
\ref{topologies-definition-big-small-Zariski}.
\item If $\tau \in \{\etale, Zariski\}$, then the
{\it small $\tau$-site of $S$}
is the full subcategory $S_\tau$ of $(\Sch/S)_\tau$ whose objects
are schemes $T$ over $S$ whose structure morphism $T \to S$ is \'etale,
resp.\ an open immersion. A covering in $S_\tau$ is a covering
$\{U_i \to U\}$ in $(\Sch/S)_\tau$
such that $U$ is an object of $S_\tau$.
\end{enumerate}
\end{definition}

\noindent
The underlying category of the site $(\Sch/S)_\tau$ has reasonable
``closure'' properties, i.e., given a scheme $T$ in it any locally closed
subscheme of $T$ is isomorphic to an object of $(\Sch/S)_\tau$.
Other such closure properties are: closed under fibre products of schemes,
taking countable disjoint unions,
taking finite type schemes over a given scheme, given an affine scheme
$\Spec(R)$ one can complete, localize, or take the quotient of $R$
by an ideal while staying inside the category, etc.
On the other hand, for example arbitrary disjoint unions
of schemes in $(\Sch/S)_\tau$ will take you outside of it.
Also note that, given an object $T$ of $(\Sch/S)_\tau$ there will exist
$\tau$-coverings $\{T_i \to T\}_{i \in I}$ (as in
Definition \ref{definition-tau-covering})
which are not coverings in $(\Sch/S)_\tau$ for example
because the schemes $T_i$ are not objects of the category
$(\Sch/S)_\tau$. But our choice of the sites $(\Sch/S)_\tau$
is such that there always does exist
a covering $\{U_j \to T\}_{j \in J}$ of $(\Sch/S)_\tau$ which refines
the covering $\{T_i \to T\}_{i \in I}$, see
Topologies, Lemmas
\ref{topologies-lemma-fppf-induced},
\ref{topologies-lemma-syntomic-induced},
\ref{topologies-lemma-smooth-induced},
\ref{topologies-lemma-etale-induced}, and
\ref{topologies-lemma-zariski-induced}.
We will mostly ignore these issues in this chapter.

\medskip\noindent
If $\mathcal{F}$ is a sheaf on $(\Sch/S)_\tau$ or $S_\tau$, then
we denote
$$
H^p_\tau(U, \mathcal{F}), \text{ in particular }
H^p_\tau(S, \mathcal{F})
$$
the cohomology groups of $\mathcal{F}$ over the object $U$ of the site, see
Section \ref{section-cohomology}. Thus we have
$H^p_{fppf}(S, \mathcal{F})$,
$H^p_{syntomic}(S, \mathcal{F})$,
$H^p_{smooth}(S, \mathcal{F})$,
$H^p_\etale(S, \mathcal{F})$, and
$H^p_{Zar}(S, \mathcal{F})$. The last two are potentially ambiguous since
they might refer to either the big or small \'etale or Zariski site. However,
this ambiguity is harmless by the following lemma.

\begin{lemma}
\label{lemma-compare-cohomology-big-small}
Let $\tau \in \{\etale, Zariski\}$.
If $\mathcal{F}$ is an abelian sheaf defined on
$(\Sch/S)_\tau$, then
the cohomology groups of $\mathcal{F}$ over $S$ agree with the cohomology
groups of $\mathcal{F}|_{S_\tau}$ over $S$.
\end{lemma}

\begin{proof}
By
Topologies, Lemmas \ref{topologies-lemma-at-the-bottom} and
\ref{topologies-lemma-at-the-bottom-etale}
the functors $S_\tau \to (\Sch/S)_\tau$
satisfy the hypotheses of
Sites, Lemma \ref{sites-lemma-bigger-site}.
Hence our lemma follows from
Cohomology on Sites, Lemma \ref{sites-cohomology-lemma-cohomology-bigger-site}.
\end{proof}

\noindent
The category of sheaves on the big or small \'etale site of $S$ depends only
on the full subcategory of $(\Sch/S)_\etale$ or $S_\etale$ consisting of
affines and one only needs to consider the standard \'etale coverings between
them (as defined below). This gives rise to sites
$(\textit{Aff}/S)_\etale$ and $S_{affine, \etale}$, see
Topologies, Definition \ref{topologies-definition-big-small-etale}.
The comparison results are proven in Topologies,
Lemmas \ref{topologies-lemma-affine-big-site-etale} and
\ref{topologies-lemma-alternative}. Here is our definition of
standard coverings in some of the topologies we will consider in this chapter.

\begin{definition}
\label{definition-standard-tau}
(See
Topologies, Definitions
\ref{topologies-definition-standard-fppf},
\ref{topologies-definition-standard-syntomic},
\ref{topologies-definition-standard-smooth},
\ref{topologies-definition-standard-etale}, and
\ref{topologies-definition-standard-Zariski}.)
Let $\tau \in \{fppf, syntomic, smooth, \etale, Zariski\}$.
Let $T$ be an affine scheme.
A {\it standard $\tau$-covering} of $T$ is a family
$\{f_j : U_j \to T\}_{j = 1, \ldots, m}$ with each $U_j$ is affine,
and each $f_j$ flat and of finite presentation,
standard syntomic, standard smooth, \'etale, resp.\ the immersion of a
standard principal open in $T$ and $T = \bigcup f_j(U_j)$.
\end{definition}

\begin{lemma}
\label{lemma-tau-affine}
Let $\tau \in \{fppf, syntomic, smooth, \etale, Zariski\}$.
Any $\tau$-covering of an affine scheme can be refined by a
standard $\tau$-covering.
\end{lemma}

\begin{proof}
See
Topologies, Lemmas
\ref{topologies-lemma-fppf-affine},
\ref{topologies-lemma-syntomic-affine},
\ref{topologies-lemma-smooth-affine},
\ref{topologies-lemma-etale-affine}, and
\ref{topologies-lemma-zariski-affine}.
\end{proof}

\noindent
For completeness we state and prove the invariance under choice of partial
universe of the cohomology groups we are considering. We will prove invariance
of the small \'etale topos in
Lemma \ref{lemma-etale-topos-independent-partial-universe} below.
For notation and terminology used in this lemma we refer to
Topologies, Section \ref{topologies-section-change-alpha}.

\begin{lemma}
\label{lemma-cohomology-enlarge-partial-universe}
Let $\tau \in \{fppf, syntomic, smooth, \etale, Zariski\}$.
Let $S$ be a scheme.
Let $(\Sch/S)_\tau$ and $(\Sch'/S)_\tau$ be two
big $\tau$-sites of $S$, and assume that the first is contained in the second.
In this case
\begin{enumerate}
\item for any abelian sheaf $\mathcal{F}'$ defined on $(\Sch'/S)_\tau$ and
any object $U$ of $(\Sch/S)_\tau$ we have
$$
H^p_\tau(U, \mathcal{F}'|_{(\Sch/S)_\tau}) =
H^p_\tau(U, \mathcal{F}')
$$
In words: the cohomology of $\mathcal{F}'$ over $U$ computed in the bigger site
agrees with the cohomology of $\mathcal{F}'$ restricted to the smaller site
over $U$.
\item for any abelian sheaf $\mathcal{F}$ on $(\Sch/S)_\tau$ there is an
abelian sheaf $\mathcal{F}'$ on $(\Sch/S)_\tau'$ whose restriction to
$(\Sch/S)_\tau$ is isomorphic to $\mathcal{F}$.
\end{enumerate}
\end{lemma}

\begin{proof}
By Topologies, Lemma \ref{topologies-lemma-change-alpha} the inclusion functor
$(\Sch/S)_\tau \to (\Sch'/S)_\tau$ satisfies the assumptions of
Sites, Lemma \ref{sites-lemma-bigger-site}. This implies (2) and (1)
follows from
Cohomology on Sites, Lemma \ref{sites-cohomology-lemma-cohomology-bigger-site}.
\end{proof}




\section{The \'etale topos}
\label{section-etale-topos}

\noindent
A {\it topos} is the category of sheaves of sets on a site, see
Sites, Definition \ref{sites-definition-topos}. Hence it is customary
to refer to the use the phrase ``\'etale topos of a scheme'' to refer to
the category of sheaves on the small \'etale site of a scheme.
Here is the formal definition.

\begin{definition}
\label{definition-etale-topos}
Let $S$ be a scheme.
\begin{enumerate}
\item The {\it \'etale topos}, or the {\it small \'etale topos}
of $S$ is the category $\Sh(S_\etale)$ of sheaves of sets on
the small \'etale site of $S$.
\item The {\it Zariski topos}, or the {\it small Zariski topos}
of $S$ is the category $\Sh(S_{Zar})$ of sheaves of sets on the
small Zariski site of $S$.
\item For $\tau \in \{fppf, syntomic, smooth, \etale, Zariski\}$ a
{\it big $\tau$-topos} is the category of sheaves of set on a
big $\tau$-topos of $S$.
\end{enumerate}
\end{definition}

\noindent
Note that the small Zariski topos of $S$ is simply the category of sheaves
of sets on the underlying topological space of $S$, see
Topologies, Lemma \ref{topologies-lemma-Zariski-usual}.
Whereas the small \'etale topos does not depend on the choices made in the
construction of the small \'etale site, in general the big topoi do depend
on those choices.

\medskip\noindent
It turns out that the big or small \'etale topos only depends on the
full subcategory of $(\Sch/S)_\etale$ or $S_\etale$ consisting of
affines, see
Topologies, Lemmas \ref{topologies-lemma-affine-big-site-etale} and
\ref{topologies-lemma-alternative}.
We will use this for example in the proof of the following lemma.

\begin{lemma}
\label{lemma-etale-topos-independent-partial-universe}
Let $S$ be a scheme. The \'etale topos of $S$ is independent
(up to canonical equivalence) of the construction of the small
\'etale site in Definition \ref{definition-tau-site}.
\end{lemma}

\begin{proof}
We have to show, given two big \'etale sites
$\Sch_\etale$ and $\Sch_\etale'$ containing
$S$, then $\Sh(S_\etale) \cong \Sh(S_\etale')$
with obvious notation. By Topologies, Lemma \ref{topologies-lemma-contained-in}
we may assume $\Sch_\etale \subset \Sch_\etale'$.
By Sets, Lemma \ref{sets-lemma-what-is-in-it}
any affine scheme \'etale over $S$ is isomorphic to an object
of both $\Sch_\etale$ and $\Sch_\etale'$.
Thus the induced functor
$S_{affine, \etale} \to S_{affine, \etale}'$
is an equivalence. Moreover, it is clear that both this functor
and a quasi-inverse map transform standard \'etale coverings into
standard \'etale coverings.
Hence the result follows from
Topologies, Lemma \ref{topologies-lemma-alternative}.
\end{proof}





\section{Cohomology of quasi-coherent sheaves}
\label{section-cohomology-quasi-coherent}
%9.22.09

\noindent
We start with a simple lemma (which holds in greater generality than
stated). It says that the {\v C}ech complex of a standard covering is
equal to the {\v C}ech complex of an fpqc covering of the form
$\{\Spec(B) \to \Spec(A)\}$ with $A \to B$ faithfully flat.

\begin{lemma}
\label{lemma-cech-complex}
Let $\tau \in \{fppf, syntomic, smooth, \etale, Zariski\}$.
Let $S$ be a scheme.
Let $\mathcal{F}$ be an abelian sheaf on $(\Sch/S)_\tau$, or on
$S_\tau$ in case $\tau = \etale$, and let
$\mathcal{U} = \{U_i \to U\}_{i \in I}$
be a standard $\tau$-covering of this site.
Let $V = \coprod_{i \in I} U_i$. Then
\begin{enumerate}
\item $V$ is an affine scheme,
\item $\mathcal{V} = \{V \to U\}$ is an fpqc covering
and also a $\tau$-covering unless $\tau = Zariski$,
\item the {\v C}ech complexes
$\check{\mathcal{C}}^\bullet (\mathcal{U}, \mathcal{F})$ and
$\check{\mathcal{C}}^\bullet (\mathcal{V}, \mathcal{F})$ agree.
\end{enumerate}
\end{lemma}

\begin{proof}
The definition of a standard $\tau$-covering is given in
Topologies, Definition
\ref{topologies-definition-standard-Zariski},
\ref{topologies-definition-standard-etale},
\ref{topologies-definition-standard-smooth},
\ref{topologies-definition-standard-syntomic}, and
\ref{topologies-definition-standard-fppf}.
By definition each of the schemes
$U_i$ is affine and $I$ is a finite set. Hence $V$ is an affine scheme.
It is clear that $V \to U$ is flat and surjective, hence
$\mathcal{V}$ is an fpqc covering, see
Example \ref{example-fpqc-coverings}.
Excepting the Zariski case, the covering $\mathcal{V}$
is also a $\tau$-covering, see
Topologies, Definition
\ref{topologies-definition-etale-covering},
\ref{topologies-definition-smooth-covering},
\ref{topologies-definition-syntomic-covering}, and
\ref{topologies-definition-fppf-covering}.

\medskip\noindent
Note that $\mathcal{U}$ is a refinement of $\mathcal{V}$
and hence there is a map of {\v C}ech complexes
$\check{\mathcal{C}}^\bullet (\mathcal{V}, \mathcal{F}) \to
\check{\mathcal{C}}^\bullet (\mathcal{U}, \mathcal{F})$, see
Cohomology on Sites,
Equation (\ref{sites-cohomology-equation-map-cech-complexes}).
Next, we observe that if $T = \coprod_{j \in J} T_j$ is a
disjoint union of schemes in the site on which $\mathcal{F}$ is defined
then the family of morphisms with fixed target
$\{T_j \to T\}_{j \in J}$ is a Zariski covering, and so
\begin{equation}
\label{equation-sheaf-coprod}
\mathcal{F}(T) =
\mathcal{F}(\coprod\nolimits_{j \in J} T_j) =
\prod\nolimits_{j \in J} \mathcal{F}(T_j)
\end{equation}
by the sheaf condition of $\mathcal{F}$.
This implies the map of {\v C}ech complexes above is an isomorphism
in each degree because
$$
V \times_U \ldots \times_U V
=
\coprod\nolimits_{i_0, \ldots i_p} U_{i_0} \times_U \ldots \times_U U_{i_p}
$$
as schemes.
\end{proof}

\noindent
Note that Equality (\ref{equation-sheaf-coprod})
is false for a general presheaf. Even for sheaves it does not hold on any
site, since coproducts may not lead to coverings, and may not be disjoint.
But it does for all the usual ones (at least all the ones we will study).

\begin{remark}
\label{remark-refinement}
In the statement of Lemma \ref{lemma-cech-complex} the covering $\mathcal{U}$
is a refinement of $\mathcal{V}$ but not the other way around. Coverings
of the form $\{V \to U\}$ do not form an initial subcategory of the
category of all coverings of $U$. Yet it is still true that
we can compute {\v C}ech cohomology $\check H^n(U, \mathcal{F})$ (which
is defined as the colimit over the opposite of the category of
coverings $\mathcal{U}$ of $U$ of the {\v C}ech cohomology groups of
$\mathcal{F}$ with respect to $\mathcal{U}$) in terms of the coverings
$\{V \to U\}$. We will formulate a precise lemma (it only works for sheaves)
and add it here if we ever need it.
\end{remark}

\begin{lemma}[Locality of cohomology]
\label{lemma-locality-cohomology}
Let $\mathcal{C}$ be a site, $\mathcal{F}$ an abelian sheaf on $\mathcal{C}$,
$U$ an object of $\mathcal{C}$, $p > 0$ an integer and $\xi \in
H^p(U, \mathcal{F})$. Then there exists a covering
$\mathcal{U} = \{U_i \to U\}_{i \in I}$ of $U$ in $\mathcal{C}$
such that $\xi |_{U_i} = 0$ for all $i \in I$.
\end{lemma}

\begin{proof}
Choose an injective resolution $\mathcal{F} \to \mathcal{I}^\bullet$. Then
$\xi$ is represented by a cocycle $\tilde{\xi} \in \mathcal{I}^p(U)$
with $d^p(\tilde{\xi}) = 0$. By assumption, the sequence
$\mathcal{I}^{p - 1} \to \mathcal{I}^p \to \mathcal{I}^{p + 1}$ in exact in
$\textit{Ab}(\mathcal{C})$, which means that there exists a covering
$\mathcal{U} = \{U_i \to U\}_{i \in I}$ such that
$\tilde{\xi}|_{U_i} = d^{p - 1}(\xi_i)$ for some
$\xi_i \in \mathcal{I}^{p-1}(U_i)$. Since
the cohomology class $\xi|_{U_i}$ is represented by the cocycle
$\tilde{\xi}|_{U_i}$ which is a coboundary, it vanishes.
For more details see
Cohomology on Sites,
Lemma \ref{sites-cohomology-lemma-kill-cohomology-class-on-covering}.
\end{proof}

\begin{theorem}
\label{theorem-zariski-fpqc-quasi-coherent}
Let $S$ be a scheme and $\mathcal{F}$ a quasi-coherent $\mathcal{O}_S$-module.
Let $\mathcal{C}$ be either $(\Sch/S)_\tau$ for
$\tau \in \{fppf, syntomic, smooth, \etale, Zariski\}$ or
$S_\etale$. Then
$$
H^p(S, \mathcal{F}) = H^p_\tau(S, \mathcal{F}^a)
$$
for all $p \geq 0$ where
\begin{enumerate}
\item the left hand side indicates the usual cohomology of the sheaf
$\mathcal{F}$ on the underlying topological space of the scheme $S$, and
\item the right hand side indicates cohomology
of the abelian sheaf $\mathcal{F}^a$ (see
Proposition \ref{proposition-quasi-coherent-sheaf-fpqc})
on the site $\mathcal{C}$.
\end{enumerate}
\end{theorem}

\begin{proof}
We are going to show that
$H^p(U, f^*\mathcal{F}) = H^p_\tau(U, \mathcal{F}^a)$
for any object $f : U \to S$ of the site $\mathcal{C}$.
The result is true for $p = 0$ by the sheaf property.

\medskip\noindent
Assume that $U$ is affine. Then we want to prove that
$H^p_\tau(U, \mathcal{F}^a) = 0$ for all $p > 0$. We use induction on $p$.
\begin{enumerate}
\item[p = 1]
Pick $\xi \in H^1_\tau(U, \mathcal{F}^a)$.
By Lemma \ref{lemma-locality-cohomology},
there exists an fpqc covering $\mathcal{U} = \{U_i \to U\}_{i \in I}$
such that $\xi|_{U_i} = 0$ for all $i \in I$. Up to refining
$\mathcal{U}$, we may assume that $\mathcal{U}$ is a standard
$\tau$-covering. Applying the spectral sequence of
Theorem \ref{theorem-cech-ss},
we see that $\xi$ comes from a cohomology class
$\check \xi \in \check H^1(\mathcal{U}, \mathcal{F}^a)$.
Consider the covering $\mathcal{V} = \{\coprod_{i\in I} U_i \to U\}$. By
Lemma \ref{lemma-cech-complex},
$\check H^\bullet(\mathcal{U}, \mathcal{F}^a) =
\check H^\bullet(\mathcal{V}, \mathcal{F}^a)$.
On the other hand, since $\mathcal{V}$ is a covering of the form
$\{\Spec(B) \to \Spec(A)\}$ and $f^*\mathcal{F} = \widetilde{M}$
for some $A$-module $M$, we see the {\v C}ech complex
$\check{\mathcal{C}}^\bullet(\mathcal{V}, \mathcal{F})$
is none other than the complex $(B/A)_\bullet \otimes_A M$.
Now by Lemma \ref{lemma-descent-modules},
$H^p((B/A)_\bullet \otimes_A M) = 0$ for $p > 0$, hence $\check \xi = 0$
and so $\xi = 0$.
\item[p > 1]
Pick $\xi \in H^p_\tau(U, \mathcal{F}^a)$. By
Lemma \ref{lemma-locality-cohomology},
there exists an fpqc covering $\mathcal{U} = \{U_i \to U\}_{i \in I}$
such that $\xi|_{U_i} = 0$ for all $i \in I$. Up to refining
$\mathcal{U}$, we may assume that $\mathcal{U}$ is a standard
$\tau$-covering. We apply the spectral sequence of
Theorem \ref{theorem-cech-ss}.
Observe that the intersections $U_{i_0} \times_U \ldots \times_U U_{i_p}$
are affine, so that by induction hypothesis the cohomology groups
$$
E_2^{p, q} = \check H^p(\mathcal{U}, \underline{H}^q(\mathcal{F}^a))
$$
vanish for all $0 < q < p$. We see that $\xi$ must come from a
$\check \xi \in \check H^p(\mathcal{U}, \mathcal{F}^a)$. Replacing
$\mathcal{U}$ with the covering $\mathcal{V}$ containing only one morphism
and using Lemma \ref{lemma-descent-modules} again,
we see that the {\v C}ech cohomology class $\check \xi$ must be zero,
hence $\xi = 0$.
\end{enumerate}
Next, assume that $U$ is separated. Choose an affine open covering
$U = \bigcup_{i \in I} U_i$ of $U$. The family
$\mathcal{U} = \{U_i \to U\}_{i \in I}$ is then an fpqc covering,
and all the intersections
$U_{i_0} \times_U \ldots \times_U U_{i_p}$ are affine
since $U$ is separated. So all rows of the spectral sequence of
Theorem \ref{theorem-cech-ss}
are zero, except the zeroth row. Therefore
$$
H^p_\tau(U, \mathcal{F}^a) =
\check H^p(\mathcal{U}, \mathcal{F}^a) =
\check H^p(\mathcal{U}, \mathcal{F}) = H^p(U, \mathcal{F})
$$
where the last equality results from standard scheme theory, see
Cohomology of Schemes, Lemma
\ref{coherent-lemma-cech-cohomology-quasi-coherent}.

\medskip\noindent
The general case is technical and (to extend the proof as given here)
requires a discussion about maps of spectral sequences, so we won't treat it.
It follows from
Descent, Proposition \ref{descent-proposition-same-cohomology-quasi-coherent}
(whose proof takes a slightly different approach) combined with
Cohomology on Sites, Lemma \ref{sites-cohomology-lemma-cohomology-of-open}.
\end{proof}

\begin{remark}
\label{remark-right-derived-global-sections}
Comment on Theorem \ref{theorem-zariski-fpqc-quasi-coherent}.
Since $S$ is a final object in the category $\mathcal{C}$, the cohomology
groups on the right-hand side are merely the right derived functors of the
global sections functor. In fact the proof shows that
$H^p(U, f^*\mathcal{F}) = H^p_\tau(U, \mathcal{F}^a)$
for any object $f : U \to S$ of the site $\mathcal{C}$.
\end{remark}





\section{Examples of sheaves}
\label{section-examples-sheaves}

\noindent
Let $S$ and $\tau$ be as in Section \ref{section-big-small}.
We have already seen that any representable presheaf is a sheaf on
$(\Sch/S)_\tau$ or $S_\tau$, see
Lemma \ref{lemma-representable-sheaf-fpqc}
and
Remark \ref{remark-fpqc-finest}.
Here are some special cases.

\begin{definition}
\label{definition-additive-sheaf}
On any of the sites $(\Sch/S)_\tau$ or $S_\tau$ of
Section \ref{section-big-small}.
\begin{enumerate}
\item The sheaf $T \mapsto \Gamma(T, \mathcal{O}_T)$ is denoted
$\mathcal{O}_S$, or $\mathbf{G}_a$, or $\mathbf{G}_{a, S}$ if we
want to indicate the base scheme.
\item Similarly, the sheaf
$T \mapsto \Gamma(T, \mathcal{O}^*_T)$ is denoted $\mathcal{O}_S^*$, or
$\mathbf{G}_m$, or $\mathbf{G}_{m, S}$ if we want
to indicate the base scheme.
\item The {\it constant sheaf} $\underline{\mathbf{Z}/n\mathbf{Z}}$ on any
site is the sheafification of the constant presheaf
$U \mapsto \mathbf{Z}/n\mathbf{Z}$.
\end{enumerate}
\end{definition}

\noindent
The first is a sheaf by
Theorem \ref{theorem-quasi-coherent}
for example. The second is a sub presheaf of the first, which is easily seen
to be a sheaf itself. The third is a sheaf by definition.
Note that each of these sheaves is representable.
The first and second by the schemes $\mathbf{G}_{a, S}$ and
$\mathbf{G}_{m, S}$, see
Groupoids, Section \ref{groupoids-section-group-schemes}.
The third by the finite \'etale group scheme $\mathbf{Z}/n\mathbf{Z}_S$
sometimes denoted $(\mathbf{Z}/n\mathbf{Z})_S$
which is just $n$ copies of $S$ endowed
with the obvious group scheme structure over $S$, see
Groupoids, Example \ref{groupoids-example-constant-group}
and the following remark.

\begin{remark}
\label{remark-constant-locally-constant-maps}
Let $G$ be an abstract group.
On any of the sites $(\Sch/S)_\tau$ or $S_\tau$ of
Section \ref{section-big-small}
the sheafification $\underline{G}$
of the constant presheaf associated to $G$ in the
{\it Zariski topology} of the site already gives
$$
\Gamma(U, \underline{G}) =
\{\text{Zariski locally constant maps }U \to G\}
$$
This Zariski sheaf is representable by the group scheme $G_S$ according to
Groupoids, Example \ref{groupoids-example-constant-group}.
By
Lemma \ref{lemma-representable-sheaf-fpqc}
any representable presheaf satisfies the sheaf condition for the
$\tau$-topology as well, and hence we conclude that the Zariski
sheafification $\underline{G}$ above is also the $\tau$-sheafification.
\end{remark}

\begin{definition}
\label{definition-structure-sheaf}
Let $S$ be a scheme. The {\it structure sheaf} of $S$ is the sheaf of rings
$\mathcal{O}_S$
on any of the sites $S_{Zar}$, $S_\etale$, or $(\Sch/S)_\tau$
discussed above.
\end{definition}

\noindent
If there is some possible confusion as to which site we are working on
then we will indicate this by using indices. For example we may use
$\mathcal{O}_{S_\etale}$ to stress the fact that we are working on the
small \'etale site of $S$.

\begin{remark}
\label{remark-special-case-fpqc-cohomology-quasi-coherent}
In the terminology introduced above a special case of
Theorem \ref{theorem-zariski-fpqc-quasi-coherent}
is
$$
H_{fppf}^p(X, \mathbf{G}_a) =
H_\etale^p(X, \mathbf{G}_a) =
H_{Zar}^p(X, \mathbf{G}_a) =
H^p(X, \mathcal{O}_X)
$$
for all $p \geq 0$. Moreover, we could use the notation
$H^p_{fppf}(X, \mathcal{O}_X)$ to indicate the cohomology of the
structure sheaf on the big fppf site of $X$.
\end{remark}




\section{Picard groups}
\label{section-picard-groups}

\noindent
The following theorem is sometimes called ``Hilbert 90''.

\begin{theorem}
\label{theorem-picard-group}
For any scheme $X$ we have canonical identifications
\begin{align*}
H_{fppf}^1(X, \mathbf{G}_m) & = H^1_{syntomic}(X, \mathbf{G}_m) \\
& = H^1_{smooth}(X, \mathbf{G}_m) \\
& = H_\etale^1(X, \mathbf{G}_m) \\
& = H^1_{Zar}(X, \mathbf{G}_m) \\
& = \Pic(X) \\
& = H^1(X, \mathcal{O}_X^*)
\end{align*}
\end{theorem}

\begin{proof}
Let $\tau$ be one of the topologies considered in
Section \ref{section-big-small}.
By
Cohomology on Sites, Lemma
\ref{sites-cohomology-lemma-h1-invertible}
we see that
$H^1_\tau(X, \mathbf{G}_m) =
H^1_\tau(X, \mathcal{O}_\tau^*) =
\Pic(\mathcal{O}_\tau)$
where $\mathcal{O}_\tau$ is the structure sheaf of the site
$(\Sch/X)_\tau$. Now an invertible $\mathcal{O}_\tau$-module
is a quasi-coherent $\mathcal{O}_\tau$-module.
By Theorem \ref{theorem-quasi-coherent} or the more precise
Descent, Proposition \ref{descent-proposition-equivalence-quasi-coherent}
we see that $\Pic(\mathcal{O}_\tau) = \Pic(X)$.
The last equality is proved in the same way.
\end{proof}






\section{The \'etale site}
\label{section-etale-site}

\noindent
At this point we start exploring the \'etale site of a scheme in
more detail. As a first step we discuss a little the notion of an
\'etale morphism.





\section{\'Etale morphisms}
\label{section-etale-morphism}

\noindent
For more details, see
Morphisms, Section \ref{morphisms-section-etale}
for the formal definition and
\'Etale Morphisms, Sections
\ref{etale-section-etale-morphisms},
\ref{etale-section-structure-etale-map},
\ref{etale-section-etale-smooth},
\ref{etale-section-topological-etale},
\ref{etale-section-functorial-etale}, and
\ref{etale-section-properties-permanence}
for a survey of interesting properties of \'etale morphisms.

\medskip\noindent
Recall that an algebra $A$ over an algebraically closed field $k$ is
{\it smooth} if it is of finite type and the module of differentials
$\Omega_{A/k}$ is finite locally free of rank equal to the dimension.
A scheme $X$ over $k$ is {\it smooth} over $k$ if it is locally of finite
type and each affine open is the spectrum of a smooth $k$-algebra.
If $k$ is not algebraically closed then a $k$-algebra $A$ is
a smooth $k$-algebra if $A \otimes_k \overline{k}$ is a smooth
$\overline{k}$-algebra. A ring map $A \to B$ is smooth if it is
flat, finitely presented, and for all primes $\mathfrak p \subset A$
the fibre ring $\kappa(\mathfrak p) \otimes_A B$ is smooth over the residue
field $\kappa(\mathfrak p)$. More generally, a morphism of schemes is
{\it smooth} if it is flat, locally of finite presentation, and the
geometric fibers are smooth.

\medskip\noindent
For these facts please see
Morphisms, Section \ref{morphisms-section-smooth}.
Using this we may define an \'etale morphism as follows.

\begin{definition}
\label{definition-etale-morphism}
A morphism of schemes is {\it \'etale} if it is smooth of relative dimension 0.
\end{definition}

\noindent
In particular, a morphism of schemes $X \to S$ is \'etale if it is smooth
and $\Omega_{X/S} = 0$.

\begin{proposition}
\label{proposition-etale-morphisms}
Facts on \'etale morphisms.
\begin{enumerate}
\item Let $k$ be a field. A morphism of schemes $U \to \Spec(k)$ is
\'etale if and only if $U \cong \coprod_{i \in I} \Spec(k_i)$
such that for each $i \in I$
the ring $k_i$ is a field which is a finite separable extension of $k$.
\item Let $\varphi : U \to S$ be a morphism of schemes. The following
conditions are equivalent:
\begin{enumerate}
\item $\varphi$ is \'etale,
\item $\varphi$ is locally finitely presented, flat, and all its fibres are
\'etale,
\item $\varphi$ is flat, unramified and locally of finite presentation.
\end{enumerate}
\item A ring map $A \to B$ is \'etale if and only if
$B \cong A[x_1, \ldots, x_n]/(f_1, \ldots, f_n)$
such that $\Delta = \det \left( \frac{\partial f_i}{\partial x_j} \right)$
is invertible in $B$.
\item The base change of an \'etale morphism is \'etale.
\item Compositions of \'etale morphisms are \'etale.
\item Fibre products and products of \'etale morphisms are \'etale.
\item An \'etale morphism has relative dimension 0.
\item Let $Y \to X$ be an \'etale morphism.
If $X$ is reduced (respectively regular) then so is $Y$.
\item \'Etale morphisms are open.
\item If $X \to S$ and $Y \to S$ are \'etale, then any
$S$-morphism $X \to Y$ is also \'etale.
\end{enumerate}
\end{proposition}

\begin{proof}
We have proved these facts (and more) in the preceding chapters.
Here is a list of references:
(1) Morphisms, Lemma \ref{morphisms-lemma-etale-over-field}.
(2) Morphisms, Lemmas \ref{morphisms-lemma-etale-flat-etale-fibres}
and \ref{morphisms-lemma-flat-unramified-etale}.
(3) Algebra, Lemma \ref{algebra-lemma-etale-standard-smooth}.
(4) Morphisms, Lemma \ref{morphisms-lemma-base-change-etale}.
(5) Morphisms, Lemma \ref{morphisms-lemma-composition-etale}.
(6) Follows formally from (4) and (5).
(7) Morphisms, Lemmas \ref{morphisms-lemma-etale-locally-quasi-finite}
and \ref{morphisms-lemma-locally-quasi-finite-rel-dimension-0}.
(8) See Algebra, Lemmas \ref{algebra-lemma-reduced-goes-up} and
\ref{algebra-lemma-Rk-goes-up}, see also more results of this kind
in \'Etale Morphisms, Section \ref{etale-section-properties-permanence}.
(9) See Morphisms, Lemma \ref{morphisms-lemma-fppf-open} and
\ref{morphisms-lemma-etale-flat}.
(10) See Morphisms, Lemma \ref{morphisms-lemma-etale-permanence}.
\end{proof}

\begin{definition}
\label{definition-standard-etale}
A ring map $A \to B$ is called {\it standard \'etale} if
$B \cong \left(A[t]/(f)\right)_g$ with $f, g \in A[t]$, with $f$ monic,
and $\text{d}f/\text{d}t$ invertible in $B$.
\end{definition}

\noindent
It is true that a standard \'etale ring map is \'etale. Namely, suppose
that $B = \left(A[t]/(f)\right)_g$ with $f, g \in A[t]$, with $f$ monic,
and $\text{d}f/\text{d}t$ invertible in $B$. Then $A[t]/(f)$ is a finite
free $A$-module of rank equal to the degree of the monic polynomial $f$.
Hence $B$, as a localization of this free algebra is finitely presented
and flat over $A$. To finish the proof that $B$ is \'etale it suffices
to show that the fibre rings
$$
\kappa(\mathfrak p) \otimes_A B
\cong
\kappa(\mathfrak p) \otimes_A (A[t]/(f))_g
\cong
\kappa(\mathfrak p)[t, 1/\overline{g}]/(\overline{f})
$$
are finite products of finite separable field extensions.
Here $\overline{f}, \overline{g} \in \kappa(\mathfrak p)[t]$ are
the images of $f$ and $g$. Let
$$
\overline{f} = \overline{f}_1 \ldots \overline{f}_a
\overline{f}_{a + 1}^{e_1} \ldots \overline{f}_{a + b}^{e_b}
$$
be the factorization of $\overline{f}$ into powers of pairwise distinct
irreducible monic factors $\overline{f}_i$ with $e_1, \ldots, e_b > 1$.
By assumption $\text{d}\overline{f}/\text{d}t$ is invertible in
$\kappa(\mathfrak p)[t, 1/\overline{g}]$. Hence we see that
at least all the $\overline{f}_i$, $i > a$ are invertible. We conclude
that
$$
\kappa(\mathfrak p)[t, 1/\overline{g}]/(\overline{f})
\cong
\prod\nolimits_{i \in I} \kappa(\mathfrak p)[t]/(\overline{f}_i)
$$
where $I \subset \{1, \ldots, a\}$ is the subset of indices $i$ such that
$\overline{f}_i$ does not divide $\overline{g}$. Moreover, the image of
$\text{d}\overline{f}/\text{d}t$ in the factor
$\kappa(\mathfrak p)[t]/(\overline{f}_i)$ is clearly equal to a
unit times $\text{d}\overline{f}_i/\text{d}t$. Hence we conclude that
$\kappa_i = \kappa(\mathfrak p)[t]/(\overline{f}_i)$ is a finite field
extension of $\kappa(\mathfrak p)$ generated by one element whose
minimal polynomial is separable, i.e., the field extension
$\kappa_i/\kappa(\mathfrak p)$ is finite separable as desired.

\medskip\noindent
It turns out that any \'etale ring map is locally standard \'etale.
To formulate this we introduce the following notation.
A ring map $A \to B$ is {\it \'etale at a prime $\mathfrak q$} of $B$ if there
exists $h \in B$, $h \not \in \mathfrak q$ such that $A \to B_h$ is \'etale.
Here is the result.

\begin{theorem}
\label{theorem-standard-etale}
A ring map $A \to B$ is \'etale at a prime $\mathfrak q$ if and only if there
exists $g \in B$, $g \not \in \mathfrak q$ such that $B_g$ is standard
\'etale over $A$.
\end{theorem}

\begin{proof}
See
Algebra, Proposition \ref{algebra-proposition-etale-locally-standard}.
\end{proof}





\section{\'Etale coverings}
\label{section-etale-covering}

\noindent
We recall the definition.

\begin{definition}
\label{definition-etale-covering}
An {\it \'etale covering} of a scheme $U$ is a family of morphisms
of schemes
$\{\varphi_i : U_i \to U\}_{i \in I}$ such that
\begin{enumerate}
\item each $\varphi_i$ is an \'etale morphism,
\item the $U_i$ cover $U$, i.e., $U = \bigcup_{i\in I}\varphi_i(U_i)$.
\end{enumerate}
\end{definition}

\begin{lemma}
\label{lemma-etale-fpqc}
Any \'etale covering is an fpqc covering.
\end{lemma}

\begin{proof}
(See also
Topologies,
Lemma \ref{topologies-lemma-zariski-etale-smooth-syntomic-fppf-fpqc}.)
Let $\{\varphi_i : U_i \to U\}_{i \in I}$ be an \'etale covering.
Since an \'etale morphism is flat, and the elements of the covering should
cover its target, the property fp (faithfully flat) is satisfied.
To check the property qc (quasi-compact), let $V \subset U$ be an affine
open, and write $\varphi_i^{-1}(V) = \bigcup_{j \in J_i} V_{ij}$
for some affine opens $V_{ij} \subset U_i$. Since $\varphi_i$ is open
(as \'etale morphisms are open), we see that
$V = \bigcup_{i\in I} \bigcup_{j \in J_i} \varphi_i(V_{ij})$
is an open covering of $V$.
Further, since $V$ is quasi-compact, this covering has a finite
refinement.
\end{proof}

\noindent
So any statement which is true for fpqc coverings
remains true {\it a fortiori} for \'etale coverings. For
instance, the \'etale site is subcanonical.

\begin{definition}
\label{definition-big-etale-site}
(For more details see Section \ref{section-big-small}, or
Topologies, Section \ref{topologies-section-etale}.)
Let $S$ be a scheme.
The {\it big \'etale site over $S$} is the site
$(\Sch/S)_\etale$, see
Definition \ref{definition-tau-site}.
The {\it small \'etale site over $S$} is the site $S_\etale$, see
Definition \ref{definition-tau-site}.
We define similarly the {\it big} and {\it small Zariski sites} on $S$,
denoted $(\Sch/S)_{Zar}$ and $S_{Zar}$.
\end{definition}

\noindent
Loosely speaking the big \'etale site of $S$ is made up out of schemes over $S$
and coverings the \'etale coverings. The small \'etale site of $S$ is made up
out of schemes \'etale over $S$ with coverings the \'etale coverings.
Actually any morphism between objects of $S_\etale$ is \'etale, in
virtue of
Proposition \ref{proposition-etale-morphisms},
hence to check that $\{U_i \to U\}_{i \in I}$ in $S_\etale$
is a covering it suffices to check that $\coprod U_i \to U$ is surjective.

\medskip\noindent
The small \'etale site has fewer objects than the big \'etale site, it
contains only the ``opens'' of the \'etale topology on $S$. It is a full
subcategory of the big \'etale site, and its topology is induced from the
topology on the big site. Hence it is true that the restriction functor
from the big \'etale site to the small one is exact and maps injectives to
injectives. This has the following consequence.

\begin{proposition}
\label{proposition-cohomology-restrict-small-site}
Let $S$ be a scheme and $\mathcal{F}$ an abelian sheaf on
$(\Sch/S)_\etale$.
Then $\mathcal{F}|_{S_\etale}$ is a sheaf on $S_\etale$ and
$$
H^p_\etale(S, \mathcal{F}|_{S_\etale}) =
H^p_\etale(S, \mathcal{F})
$$
for all $p \geq 0$.
\end{proposition}

\begin{proof}
This is a special case of Lemma \ref{lemma-compare-cohomology-big-small}.
\end{proof}

\noindent
In accordance with the general notation introduced in
Section \ref{section-big-small}
we write $H_\etale^p(S, \mathcal{F})$ for the above cohomology group.





%9.24.09
\section{Kummer theory}
\label{section-kummer}

\noindent
Let $n \in \mathbf{N}$ and consider the functor $\mu_n$ defined by
$$
\begin{matrix}
\Sch^{opp} & \longrightarrow & \textit{Ab} \\
S & \longmapsto &
\mu_n(S)
=
\{t \in \Gamma(S, \mathcal{O}_S^*) \mid t^n = 1 \}.
\end{matrix}
$$
By
Groupoids, Example \ref{groupoids-example-roots-of-unity}
this is a representable functor, and the scheme representing it
is denoted $\mu_n$ also. By
Lemma \ref{lemma-representable-sheaf-fpqc}
this functor satisfies the sheaf condition for the fpqc topology
(in particular, it also satisfies the sheaf condition for the
\'etale, Zariski, etc topology).

\begin{lemma}
\label{lemma-kummer-sequence}
If $n\in \mathcal{O}_S^*$ then
$$
0 \to
\mu_{n, S} \to
\mathbf{G}_{m, S} \xrightarrow{(\cdot)^n}
\mathbf{G}_{m, S} \to 0
$$
is a short exact sequence of sheaves on both the small and
big \'etale site of $S$.
\end{lemma}

\begin{proof}
By definition the sheaf $\mu_{n, S}$ is the kernel of the map
$(\cdot)^n$. Hence it suffices to show that the last map is surjective.
Let $U$ be a scheme over $S$. Let
$f \in \mathbf{G}_m(U) = \Gamma(U, \mathcal{O}_U^*)$.
We need to show that we can find an \'etale cover of
$U$ over the members of which the restriction of $f$ is an $n$th power.
Set
$$
U' =
\underline{\Spec}_U(\mathcal{O}_U[T]/(T^n-f))
\xrightarrow{\pi}
U.
$$
(See
Constructions, Section \ref{constructions-section-spec-via-glueing} or
\ref{constructions-section-spec}
for a discussion of the relative spectrum.)
Let $\Spec(A) \subset U$ be an affine open, and say $f|_{\Spec(A)}$ corresponds
to the unit $a \in A^*$. Then $\pi^{-1}(\Spec(A)) = \Spec(B)$ with
$B = A[T]/(T^n - a)$. The ring map $A \to B$ is finite free of rank $n$,
hence it is faithfully flat, and hence we conclude that
$\Spec(B) \to \Spec(A)$ is surjective. Since this holds for every
affine open in $U$ we conclude that $\pi$ is surjective.
In addition, $n$ and $T^{n - 1}$ are invertible in $B$, so
$nT^{n-1} \in B^*$ and the ring map $A \to B$ is standard \'etale,
in particular \'etale. Since this holds for every affine open of $U$
we conclude that $\pi$ is \'etale. Hence
$\mathcal{U} = \{\pi : U' \to U\}$ is an \'etale covering.
Moreover, $f|_{U'} = (f')^n$ where $f'$ is the class of $T$
in $\Gamma(U', \mathcal{O}_{U'}^*)$, so $\mathcal{U}$ has the desired property.
\end{proof}

\begin{remark}
\label{remark-no-kummer-sequence-zariski}
Lemma \ref{lemma-kummer-sequence} is false when ``\'etale'' is replaced
with ``Zariski''.
Since the \'etale topology is coarser than the smooth topology, see
Topologies, Lemma \ref{topologies-lemma-zariski-etale-smooth}
it follows that the sequence is also exact in the smooth topology.
\end{remark}

\noindent
By
Theorem \ref{theorem-picard-group}
and
Lemma \ref{lemma-kummer-sequence}
and general properties of cohomology we obtain
the long exact cohomology sequence
$$
\xymatrix{
0 \ar[r] &
H_\etale^0(S, \mu_{n, S}) \ar[r] &
\Gamma(S, \mathcal{O}_S^*) \ar[r]^{(\cdot)^n} &
\Gamma(S, \mathcal{O}_S^*) \ar@(rd, ul)[rdllllr]
\\
& H_\etale^1(S, \mu_{n, S}) \ar[r] &
\Pic(S) \ar[r]^{(\cdot)^n} &
\Pic(S) \ar@(rd, ul)[rdllllr] \\
& H_\etale^2(S, \mu_{n, S}) \ar[r] &
\ldots
}
$$
at least if $n$ is invertible on $S$. When $n$ is not invertible on $S$
we can apply the following lemma.

\begin{lemma}
\label{lemma-kummer-sequence-syntomic}
For any $n \in \mathbf{N}$ the sequence
$$
0 \to
\mu_{n, S} \to
\mathbf{G}_{m, S} \xrightarrow{(\cdot)^n}
\mathbf{G}_{m, S} \to 0
$$
is a short exact sequence of sheaves on the site
$(\Sch/S)_{fppf}$ and $(\Sch/S)_{syntomic}$.
\end{lemma}

\begin{proof}
By definition the sheaf $\mu_{n, S}$ is the kernel of the map
$(\cdot)^n$. Hence it suffices to show that the last map is surjective.
Since the syntomic topology is weaker than the fppf topology, see
Topologies, Lemma \ref{topologies-lemma-zariski-etale-smooth-syntomic-fppf},
it suffices to prove this for the syntomic topology.
Let $U$ be a scheme over $S$. Let
$f \in \mathbf{G}_m(U) = \Gamma(U, \mathcal{O}_U^*)$.
We need to show that we can find a syntomic cover of
$U$ over the members of which the restriction of $f$ is an $n$th power.
Set
$$
U' =
\underline{\Spec}_U(\mathcal{O}_U[T]/(T^n-f))
\xrightarrow{\pi}
U.
$$
(See
Constructions, Section \ref{constructions-section-spec-via-glueing} or
\ref{constructions-section-spec}
for a discussion of the relative spectrum.)
Let $\Spec(A) \subset U$ be an affine open, and say $f|_{\Spec(A)}$ corresponds
to the unit $a \in A^*$. Then $\pi^{-1}(\Spec(A)) = \Spec(B)$ with
$B = A[T]/(T^n - a)$. The ring map $A \to B$ is finite free of rank $n$,
hence it is faithfully flat, and hence we conclude that
$\Spec(B) \to \Spec(A)$ is surjective. Since this holds for every
affine open in $U$ we conclude that $\pi$ is surjective.
In addition, $B$ is a global relative complete intersection over $A$, so
the ring map $A \to B$ is standard syntomic,
in particular syntomic. Since this holds for every affine open of $U$
we conclude that $\pi$ is syntomic. Hence
$\mathcal{U} = \{\pi : U' \to U\}$ is a syntomic covering.
Moreover, $f|_{U'} = (f')^n$ where $f'$ is the class of $T$
in $\Gamma(U', \mathcal{O}_{U'}^*)$, so $\mathcal{U}$ has the desired property.
\end{proof}

\begin{remark}
\label{remark-no-kummer-sequence-smooth-etale-zariski}
Lemma \ref{lemma-kummer-sequence-syntomic}
is false for the smooth, \'etale, or Zariski topology.
\end{remark}

\noindent
By
Theorem \ref{theorem-picard-group}
and
Lemma \ref{lemma-kummer-sequence-syntomic}
and general properties of cohomology we obtain
the long exact cohomology sequence
$$
\xymatrix{
0 \ar[r] &
H_{fppf}^0(S, \mu_{n, S}) \ar[r] &
\Gamma(S, \mathcal{O}_S^*) \ar[r]^{(\cdot)^n} &
\Gamma(S, \mathcal{O}_S^*) \ar@(rd, ul)[rdllllr]
\\
& H_{fppf}^1(S, \mu_{n, S}) \ar[r] &
\Pic(S) \ar[r]^{(\cdot)^n} &
\Pic(S) \ar@(rd, ul)[rdllllr] \\
& H_{fppf}^2(S, \mu_{n, S}) \ar[r] &
\ldots
}
$$
for any scheme $S$ and any integer $n$. Of course there is a similar sequence
with syntomic cohomology.

\medskip\noindent
Let $n \in \mathbf{N}$ and let $S$ be any scheme.
There is another more direct way to describe the first cohomology group with
values in $\mu_n$. Consider pairs
$(\mathcal{L}, \alpha)$ where $\mathcal{L}$ is an invertible sheaf on $S$
and $\alpha : \mathcal{L}^{\otimes n} \to \mathcal{O}_S$ is a trivialization
of the $n$th tensor power of $\mathcal{L}$.
Let $(\mathcal{L}', \alpha')$ be a second such pair.
An isomorphism $\varphi : (\mathcal{L}, \alpha) \to (\mathcal{L}', \alpha')$
is an isomorphism $\varphi : \mathcal{L} \to \mathcal{L}'$ of invertible
sheaves such that the diagram
$$
\xymatrix{
\mathcal{L}^{\otimes n} \ar[d]_{\varphi^{\otimes n}} \ar[r]_\alpha &
\mathcal{O}_S \ar[d]^1 \\
(\mathcal{L}')^{\otimes n} \ar[r]^{\alpha'} &
\mathcal{O}_S \\
}
$$
commutes. Thus we have
\begin{equation}
\label{equation-isomorphisms-pairs}
\mathit{Isom}_S((\mathcal{L}, \alpha), (\mathcal{L}', \alpha'))
=
\left\{
\begin{matrix}
\emptyset & \text{if} & \text{they are not isomorphic} \\
H^0(S, \mu_{n, S})\cdot \varphi & \text{if} &
\varphi \text{ isomorphism of pairs}
\end{matrix}
\right.
\end{equation}
Moreover, given two pairs $(\mathcal{L}, \alpha)$, $(\mathcal{L}', \alpha')$
the tensor product
$$
(\mathcal{L}, \alpha) \otimes (\mathcal{L}', \alpha')
=
(\mathcal{L} \otimes \mathcal{L}', \alpha \otimes \alpha')
$$
is another pair. The pair $(\mathcal{O}_S, 1)$ is an identity for this
tensor product operation, and an inverse is given by
$$
(\mathcal{L}, \alpha)^{-1} = (\mathcal{L}^{\otimes -1}, \alpha^{\otimes -1}).
$$
Hence the collection of isomorphism classes of pairs forms an abelian group.
Note that
$$
(\mathcal{L}, \alpha)^{\otimes n}
=
(\mathcal{L}^{\otimes n}, \alpha^{\otimes n})
\xrightarrow{\alpha}
(\mathcal{O}_S, 1)
$$
is an isomorphism
hence every element of this group has order dividing $n$. We warn the reader
that this group is in general {\bf not} the $n$-torsion in $\Pic(S)$.

\begin{lemma}
\label{lemma-describe-h1-mun}
Let $S$ be a scheme. There is a canonical identification
$$
H_\etale^1(S, \mu_n) =
\text{group of pairs }(\mathcal{L}, \alpha)\text{ up to isomorphism as above}
$$
if $n$ is invertible on $S$. In general we have
$$
H_{fppf}^1(S, \mu_n) =
\text{group of pairs }(\mathcal{L}, \alpha)\text{ up to isomorphism as above}.
$$
The same result holds with fppf replaced by syntomic.
\end{lemma}

\begin{proof}
We first prove the second isomorphism.
Let $(\mathcal{L}, \alpha)$ be a pair as above.
Choose an affine open covering $S = \bigcup U_i$ such that
$\mathcal{L}|_{U_i} \cong \mathcal{O}_{U_i}$. Say $s_i \in \mathcal{L}(U_i)$
is a generator. Then $\alpha(s_i^{\otimes n}) = f_i \in \mathcal{O}_S^*(U_i)$.
Writing $U_i = \Spec(A_i)$ we see there exists a global
relative complete intersection $A_i \to B_i = A_i[T]/(T^n - f_i)$
such that $f_i$ maps to an $n$th power in $B_i$. In other words, setting
$V_i = \Spec(B_i)$ we obtain a syntomic covering
$\mathcal{V} = \{V_i \to S\}_{i \in I}$ and trivializations
$\varphi_i : (\mathcal{L}, \alpha)|_{V_i} \to (\mathcal{O}_{V_i}, 1)$.

\medskip\noindent
We will use this result (the existence of the covering $\mathcal{V}$)
to associate to this pair a cohomology class in
$H^1_{syntomic}(S, \mu_{n, S})$. We give two (equivalent) constructions.

\medskip\noindent
First construction: using {\v C}ech cohomology.
Over the double overlaps $V_i \times_S V_j$ we have the isomorphism
$$
(\mathcal{O}_{V_i \times_S V_j}, 1)
\xrightarrow{\text{pr}_0^*\varphi_i^{-1}}
(\mathcal{L}|_{V_i \times_S V_j}, \alpha|_{V_i \times_S V_j})
\xrightarrow{\text{pr}_1^*\varphi_j}
(\mathcal{O}_{V_i \times_S V_j}, 1)
$$
of pairs. By (\ref{equation-isomorphisms-pairs}) this is given by an
element $\zeta_{ij} \in \mu_n(V_i \times_S V_j)$. We omit the verification
that these $\zeta_{ij}$'s give a $1$-cocycle, i.e., give
an element $(\zeta_{i_0i_1}) \in \check C(\mathcal{V}, \mu_n)$
with $d(\zeta_{i_0i_1}) = 0$. Thus its class is an element in
$\check H^1(\mathcal{V}, \mu_n)$ and by
Theorem \ref{theorem-cech-ss}
it maps to a cohomology class in $H^1_{syntomic}(S, \mu_{n, S})$.

\medskip\noindent
Second construction: Using torsors. Consider the presheaf
$$
\mu_n(\mathcal{L}, \alpha) :
U
\longmapsto
\mathit{Isom}_U((\mathcal{O}_U, 1), (\mathcal{L}, \alpha)|_U)
$$
on $(\Sch/S)_{syntomic}$.
We may view this as a subpresheaf of
$\SheafHom_\mathcal{O}(\mathcal{O}, \mathcal{L})$ (internal hom
sheaf, see
Modules on Sites, Section \ref{sites-modules-section-internal-hom}).
Since the conditions defining this subpresheaf are local, we see that it is
a sheaf.
By (\ref{equation-isomorphisms-pairs}) this sheaf has a free action of
the sheaf $\mu_{n, S}$. Hence the only thing we have to check is that
it locally has sections. This is true because of the existence of the
trivializing cover $\mathcal{V}$. Hence $\mu_n(\mathcal{L}, \alpha)$
is a $\mu_{n, S}$-torsor and by
Cohomology on Sites, Lemma \ref{sites-cohomology-lemma-torsors-h1}
we obtain a corresponding element of $H^1_{syntomic}(S, \mu_{n, S})$.

\medskip\noindent
Ok, now we have to still show the following
\begin{enumerate}
\item The two constructions give the same cohomology class.
\item Isomorphic pairs give rise to the same cohomology class.
\item The cohomology class of
$(\mathcal{L}, \alpha) \otimes (\mathcal{L}', \alpha')$
is the sum of the cohomology classes of
$(\mathcal{L}, \alpha)$ and $(\mathcal{L}', \alpha')$.
\item If the cohomology class is trivial, then the pair is trivial.
\item Any element of $H^1_{syntomic}(S, \mu_{n, S})$ is the
cohomology class of a pair.
\end{enumerate}
We omit the proof of (1). Part (2) is clear from the second construction,
since isomorphic torsors give the same cohomology classes.
Part (3) is clear from the first construction, since the resulting
{\v C}ech classes add up. Part (4) is clear from the second construction
since a torsor is trivial if and only if it has a global section, see
Cohomology on Sites, Lemma \ref{sites-cohomology-lemma-trivial-torsor}.

\medskip\noindent
Part (5) can be seen as follows (although a direct proof would be
preferable). Suppose $\xi \in H^1_{syntomic}(S, \mu_{n, S})$.
Then $\xi$ maps to an element
$\overline{\xi} \in H^1_{syntomic}(S, \mathbf{G}_{m, S})$
with $n \overline{\xi} = 0$. By
Theorem \ref{theorem-picard-group}
we see that $\overline{\xi}$ corresponds to an invertible sheaf $\mathcal{L}$
whose $n$th tensor power is isomorphic to $\mathcal{O}_S$.
Hence there exists a pair $(\mathcal{L}, \alpha')$ whose cohomology
class $\xi'$ has the same image $\overline{\xi'}$ in
$H^1_{syntomic}(S, \mathbf{G}_{m, S})$. Thus it suffices to show
that $\xi - \xi'$ is the class of a pair. By construction, and the
long exact cohomology sequence above, we see that
$\xi - \xi' = \partial(f)$ for some $f \in H^0(S, \mathcal{O}_S^*)$.
Consider the pair $(\mathcal{O}_S, f)$. We omit the verification
that the cohomology class of this pair is $\partial(f)$, which
finishes the proof of the first identification (with fppf replaced
with syntomic).

\medskip\noindent
To see the first, note that if $n$ is invertible on $S$, then the
covering $\mathcal{V}$ constructed in the first part of the proof
is actually an \'etale covering (compare with the proof of
Lemma \ref{lemma-kummer-sequence}). The rest of the proof is independent
of the topology, apart from the very last argument which uses that
the Kummer sequence is exact, i.e., uses Lemma \ref{lemma-kummer-sequence}.
\end{proof}






\section{Neighborhoods, stalks and points}
\label{section-stalks}

\noindent
We can associate to any geometric point of $S$ a stalk functor which is
exact. A map of sheaves on $S_\etale$ is an isomorphism if and only
if it
is an isomorphism on all these stalks. A complex of abelian sheaves is
exact if and only if the complex of stalks is exact at all geometric points.
Altogether this means that the small \'etale site of a scheme $S$
has enough points. It also turns out that any point of the small \'etale topos
of $S$ (an abstract notion) is given by a geometric point.
Thus in some sense the small \'etale topos of $S$ can be understood in
terms of geometric points and neighbourhoods.

\begin{definition}
\label{definition-geometric-point}
Let $S$ be a scheme.
\begin{enumerate}
\item A {\it geometric point} of $S$ is a morphism
$\Spec(k) \to S$ where $k$ is algebraically closed.
Such a point is usually denoted $\overline{s}$, i.e., by an overlined
small case letter. We often use $\overline{s}$ to denote the scheme
$\Spec(k)$ as well as the morphism, and we use $\kappa(\overline{s})$
to denote $k$.
\item We say $\overline{s}$ {\it lies over} $s$
to indicate that $s \in S$ is the image of $\overline{s}$.
\item An {\it \'etale neighborhood} of a geometric point $\overline{s}$
of $S$ is a commutative diagram
$$
\xymatrix{
& U \ar[d]^\varphi \\
{\overline{s}} \ar[r]^{\overline{s}} \ar[ur]^{\bar u} & S
}
$$
where $\varphi$ is an \'etale morphism of schemes.
We write $(U, \overline{u}) \to (S, \overline{s})$.
\item A {\it morphism of \'etale neighborhoods}
$(U, \overline{u}) \to (U', \overline{u}')$
is an $S$-morphism $h: U \to U'$
such that $\overline{u}' = h \circ \overline{u}$.
\end{enumerate}
\end{definition}

\begin{remark}
\label{remark-etale-between-etale}
Since $U$ and $U'$ are \'etale over $S$, any $S$-morphism
between them is also \'etale, see
Proposition \ref{proposition-etale-morphisms}.
In particular all morphisms of \'etale neighborhoods are \'etale.
\end{remark}

\begin{remark}
\label{remark-etale-neighbourhoods}
Let $S$ be a scheme and $s \in S$ a point. In
More on Morphisms,
Definition \ref{more-morphisms-definition-etale-neighbourhood}
we defined the notion of an \'etale neighbourhood $(U, u) \to (S, s)$
of $(S, s)$. If $\overline{s}$ is a geometric point of $S$ lying over
$s$, then any \'etale neighbourhood $(U, \overline{u}) \to (S, \overline{s})$
gives rise to an \'etale neighbourhood $(U, u)$ of $(S, s)$ by taking
$u \in U$ to be the unique point of $U$ such that $\overline{u}$
lies over $u$. Conversely, given an \'etale neighbourhood $(U, u)$
of $(S, s)$ the residue field extension $\kappa(u)/\kappa(s)$
is finite separable (see
Proposition \ref{proposition-etale-morphisms})
and hence we can find an embedding $\kappa(u) \subset \kappa(\overline{s})$
over $\kappa(s)$. In other words, we can find a geometric point
$\overline{u}$ of $U$ lying over $u$ such that $(U, \overline{u})$
is an \'etale neighbourhood of $(S, \overline{s})$.
We will use these observations to go between the two types of
\'etale neighbourhoods.
\end{remark}

\begin{lemma}
\label{lemma-cofinal-etale}
Let $S$ be a scheme, and let $\overline{s}$ be a geometric point of $S$.
The category of \'etale neighborhoods is cofiltered. More precisely:
\begin{enumerate}
\item Let $(U_i, \overline{u}_i)_{i = 1, 2}$ be two \'etale neighborhoods of
$\overline{s}$ in $S$. Then there exists a third \'etale neighborhood
$(U, \overline{u})$ and morphisms
$(U, \overline{u}) \to (U_i, \overline{u}_i)$, $i = 1, 2$.
\item Let $h_1, h_2: (U, \overline{u}) \to (U', \overline{u}')$ be two
morphisms between \'etale neighborhoods of $\overline{s}$. Then there exist an
\'etale neighborhood $(U'', \overline{u}'')$ and a morphism
$h : (U'', \overline{u}'') \to (U, \overline{u})$
which equalizes $h_1$ and $h_2$, i.e., such that
$h_1 \circ h = h_2 \circ h$.
\end{enumerate}
\end{lemma}

\begin{proof}
For part (1), consider the fibre product $U = U_1 \times_S U_2$.
It is \'etale over both $U_1$ and $U_2$ because \'etale morphisms are
preserved under base change, see
Proposition \ref{proposition-etale-morphisms}.
The map $\overline{s} \to U$ defined by $(\overline{u}_1, \overline{u}_2)$
gives it the structure of an \'etale neighborhood mapping to both
$U_1$ and $U_2$. For part (2), define $U''$ as the fibre product
$$
\xymatrix{
U'' \ar[r] \ar[d] & U \ar[d]^{(h_1, h_2)} \\
U' \ar[r]^-\Delta & U' \times_S U'.
}
$$
Since $\overline{u}$ and $\overline{u}'$ agree over $S$ with $\overline{s}$,
we see that $\overline{u}'' = (\overline{u}, \overline{u}')$ is a geometric
point of $U''$. In particular $U'' \not = \emptyset$.
Moreover, since $U'$ is \'etale over $S$, so is the fibre product
$U'\times_S U'$ (see
Proposition \ref{proposition-etale-morphisms}).
Hence the vertical arrow $(h_1, h_2)$ is \'etale by
Remark \ref{remark-etale-between-etale} above.
Therefore $U''$ is \'etale over $U'$ by base change, and hence also
\'etale over $S$ (because compositions of \'etale morphisms are \'etale).
Thus $(U'', \overline{u}'')$ is a solution to the problem.
\end{proof}

\begin{lemma}
\label{lemma-geometric-lift-to-cover}
Let $S$ be a scheme.
Let $\overline{s}$ be a geometric point of $S$.
Let $(U, \overline{u})$ be an \'etale neighborhood of $\overline{s}$.
Let $\mathcal{U} = \{\varphi_i : U_i \to U \}_{i\in I}$ be an \'etale covering.
Then there exist $i \in I$ and $\overline{u}_i : \overline{s} \to U_i$
such that $\varphi_i : (U_i, \overline{u}_i) \to (U, \overline{u})$
is a morphism of \'etale neighborhoods.
\end{lemma}

\begin{proof}
As $U = \bigcup_{i\in I} \varphi_i(U_i)$, the fibre product
$\overline{s} \times_{\overline{u}, U, \varphi_i} U_i$ is not empty
for some $i$. Then look at the cartesian diagram
$$
\xymatrix{
\overline{s} \times_{\overline{u}, U, \varphi_i} U_i
\ar[d]^{\text{pr}_1} \ar[r]_-{\text{pr}_2} & U_i
\ar[d]^{\varphi_i} \\
\Spec(k) = \overline{s} \ar@/^1pc/[u]^\sigma
\ar[r]^-{\overline{u}} & U
}
$$
The projection $\text{pr}_1$ is the base change of an \'etale morphisms so it
is \'etale, see
Proposition \ref{proposition-etale-morphisms}.
Therefore, $\overline{s} \times_{\overline{u}, U, \varphi_i} U_i$
is a disjoint union of finite separable extensions of $k$, by
Proposition \ref{proposition-etale-morphisms}. Here
$\overline{s} = \Spec(k)$. But $k$ is algebraically closed, so all
these extensions are trivial, and there exists a section $\sigma$ of
$\text{pr}_1$. The composition
$\text{pr}_2 \circ \sigma$ gives a map compatible with $\overline{u}$.
\end{proof}

\begin{definition}
\label{definition-stalk}
Let $S$ be a scheme.
Let $\mathcal{F}$ be a presheaf on $S_\etale$.
Let $\overline{s}$ be a geometric point of $S$.
The {\it stalk} of $\mathcal{F}$ at $\overline{s}$ is
$$
\mathcal{F}_{\overline{s}}
=
\colim_{(U, \overline{u})} \mathcal{F}(U)
$$
where $(U, \overline{u})$ runs over all \'etale
neighborhoods of $\overline{s}$ in $S$.
\end{definition}

\noindent
By Lemma \ref{lemma-cofinal-etale}, this colimit is over a filtered
index category, namely the opposite of the category of \'etale neighbourhoods.
In other words, an element of $\mathcal{F}_{\overline{s}}$ can be
thought of as a triple $(U, \overline{u}, \sigma)$ where
$\sigma \in \mathcal{F}(U)$. Two triples
$(U, \overline{u}, \sigma)$, $(U', \overline{u}', \sigma')$
define the same element of the stalk if there exists a third
\'etale neighbourhood $(U'', \overline{u}'')$ and morphisms of \'etale
neighbourhoods $h : (U'', \overline{u}'') \to (U, \overline{u})$,
$h' : (U'', \overline{u}'') \to (U', \overline{u}')$ such that
$h^*\sigma = (h')^*\sigma'$ in $\mathcal{F}(U'')$. See
Categories, Section \ref{categories-section-directed-colimits}.

\begin{lemma}
\label{lemma-stalk-gives-point}
Let $S$ be a scheme. Let $\overline{s}$ be a geometric point of $S$.
Consider the functor
\begin{align*}
u : S_\etale & \longrightarrow \textit{Sets}, \\
U & \longmapsto
|U_{\overline{s}}|
=
\{\overline{u} \text{ such that }(U, \overline{u})
\text{ is an \'etale neighbourhood of }\overline{s}\}.
\end{align*}
Here $|U_{\overline{s}}|$ denotes the underlying set of the geometric fibre.
Then $u$ defines a point $p$ of the site $S_\etale$
(Sites, Definition \ref{sites-definition-point})
and its associated stalk functor $\mathcal{F} \mapsto \mathcal{F}_p$
(Sites, Equation \ref{sites-equation-stalk})
is the functor $\mathcal{F} \mapsto \mathcal{F}_{\overline{s}}$
defined above.
\end{lemma}

\begin{proof}
In the proof of
Lemma \ref{lemma-geometric-lift-to-cover}
we have seen that the scheme $U_{\overline{s}}$ is a disjoint union of
schemes isomorphic to $\overline{s}$. Thus we can also think of
$|U_{\overline{s}}|$ as the set of geometric points of $U$ lying over
$\overline{s}$, i.e., as the collection of morphisms
$\overline{u} : \overline{s} \to U$ fitting into the diagram of
Definition \ref{definition-geometric-point}.
From this it follows that $u(S)$ is a singleton, and that
$u(U \times_V W) = u(U) \times_{u(V)} u(W)$
whenever $U \to V$ and $W \to V$ are morphisms in $S_\etale$.
And, given a covering $\{U_i \to U\}_{i \in I}$ in $S_\etale$
we see that $\coprod u(U_i) \to u(U)$ is surjective by
Lemma \ref{lemma-geometric-lift-to-cover}.
Hence
Sites, Proposition \ref{sites-proposition-point-limits}
applies, so $p$ is a point of the site $S_\etale$.
Finally, our functor $\mathcal{F} \mapsto \mathcal{F}_{\overline{s}}$
is given by exactly the same colimit as the functor
$\mathcal{F} \mapsto \mathcal{F}_p$ associated to $p$ in
Sites, Equation \ref{sites-equation-stalk}
which proves the final assertion.
\end{proof}

\begin{remark}
\label{remark-map-stalks}
Let $S$ be a scheme and let $\overline{s} : \Spec(k) \to S$
and $\overline{s}' : \Spec(k') \to S$ be two geometric points of
$S$. A {\it morphism $a : \overline{s} \to \overline{s}'$ of geometric points}
is simply a morphism $a : \Spec(k) \to \Spec(k')$ such that
$\overline{s}' \circ a = \overline{s}$. Given such a morphism we obtain
a functor from the category of \'etale neighbourhoods of $\overline{s}'$
to the category of \'etale neighbourhoods of $\overline{s}$ by the rule
$(U, \overline{u}') \mapsto (U, \overline{u}' \circ a)$. Hence we obtain
a canonical map
$$
\mathcal{F}_{\overline{s}'}
=
\colim_{(U, \overline{u}')} \mathcal{F}(U)
\longrightarrow
\colim_{(U, \overline{u})} \mathcal{F}(U)
=
\mathcal{F}_{\overline{s}}
$$
from Categories, Lemma \ref{categories-lemma-functorial-colimit}. Using the
description of elements of stalks as triples this maps the element of
$\mathcal{F}_{\overline{s}'}$ represented by the triple
$(U, \overline{u}', \sigma)$ to the element of $\mathcal{F}_{\overline{s}}$
represented by the triple $(U, \overline{u}' \circ a, \sigma)$.
Since the functor above is clearly an equivalence we conclude that this
canonical map is an isomorphism of stalk functors.

\medskip\noindent
Let us make sure we have the map of stalks corresponding to $a$ pointing
in the correct direction. Note that the above means, according to
Sites, Definition \ref{sites-definition-morphism-points},
that $a$ defines a morphism $a : p \to p'$ between the points $p, p'$ of
the site $S_\etale$ associated to $\overline{s}, \overline{s}'$ by
Lemma \ref{lemma-stalk-gives-point}. There are more general morphisms of
points (corresponding to specializations of points of $S$) which we will
describe later, and which will not be isomorphisms, see Section
\ref{section-specialization}.
\end{remark}

\begin{lemma}
\label{lemma-stalk-exact}
Let $S$ be a scheme. Let $\overline{s}$ be a geometric point of $S$.
\begin{enumerate}
\item The stalk functor
$\textit{PAb}(S_\etale) \to \textit{Ab}$,
$\mathcal{F}  \mapsto  \mathcal{F}_{\overline{s}}$
is exact.
\item We have $(\mathcal{F}^\#)_{\overline{s}} = \mathcal{F}_{\overline{s}}$
for any presheaf of sets $\mathcal{F}$ on $S_\etale$.
\item The functor
$\textit{Ab}(S_\etale) \to \textit{Ab}$,
$\mathcal{F} \mapsto \mathcal{F}_{\overline{s}}$ is exact.
\item Similarly the functors
$\textit{PSh}(S_\etale) \to \textit{Sets}$ and
$\Sh(S_\etale) \to \textit{Sets}$ given by the stalk functor
$\mathcal{F} \mapsto \mathcal{F}_{\overline{s}}$ are exact (see
Categories, Definition \ref{categories-definition-exact})
and commute with arbitrary colimits.
\end{enumerate}
\end{lemma}

\begin{proof}
Before we indicate how to prove this by direct arguments
we note that the result follows from the general material in
Modules on Sites, Section \ref{sites-modules-section-stalks}.
This is true because $\mathcal{F} \mapsto \mathcal{F}_{\overline{s}}$
comes from a point of the small \'etale site of $S$, see
Lemma \ref{lemma-stalk-gives-point}.
We will only give a direct proof of (1), (2) and (3), and omit
a direct proof of (4).

\medskip\noindent
Exactness as a functor on $\textit{PAb}(S_\etale)$ is formal from the
fact that directed colimits commute with all colimits and with finite
limits. The identification of the stalks in (2) is via the map
$$
\kappa :
\mathcal{F}_{\overline{s}}
\longrightarrow
(\mathcal{F}^\#)_{\overline{s}}
$$
induced by the natural morphism $\mathcal{F}\to \mathcal{F}^\#$, see
Theorem \ref{theorem-sheafification}.
We claim that this map is an isomorphism of abelian groups. We will show
injectivity and omit the proof of surjectivity.

\medskip\noindent
Let $\sigma\in \mathcal{F}_{\overline{s}}$.
There exists an \'etale neighborhood
$(U, \overline{u})\to (S, \overline{s})$ such that $\sigma$ is the image of some
section $s \in \mathcal{F}(U)$. If $\kappa(\sigma) = 0$ in
$(\mathcal{F}^\#)_{\overline{s}}$ then there exists a morphism of \'etale
neighborhoods $(U', \overline{u}')\to (U, \overline{u})$ such that
$s|_{U'}$ is zero in $\mathcal{F}^\#(U')$. It follows there
exists an \'etale covering
$\{U_i'\to U'\}_{i\in I}$ such that $s|_{U_i'}=0$ in
$\mathcal{F}(U_i')$ for all $i$. By Lemma \ref{lemma-geometric-lift-to-cover}
there exist $i \in I$ and a morphism
$\overline{u}_i': \overline{s} \to U_i'$ such that
$(U_i', \overline{u}_i') \to (U', \overline{u}')\to (U, \overline{u})$
are morphisms of \'etale neighborhoods. Hence $\sigma = 0$
since $(U_i', \overline{u}_i') \to (U, \overline{u})$
is a morphism of \'etale neighbourhoods such that
we have $s|_{U'_i}=0$. This proves $\kappa$ is injective.

\medskip\noindent
To show that the functor $\textit{Ab}(S_\etale) \to \textit{Ab}$ is
exact, consider any short exact sequence in $\textit{Ab}(S_\etale)$:
$
0\to \mathcal{F}\to \mathcal{G}\to \mathcal H \to 0.
$
This gives us the exact sequence of presheaves
$$
0 \to \mathcal{F} \to \mathcal{G} \to \mathcal H \to
\mathcal H/^p\mathcal{G} \to 0,
$$
where $/^p$ denotes the quotient in $\textit{PAb}(S_\etale)$.
Taking stalks at
$\overline{s}$, we see that $(\mathcal H /^p\mathcal{G})_{\bar{s}} =
(\mathcal H /\mathcal{G})_{\bar{s}} = 0$, since the sheafification of
$\mathcal H/^p\mathcal{G}$ is $0$.
Therefore,
$$
0\to \mathcal{F}_{\overline{s}} \to \mathcal{G}_{\overline{s}} \to
\mathcal{H}_{\overline{s}} \to 0 = (\mathcal H/^p\mathcal{G})_{\overline{s}}
$$
is exact, since taking stalks is exact as a functor from presheaves.
\end{proof}

\begin{theorem}
\label{theorem-exactness-stalks}
Let $S$ be a scheme.
A map $a : \mathcal{F} \to \mathcal{G}$ of sheaves of sets is injective
(resp.\ surjective) if and only if the map on stalks
$a_{\overline{s}} : \mathcal{F}_{\overline{s}} \to \mathcal{G}_{\overline{s}}$
is injective (resp.\ surjective) for all geometric points of $S$.
A sequence of abelian sheaves on $S_\etale$ is exact
if and only if it is exact on all stalks at geometric points of $S$.
\end{theorem}

\begin{proof}
The necessity of exactness on stalks follows from
Lemma \ref{lemma-stalk-exact}.
For the converse, it suffices to show that a map of sheaves is surjective
(respectively injective) if and only if it is surjective (respectively
injective) on all stalks. We prove this in the case of surjectivity, and omit
the proof in the case of injectivity.

\medskip\noindent
Let $\alpha : \mathcal{F} \to \mathcal{G}$ be a map of sheaves such
that $\mathcal{F}_{\overline{s}} \to \mathcal{G}_{\overline{s}}$
is surjective for all geometric points. Fix
$U \in \Ob(S_\etale)$
and $s \in \mathcal{G}(U)$. For every $u \in U$ choose some
$\overline{u} \to U$ lying over $u$ and an \'etale neighborhood
$(V_u , \overline{v}_u) \to (U, \overline{u})$ such that
$s|_{V_u} = \alpha(s_{V_u})$ for some
$s_{V_u} \in \mathcal{F}(V_u)$.
This is possible since $\alpha$ is surjective on
stalks. Then $\{V_u \to U\}_{u \in U}$
is an \'etale covering on which the restrictions of $s$
are in the image of the map $\alpha$.
Thus, $\alpha$ is surjective, see
Sites, Section \ref{sites-section-sheaves-injective}.
\end{proof}

\begin{remarks}
\label{remarks-enough-points}
On points of the geometric sites.
\begin{enumerate}
\item Theorem \ref{theorem-exactness-stalks} says that the family of points
of $S_\etale$ given by the geometric points of $S$
(Lemma \ref{lemma-stalk-gives-point}) is conservative, see
Sites, Definition \ref{sites-definition-enough-points}.
In particular $S_\etale$ has enough points.
\item Suppose $\mathcal{F}$ is a sheaf on the big \'etale site
\label{item-stalks-big}
of $S$. Let $T \to S$ be an object of the big \'etale site of $S$,
and let $\overline{t}$ be a geometric point of $T$. Then we define
$\mathcal{F}_{\overline{t}}$ as the stalk
of the restriction $\mathcal{F}|_{T_\etale}$ of $\mathcal{F}$
to the small \'etale site of $T$. In other words, we can define
the stalk of $\mathcal{F}$ at any geometric point of any
scheme $T/S \in \Ob((\Sch/S)_\etale)$.
\item The big \'etale site of $S$ also has enough points, by
considering all geometric points of all objects of this site, see
(\ref{item-stalks-big}).
\end{enumerate}
\end{remarks}

\noindent
The following lemma should be skipped on a first reading.

\begin{lemma}
\label{lemma-points-small-etale-site}
Let $S$ be a scheme.
\begin{enumerate}
\item Let $p$ be a point of the small \'etale site
$S_\etale$ of $S$ given by a functor
$u : S_\etale \to \textit{Sets}$.
Then there exists a geometric point $\overline{s}$ of $S$ such that
$p$ is isomorphic to the point of $S_\etale$ associated to
$\overline{s}$ in
Lemma \ref{lemma-stalk-gives-point}.
\item Let $p : \Sh(pt) \to \Sh(S_\etale)$ be a point
of the small \'etale topos of $S$. Then $p$ comes from a geometric point
of $S$, i.e., the stalk functor $\mathcal{F} \mapsto \mathcal{F}_p$
is isomorphic to a stalk functor as defined in
Definition \ref{definition-stalk}.
\end{enumerate}
\end{lemma}

\begin{proof}
By
Sites, Lemma \ref{sites-lemma-point-site-topos}
there is a one to one correspondence between points of the site and points
of the associated topos, hence it suffices to prove (1).
By
Sites, Proposition \ref{sites-proposition-point-limits}
the functor $u$ has the following properties:
(a) $u(S) = \{*\}$, (b) $u(U \times_V W) = u(U) \times_{u(V)} u(W)$, and
(c) if $\{U_i \to U\}$ is an \'etale covering, then
$\coprod u(U_i) \to u(U)$ is surjective.
In particular, if $U' \subset U$ is an open subscheme,
then $u(U') \subset u(U)$. Moreover, by
Sites, Lemma \ref{sites-lemma-point-site-topos}
we can write $u(U) = p^{-1}(h_U^\#)$, in other words $u(U)$ is the
stalk of the representable sheaf $h_U$. If
$U = V \amalg W$, then we see that $h_U = (h_V \amalg h_W)^\#$ and we get
$u(U) = u(V) \amalg u(W)$ since $p^{-1}$ is exact.

\medskip\noindent
Consider the restriction of $u$ to $S_{Zar}$. By
Sites, Examples \ref{sites-example-point-topological} and
\ref{sites-example-point-topology}
there exists a unique point $s \in S$ such that for $S' \subset S$ open we
have $u(S') = \{*\}$ if $s \in S'$ and $u(S') = \emptyset$ if $s \not \in S'$.
Note that if $\varphi : U \to S$ is an object of $S_\etale$ then
$\varphi(U) \subset S$ is open (see
Proposition \ref{proposition-etale-morphisms})
and $\{U \to \varphi(U)\}$ is an \'etale covering. Hence we conclude that
$u(U) = \emptyset \Leftrightarrow s \not \in \varphi(U)$.

\medskip\noindent
Pick a geometric point $\overline{s} : \overline{s} \to S$ lying over $s$, see
Definition \ref{definition-geometric-point}
for customary abuse of notation. Suppose that $\varphi : U \to S$ is an object
of $S_\etale$ with $U$ affine. Note that $\varphi$ is separated, and
that the fibre $U_s$ of $\varphi$ over $s$ is an affine scheme over
$\Spec(\kappa(s))$ which is the spectrum of a finite product of
finite separable extensions $k_i$ of $\kappa(s)$. Hence we may apply
\'Etale Morphisms, Lemma \ref{etale-lemma-etale-etale-local-technical}
to get an \'etale neighbourhood $(V, \overline{v})$ of $(S, \overline{s})$
such that
$$
U \times_S V = U_1 \amalg \ldots \amalg U_n \amalg W
$$
with $U_i \to V$ an isomorphism and $W$ having no point lying over
$\overline{v}$. Thus we conclude that
$$
u(U) \times u(V) =
u(U \times_S V) =
u(U_1) \amalg \ldots \amalg u(U_n) \amalg u(W)
$$
and of course also $u(U_i) = u(V)$. After shrinking $V$ a bit we can
assume that $V$ has exactly one point lying over $s$, and hence $W$ has no
point lying over $s$. By the above this then gives $u(W) = \emptyset$.
Hence we obtain
$$
u(U) \times u(V) =
u(U_1) \amalg \ldots \amalg u(U_n) =
\coprod\nolimits_{i = 1, \ldots, n} u(V)
$$
where the equality signs are compatible with the maps to the set $u(V)$.
Note that $u(V) \not = \emptyset$ as $s$ is in the image of $V \to S$.
In particular, we see that in this situation $u(U)$ is a finite
set with $n$ element, namely the left hand side has fibre $u(U)$ over
any element of $u(V)$ and the right hand side has fibre of cardinality $n$.

\medskip\noindent
Consider the limit
$$
\lim_{(V, \overline{v})} u(V)
$$
over the category of \'etale neighbourhoods $(V, \overline{v})$ of
$\overline{s}$. It is clear that we get the same value when taking
the limit over the subcategory of $(V, \overline{v})$ with $V$ affine.
By the previous paragraph (applied with the roles of $V$ and $U$ switched)
we see that in this case $u(V)$ is always a finite nonempty set.
Moreover, the limit is cofiltered, see
Lemma \ref{lemma-cofinal-etale}.
Hence by
Categories, Section \ref{categories-section-codirected-limits}
the limit is nonempty. Pick an element $x$ from this limit.
This means we obtain a $x_{V, \overline{v}} \in u(V)$ for
every \'etale neighbourhood $(V, \overline{v})$ of $(S, \overline{s})$
such that for every morphism of \'etale neighbourhoods
$\varphi : (V', \overline{v}') \to (V, \overline{v})$ we have
$u(\varphi)(x_{V', \overline{v}'}) = x_{V, \overline{v}}$.

\medskip\noindent
We will use the choice of $x$ to construct a functorial bijective map
$$
c : |U_{\overline{s}}| \longrightarrow u(U)
$$
for $U \in \Ob(S_\etale)$ which will conclude the proof. See
Lemma \ref{lemma-stalk-gives-point}
and its proof for a description of $|U_{\overline{s}}|$.
First we claim that it suffices to construct the map for $U$ affine.
We omit the proof of this claim.
Assume $U \to S$ in $S_\etale$ with $U$ affine, and let
$\overline{u} : \overline{s} \to U$ be an element of $|U_{\overline{s}}|$.
Choose a $(V, \overline{v})$ such that $U \times_S V$ decomposes
as in the third paragraph of the proof.
Then the pair $(\overline{u}, \overline{v})$ gives a geometric point of
$U \times_S V$ lying over $\overline{v}$ and determines one of the
components $U_i$ of $U \times_S V$. More precisely, there exists
a section $\sigma : V \to U \times_S V$ of the projection $\text{pr}_U$
such that $(\overline{u}, \overline{v}) = \sigma \circ \overline{v}$. Set
$c(\overline{u}) = u(\text{pr}_U)(u(\sigma)(x_{V, \overline{v}})) \in u(U)$.
We have to check this is independent of the choice of $(V, \overline{v})$. By
Lemma \ref{lemma-cofinal-etale}
the category of \'etale neighbourhoods is cofiltered.
Hence it suffice to show
that given a morphism of \'etale neighbourhood
$\varphi : (V', \overline{v}') \to (V, \overline{v})$ and a choice of a
section $\sigma' : V' \to U \times_S V'$ of the projection such that
$(\overline{u}, \overline{v'}) = \sigma' \circ \overline{v}'$
we have $u(\sigma')(x_{V', \overline{v}'}) = u(\sigma)(x_{V, \overline{v}})$.
Consider the diagram
$$
\xymatrix{
V' \ar[d]^{\sigma'} \ar[r]_\varphi & V \ar[d]^\sigma \\
U \times_S V' \ar[r]^{1 \times \varphi} &
U \times_S V
}
$$
Now, it may not be the case that this diagram commutes. The reason is
that the schemes $V'$ and $V$ may not be connected, and hence
the decompositions used to construct $\sigma'$ and $\sigma$ above may
not be unique. But we do know that
$\sigma \circ \varphi \circ \overline{v}' =
(1 \times \varphi) \circ \sigma' \circ \overline{v}'$
by construction. Hence, since $U \times_S V$ is \'etale over $S$,
there exists an open neighbourhood
$V'' \subset V'$ of $\overline{v'}$ such that the diagram does
commute when restricted to $V''$, see
Morphisms, Lemma \ref{morphisms-lemma-value-at-one-point}.
This means we may extend the diagram above to
$$
\xymatrix{
V'' \ar[r] \ar[d]^{\sigma'|_{V''}} &
V' \ar[d]^{\sigma'} \ar[r]_\varphi &
V \ar[d]^\sigma \\
U \times_S V'' \ar[r] &
U \times_S V' \ar[r]^{1 \times \varphi} &
U \times_S V
}
$$
such that the left square and the outer rectangle commute.
Since $u$ is a functor this implies that
$x_{V'', \overline{v}'}$ maps to the same element in
$u(U \times_S V)$ no matter which route we take through the
diagram. On the other hand, it maps to the elements
$x_{V', \overline{v}'}$ and $x_{V, \overline{v}}$ in
$u(V')$ and $u(V)$. This implies the desired equality
$u(\sigma')(x_{V', \overline{v}'}) = u(\sigma)(x_{V, \overline{v}})$.

\medskip\noindent
In a similar manner one proves that the construction
$c : |U_{\overline{s}}| \to u(U)$ is functorial in $U$;
details omitted. And finally, by the results of the
third paragraph it is clear that the map $c$ is bijective
which ends the proof of the lemma.
\end{proof}







\section{Points in other topologies}
\label{section-points-topologies}

\noindent
In this section we briefly discuss the existence of points for some
sites other than the \'etale site of a scheme. We refer to
Sites, Section \ref{sites-section-sites-enough-points}
and
Topologies, Section \ref{topologies-section-procedure} ff
for the terminology used in this section.
All of the geometric sites have enough points.

\begin{lemma}
\label{lemma-points-fppf}
Let $S$ be a scheme. All of the following sites have enough points
$S_{affine, Zar}$, $S_{Zar}$, $S_{affine, \etale}$, $S_\etale$,
$(\Sch/S)_{Zar}$, $(\textit{Aff}/S)_{Zar}$,
$(\Sch/S)_\etale$, $(\textit{Aff}/S)_\etale$,
$(\Sch/S)_{smooth}$, $(\textit{Aff}/S)_{smooth}$,
$(\Sch/S)_{syntomic}$, $(\textit{Aff}/S)_{syntomic}$,
$(\Sch/S)_{fppf}$, and $(\textit{Aff}/S)_{fppf}$.
\end{lemma}

\begin{proof}
For each of the big sites the associated topos is equivalent to the
topos defined by the site $(\textit{Aff}/S)_\tau$, see
Topologies, Lemmas \ref{topologies-lemma-affine-big-site-Zariski},
\ref{topologies-lemma-affine-big-site-etale},
\ref{topologies-lemma-affine-big-site-smooth},
\ref{topologies-lemma-affine-big-site-syntomic}, and
\ref{topologies-lemma-affine-big-site-fppf}.
The result for the sites $(\textit{Aff}/S)_\tau$ follows immediately
from Deligne's result
Sites, Lemma \ref{sites-lemma-criterion-points}.

\medskip\noindent
The result for $S_{Zar}$ is clear. The result for $S_{affine, Zar}$
follows from Deligne's result. The result for $S_\etale$
either follows from (the proof of)
Theorem \ref{theorem-exactness-stalks}
or from Topologies, Lemma \ref{topologies-lemma-alternative}
and Deligne's result applied to $S_{affine, \etale}$.
\end{proof}

\noindent
The lemma above guarantees the existence of points, but it doesn't
tell us what these points look like. We can explicitly construct
{\it some} points as follows.
Suppose $\overline{s} : \Spec(k) \to S$ is a geometric
point with $k$ algebraically closed. Consider the functor
$$
u : (\Sch/S)_{fppf} \longrightarrow \textit{Sets},
\quad
u(U) = U(k) = \Mor_S(\Spec(k), U).
$$
Note that $U \mapsto U(k)$ commutes with finite limits as
$S(k) = \{\overline{s}\}$ and
$(U_1 \times_U U_2)(k) = U_1(k) \times_{U(k)} U_2(k)$.
Moreover, if $\{U_i \to U\}$ is an fppf covering, then
$\coprod U_i(k) \to U(k)$ is surjective.
By
Sites, Proposition \ref{sites-proposition-point-limits}
we see that $u$ defines a point $p$ of $(\Sch/S)_{fppf}$ with
stalks
$$
\mathcal{F}_p = \colim_{(U, x)} \mathcal{F}(U)
$$
where the colimit is over pairs $U \to S$, $x \in U(k)$ as usual.
But... this category has an initial object, namely
$(\Spec(k), \text{id})$, hence we see that
$$
\mathcal{F}_p = \mathcal{F}(\Spec(k))
$$
which isn't terribly interesting! In fact, in general these points won't
form a conservative family of points. A more interesting type of point
is described in the following remark.

\begin{remark}
\label{remark-points-fppf-site}
\begin{reference}
This is discussed in \cite{Schroeer}.
\end{reference}
Let $S = \Spec(A)$ be an affine scheme. Let $(p, u)$ be a point of
the site $(\textit{Aff}/S)_{fppf}$, see
Sites, Sections \ref{sites-section-points} and
\ref{sites-section-construct-points}. Let $B = \mathcal{O}_p$ be the stalk
of the structure sheaf at the point $p$. Recall that
$$
B = \colim_{(U, x)} \mathcal{O}(U) =
\colim_{(\Spec(C), x_C)} C
$$
where $x_C \in u(\Spec(C))$. It can happen that
$\Spec(B)$ is an object of $(\textit{Aff}/S)_{fppf}$
and that there is an element $x_B \in u(\Spec(B))$ mapping to
the compatible system $x_C$. In this case the system of neighbourhoods
has an initial object and it follows that
$\mathcal{F}_p = \mathcal{F}(\Spec(B))$ for any sheaf $\mathcal{F}$
on $(\textit{Aff}/S)_{fppf}$. It is straightforward
to see that if $\mathcal{F} \mapsto \mathcal{F}(\Spec(B))$ defines a point
of $\Sh((\textit{Aff}/S)_{fppf})$, then
$B$ has to be a local $A$-algebra such that for every faithfully flat,
finitely presented ring map $B \to B'$ there is a section $B' \to B$.
Conversely, for any such $A$-algebra $B$ the functor
$\mathcal{F} \mapsto \mathcal{F}(\Spec(B))$ is the stalk functor
of a point. Details omitted. It is not clear what a general point of the
site $(\textit{Aff}/S)_{fppf}$ looks like.
\end{remark}











\section{Supports of abelian sheaves}
\label{section-support}

\noindent
First we talk about supports of local sections.

\begin{lemma}
\label{lemma-support-subsheaf-final}
Let $S$ be a scheme. Let $\mathcal{F}$ be a subsheaf of the final
object of the \'etale topos of $S$ (see
Sites, Example \ref{sites-example-singleton-sheaf}).
Then there exists a unique open
$W \subset S$ such that $\mathcal{F} = h_W$.
\end{lemma}

\begin{proof}
The condition means that $\mathcal{F}(U)$ is a singleton or
empty for all $\varphi : U \to S$ in $\Ob(S_\etale)$.
In particular local sections always glue. If
$\mathcal{F}(U) \not = \emptyset$, then
$\mathcal{F}(\varphi(U)) \not = \emptyset$ because
$\{\varphi : U \to \varphi(U)\}$ is a covering.
Hence we can take
$W = \bigcup_{\varphi : U \to S, \mathcal{F}(U) \not = \emptyset} \varphi(U)$.
\end{proof}

\begin{lemma}
\label{lemma-zero-over-image}
Let $S$ be a scheme.
Let $\mathcal{F}$ be an abelian sheaf on $S_\etale$.
Let $\sigma \in \mathcal{F}(U)$ be a local section.
There exists an open subset $W \subset U$ such that
\begin{enumerate}
\item $W \subset U$ is the largest Zariski open subset of $U$ such
that $\sigma|_W = 0$,
\item for every $\varphi : V \to U$ in $S_\etale$ we have
$$
\sigma|_V = 0 \Leftrightarrow \varphi(V) \subset W,
$$
\item for every geometric point $\overline{u}$ of $U$ we have
$$
(U, \overline{u}, \sigma) = 0\text{ in }\mathcal{F}_{\overline{s}}
\Leftrightarrow
\overline{u} \in W
$$
where $\overline{s} = (U \to S) \circ \overline{u}$.
\end{enumerate}
\end{lemma}

\begin{proof}
Since $\mathcal{F}$ is a sheaf in the \'etale topology the restriction of
$\mathcal{F}$ to $U_{Zar}$ is a sheaf on $U$ in the Zariski topology.
Hence there exists a Zariski open $W$ having property (1), see
Modules, Lemma \ref{modules-lemma-support-section-closed}. Let
$\varphi : V \to U$ be an arrow of $S_\etale$. Note that
$\varphi(V) \subset U$ is an open subset and that
$\{V \to \varphi(V)\}$ is an \'etale covering. Hence if
$\sigma|_V = 0$, then by the sheaf condition for $\mathcal{F}$ we
see that $\sigma|_{\varphi(V)} = 0$. This proves (2).
To prove (3) we have to show that if $(U, \overline{u}, \sigma)$
defines the zero element of $\mathcal{F}_{\overline{s}}$, then
$\overline{u} \in W$. This is true because the assumption means
there exists a morphism of \'etale neighbourhoods
$(V, \overline{v}) \to (U, \overline{u})$ such that
$\sigma|_V = 0$. Hence by (2) we see that $V \to U$ maps into $W$, and
hence $\overline{u} \in W$.
\end{proof}

\noindent
Let $S$ be a scheme. Let $s \in S$.
Let $\mathcal{F}$ be a sheaf on $S_\etale$. By
Remark \ref{remark-map-stalks}
the isomorphism class of the stalk of the sheaf $\mathcal{F}$
at a geometric points lying over $s$ is well defined.

\begin{definition}
\label{definition-support}
Let $S$ be a scheme.
Let $\mathcal{F}$ be an abelian sheaf on $S_\etale$.
\begin{enumerate}
\item The {\it support of $\mathcal{F}$} is the set of
points $s \in S$ such that $\mathcal{F}_{\overline{s}} \not = 0$
for any (some) geometric point $\overline{s}$ lying over $s$.
\item Let $\sigma \in \mathcal{F}(U)$ be a section.
The {\it support of $\sigma$} is the closed subset $U \setminus W$, where
$W \subset U$ is the largest open subset of $U$ on which $\sigma$
restricts to zero (see
Lemma \ref{lemma-zero-over-image}).
\end{enumerate}
\end{definition}

\noindent
In general the support of an abelian sheaf is not closed.
For example, suppose that $S = \mathbf{A}^1_{\mathbf{C}}$.
Let $i_t : \Spec(\mathbf{C}) \to S$ be the inclusion of the
point $t \in \mathbf{C}$.
We will see later that $\mathcal{F}_t = i_{t, *}(\mathbf{Z}/2\mathbf{Z})$
is an abelian sheaf whose support is exactly $\{t\}$, see
Section \ref{section-closed-immersions}.
Then
$$
\bigoplus\nolimits_{n \in \mathbf{N}} \mathcal{F}_n
$$
is an abelian sheaf with support $\{1, 2, 3, \ldots\} \subset S$.
This is true because taking stalks commutes with colimits, see
Lemma \ref{lemma-stalk-exact}.
Thus an example of an abelian sheaf whose support is not closed.
Here are some basic facts on supports of sheaves and sections.

\begin{lemma}
\label{lemma-support-section-closed}
Let $S$ be a scheme.
Let $\mathcal{F}$ be an abelian sheaf on $S_\etale$.
Let $U \in \Ob(S_\etale)$ and $\sigma \in \mathcal{F}(U)$.
\begin{enumerate}
\item The support of $\sigma$ is closed in $U$.
\item The support of $\sigma + \sigma'$ is contained in the union of
the supports of $\sigma, \sigma' \in \mathcal{F}(U)$.
\item If $\varphi : \mathcal{F} \to \mathcal{G}$ is a map of
abelian sheaves on $S_\etale$, then the support of $\varphi(\sigma)$
is contained in the support of $\sigma \in \mathcal{F}(U)$.
\item The support of $\mathcal{F}$ is the union of the images of the
supports of all local sections of $\mathcal{F}$.
\item If $\mathcal{F} \to \mathcal{G}$ is surjective then the support
of $\mathcal{G}$ is a subset of the support of $\mathcal{F}$.
\item If $\mathcal{F} \to \mathcal{G}$ is injective then the support
of $\mathcal{F}$ is a subset of the support of $\mathcal{G}$.
\end{enumerate}
\end{lemma}

\begin{proof}
Part (1) holds by definition.
Parts (2) and (3) hold because they holds for the restriction of
$\mathcal{F}$ and $\mathcal{G}$ to $U_{Zar}$, see
Modules, Lemma \ref{modules-lemma-support-section-closed}.
Part (4) is a direct consequence of
Lemma \ref{lemma-zero-over-image} part (3).
Parts (5) and (6) follow from the other parts.
\end{proof}

\begin{lemma}
\label{lemma-support-sheaf-rings-closed}
The support of a sheaf of rings on $S_\etale$ is closed.
\end{lemma}

\begin{proof}
This is true because (according to our conventions)
a ring is $0$ if and only if
$1 = 0$, and hence the support of a sheaf of rings
is the support of the unit section.
\end{proof}




\section{Henselian rings}
\label{section-henselian-ring}

\noindent
We begin by stating a theorem which has already been used many times
in the Stacks project. There are many versions of this result; here we
just state the algebraic version.

\begin{theorem}
\label{theorem-quasi-finite-etale-locally}
Let $A\to B$ be finite type ring map and $\mathfrak p \subset A$ a prime
ideal. Then there exist an \'etale ring map $A \to A'$ and a prime
$\mathfrak p' \subset A'$ lying over $\mathfrak p$ such that
\begin{enumerate}
\item
$\kappa(\mathfrak p) = \kappa(\mathfrak p')$,
\item
$ B \otimes_A A' = B_1\times \ldots \times B_r \times C$,
\item
$ A'\to B_i$ is finite and there exists a unique prime $q_i\subset B_i$ lying
over $\mathfrak p'$, and
\item all irreducible components of the fibre
$\Spec(C \otimes_{A'} \kappa(\mathfrak p'))$ of $C$ over $\mathfrak p'$
have dimension at least 1.
\end{enumerate}
\end{theorem}

\begin{proof}
See Algebra, Lemma \ref{algebra-lemma-etale-makes-quasi-finite-finite}, or
see \cite[Th\'eor\`eme 18.12.1]{EGA4}. For a slew of versions in terms of
morphisms of schemes, see
More on Morphisms, Section \ref{more-morphisms-section-etale-localization}.
\end{proof}

\noindent
Recall Hensel's lemma.
There are many versions of this lemma. Here are two:
\begin{enumerate}
\item[(f)] if $f\in \mathbf{Z}_p[T]$ monic and
$f \bmod p = g_0 h_0$ with $gcd(g_0, h_0) = 1$ then $f$ factors
as $f = gh$ with $\bar g = g_0$ and $\bar h = h_0$,
\item[(r)] if $f \in \mathbf{Z}_p[T]$, monic $a_0 \in \mathbf{F}_p$,
$\bar f(a_0) =0$ but $\bar f'(a_0) \neq 0$
then there exists $a \in \mathbf{Z}_p$ with
$f(a) = 0$ and $\bar a = a_0$.
\end{enumerate}
Both versions are true (we will see this later). The first version
asks for lifts of factorizations into coprime parts,
and the second version asks for lifts of simple roots
modulo the maximal ideal. It turns out that requiring
these conditions for a general local ring are equivalent, and are
equivalent to many other conditions. We use the root lifting
property as the definition of a henselian local ring as it is
often the easiest one to check.

%10.01.09
\begin{definition}
\label{definition-henselian}
(See Algebra, Definition \ref{algebra-definition-henselian}.)
A local ring $(R, \mathfrak m, \kappa)$ is called
{\it henselian} if for all
$f \in R[T]$ monic, for all $a_0 \in \kappa$ such that
$\bar f(a_0) = 0$ and $\bar f'(a_0) \neq 0$, there exists
an $a \in R$ such that $f(a) = 0$ and $a \bmod \mathfrak m = a_0$.
\end{definition}

\noindent
A good example of henselian local rings
to keep in mind is complete local rings.
Recall
(Algebra, Definition \ref{algebra-definition-complete-local-ring})
that a complete local ring is a local ring $(R, \mathfrak m)$ such that
$R \cong \lim_n R/\mathfrak m^n$, i.e., it is complete and separated
for the $\mathfrak m$-adic topology.

\begin{theorem}
\label{theorem-hensel}
Complete local rings are henselian.
\end{theorem}

\begin{proof}
Newton's method. See
Algebra, Lemma \ref{algebra-lemma-complete-henselian}.
\end{proof}

\begin{theorem}
\label{theorem-henselian}
Let $(R, \mathfrak m, \kappa)$ be a local ring. The following are equivalent:
\begin{enumerate}
\item $R$ is henselian,
\item for any $f\in R[T]$ and any factorization $\bar f = g_0 h_0$ in
$\kappa[T]$ with $\gcd(g_0, h_0)=1$, there exists a factorization $f = gh$ in
$R[T]$ with $\bar g = g_0$ and $\bar h = h_0$,
\item any finite $R$-algebra $S$ is isomorphic to a finite product of
local rings finite over $R$,
\item any finite type $R$-algebra $A$ is isomorphic to a product
$A \cong A' \times C$ where $A' \cong A_1 \times \ldots \times A_r$
is a product of finite local $R$-algebras and all the irreducible
components of $C \otimes_R \kappa$ have dimension at least 1,
\item if $A$ is an \'etale $R$-algebra and $\mathfrak n$ is a maximal ideal of
$A$ lying over $\mathfrak m$ such that $\kappa \cong A/\mathfrak n$, then there
exists an isomorphism $\varphi : A \cong R \times A'$ such that
$\varphi(\mathfrak n) = \mathfrak m \times A' \subset R \times A'$.
\end{enumerate}
\end{theorem}

\begin{proof}
This is just a subset of the results from
Algebra, Lemma \ref{algebra-lemma-characterize-henselian}.
Note that part (5) above corresponds to part (8) of
Algebra, Lemma \ref{algebra-lemma-characterize-henselian}
but is formulated slightly differently.
\end{proof}

\begin{lemma}
\label{lemma-finite-over-henselian}
If $R$ is henselian and $A$ is a finite $R$-algebra, then $A$ is a finite
product of henselian local rings.
\end{lemma}

\begin{proof}
See
Algebra, Lemma \ref{algebra-lemma-finite-over-henselian}.
\end{proof}

\begin{definition}
\label{definition-strictly-henselian}
A local ring $R$ is called {\it strictly henselian} if it is henselian and its
residue field is separably closed.
\end{definition}

\begin{example}
\label{example-powerseries}
In the case $R = \mathbf{C}[[t]]$, the \'etale $R$-algebras are finite products
of the trivial extension $R \to R$ and the extensions
$R \to R[X, X^{-1}]/(X^n-t)$.
The latter ones factor through the open $D(t) \subset \Spec(R)$, so any
\'etale covering can be refined by the covering
$\{\text{id} : \Spec(R) \to \Spec(R)\}$. We will see below that
this is a somewhat general fact on \'etale coverings of spectra of henselian
rings. This will show that higher \'etale cohomology of the spectrum of a
strictly henselian ring is zero.
\end{example}

\begin{theorem}
\label{theorem-henselization}
Let $(R, \mathfrak m, \kappa)$ be a local ring and
$\kappa\subset\kappa^{sep}$ a separable algebraic closure.
There exist canonical flat local ring maps $R \to R^h \to R^{sh}$ where
\begin{enumerate}
\item $R^h$, $R^{sh}$ are filtered colimits of \'etale $R$-algebras,
\item $R^h$ is henselian, $R^{sh}$ is strictly henselian,
\item $\mathfrak m R^h$ (resp.\ $\mathfrak m R^{sh}$) is the
maximal ideal of $R^h$ (resp.\ $R^{sh}$), and
\item $\kappa = R^h/\mathfrak m R^h$, and
$\kappa^{sep} = R^{sh}/\mathfrak m R^{sh}$ as extensions of $\kappa$.
\end{enumerate}
\end{theorem}

\begin{proof}
The structure of $R^h$ and $R^{sh}$ is described in
Algebra, Lemmas \ref{algebra-lemma-henselization} and
\ref{algebra-lemma-strict-henselization}.
\end{proof}

\noindent
The rings constructed in Theorem \ref{theorem-henselization}
are called respectively the {\it henselization} and the
{\it strict henselization} of the local ring $R$, see
Algebra, Definition \ref{algebra-definition-henselization}.
Many of the properties of $R$ are reflected in its (strict) henselization,
see More on Algebra,
Section \ref{more-algebra-section-permanence-henselization}.




\section{Stalks of the structure sheaf}
\label{section-stalks-structure-sheaf}

\noindent
In this section we identify the stalk of the structure sheaf at a geometric
point with the strict henselization of the local ring at the corresponding
``usual'' point.

\begin{lemma}
\label{lemma-describe-etale-local-ring}
\begin{slogan}
The stalk of the structure sheaf of a scheme
in the etale topology is the strict henselization.
\end{slogan}
Let $S$ be a scheme.
Let $\overline{s}$ be a geometric point of $S$ lying over $s \in S$.
Let $\kappa = \kappa(s)$ and let
$\kappa \subset \kappa^{sep} \subset \kappa(\overline{s})$ denote
the separable algebraic closure of $\kappa$ in $\kappa(\overline{s})$.
Then there is a canonical identification
$$
(\mathcal{O}_{S, s})^{sh}
\cong
(\mathcal{O}_S)_{\overline{s}}
$$
where the left hand side is the strict henselization of the local ring
$\mathcal{O}_{S, s}$ as described in
Theorem \ref{theorem-henselization}
and right hand side is the stalk of the structure sheaf
$\mathcal{O}_S$ on $S_\etale$ at
the geometric point $\overline{s}$.
\end{lemma}

\begin{proof}
Let $\Spec(A) \subset S$ be an affine neighbourhood of $s$.
Let $\mathfrak p \subset A$ be the prime ideal corresponding to $s$.
With these choices we have canonical isomorphisms
$\mathcal{O}_{S, s} = A_{\mathfrak p}$ and $\kappa(s) = \kappa(\mathfrak p)$.
Thus we have
$\kappa(\mathfrak p) \subset \kappa^{sep} \subset \kappa(\overline{s})$.
Recall that
$$
(\mathcal{O}_S)_{\overline{s}} =
\colim_{(U, \overline{u})} \mathcal{O}(U)
$$
where the limit is over the \'etale neighbourhoods of $(S, \overline{s})$.
A cofinal system is given by those \'etale neighbourhoods $(U, \overline{u})$
such that $U$ is affine and $U \to S$ factors through $\Spec(A)$.
In other words, we see that
$$
(\mathcal{O}_S)_{\overline{s}} = \colim_{(B, \mathfrak q, \phi)} B
$$
where the colimit is over \'etale $A$-algebras $B$ endowed with a prime
$\mathfrak q$ lying over $\mathfrak p$ and a
$\kappa(\mathfrak p)$-algebra map
$\phi : \kappa(\mathfrak q) \to \kappa(\overline{s})$.
Note that since $\kappa(\mathfrak q)$ is finite separable over
$\kappa(\mathfrak p)$ the image of $\phi$ is contained in $\kappa^{sep}$.
Via these translations the result of the lemma is equivalent
to the result of
Algebra, Lemma \ref{algebra-lemma-strict-henselization-different}.
\end{proof}

\begin{definition}
\label{definition-etale-local-rings}
Let $S$ be a scheme. Let $\overline{s}$ be a geometric point of $S$
lying over the point $s \in S$.
\begin{enumerate}
\item The {\it \'etale local ring of $S$ at $\overline{s}$}
is the stalk of the structure sheaf $\mathcal{O}_S$ on $S_\etale$
at $\overline{s}$. We sometimes call this the
{\it strict henselization of $\mathcal{O}_{S, s}$} relative
to the geometric point $\overline{s}$.
Notation used: $\mathcal{O}_{S, \overline{s}}^{sh}$.
\item The {\it henselization of $\mathcal{O}_{S, s}$} is the
henselization of the local ring of $S$ at $s$. See
Algebra, Definition \ref{algebra-definition-henselization},
and
Theorem \ref{theorem-henselization}.
Notation: $\mathcal{O}_{S, s}^h$.
\item The {\it strict henselization of $S$ at $\overline{s}$}
is the scheme $\Spec(\mathcal{O}_{S, \overline{s}}^{sh})$.
\item The {\it henselization of $S$ at $s$} is the scheme
$\Spec(\mathcal{O}_{S, s}^h)$.
\end{enumerate}
\end{definition}

\noindent
Let $f : T \to S$ be a morphism of schemes. Let $\overline{t}$
be a geometric point of $T$ with image $\overline{s}$ in $S$.
Let $t \in T$ and $s \in S$ be their images.
Then we obtain a canonical commutative diagram
$$
\xymatrix{
\Spec(\mathcal{O}^{sh}_{T, \overline{t}}) \ar[r] \ar[d] &
\Spec(\mathcal{O}^h_{T, t}) \ar[r] \ar[d] &
T \ar[d]^f \\
\Spec(\mathcal{O}^{sh}_{S, \overline{s}}) \ar[r] &
\Spec(\mathcal{O}^h_{S, s}) \ar[r] &
S
}
$$
of henselizations and strict henselizations of $T$ and $S$. You can
prove this by choosing affine neighbourhoods of $t$ and $s$ and using
the functoriality of (strict) henselizations given by
Algebra, Lemmas \ref{algebra-lemma-henselian-functorial-improve} and
\ref{algebra-lemma-strictly-henselian-functorial-improve}.

\begin{lemma}
\label{lemma-describe-henselization}
Let $S$ be a scheme. Let $s \in S$. Then we have
$$
\mathcal{O}_{S, s}^h =
\colim_{(U, u)} \mathcal{O}(U)
$$
where the colimit is over the filtered category of
\'etale neighbourhoods $(U, u)$ of $(S, s)$ such that
$\kappa(s) = \kappa(u)$.
\end{lemma}

\begin{proof}
This lemma is a copy of
More on Morphisms, Lemma \ref{more-morphisms-lemma-describe-henselization}.
\end{proof}

\begin{remark}
\label{remark-henselization-Noetherian}
Let $S$ be a scheme. Let $s \in S$.
If $S$ is locally Noetherian then $\mathcal{O}_{S, s}^h$
is also Noetherian and it has the same completion:
$$
\widehat{\mathcal{O}_{S, s}} \cong \widehat{\mathcal{O}_{S, s}^h}.
$$
In particular,
$\mathcal{O}_{S, s} \subset
\mathcal{O}_{S, s}^h \subset
\widehat{\mathcal{O}_{S, s}}$.
The henselization of $\mathcal{O}_{S, s}$ is in general much
smaller than its completion and inherits many of its properties.
For example, if $\mathcal{O}_{S, s}$ is reduced, then so is
$\mathcal{O}_{S, s}^h$, but this is not true for the completion in general.
Insert future references here.
\end{remark}

\begin{lemma}
\label{lemma-etale-site-locally-ringed}
Let $S$ be a scheme. The small \'etale site $S_\etale$ endowed with
its structure sheaf $\mathcal{O}_S$ is a locally ringed site, see
Modules on Sites, Definition \ref{sites-modules-definition-locally-ringed}.
\end{lemma}

\begin{proof}
This follows because the stalks
$(\mathcal{O}_S)_{\overline{s}} = \mathcal{O}^{sh}_{S, \overline{s}}$ are
local, and because $S_\etale$ has enough points, see
Lemma \ref{lemma-describe-etale-local-ring},
Theorem \ref{theorem-exactness-stalks},
and
Remarks \ref{remarks-enough-points}.
See
Modules on Sites, Lemmas \ref{sites-modules-lemma-locally-ringed-stalk} and
\ref{sites-modules-lemma-ringed-stalk-not-zero}
for the fact that this implies the small \'etale site is locally ringed.
\end{proof}






\section{Functoriality of small \'etale topos}
\label{section-functoriality}

\noindent
So far we haven't yet discussed the functoriality of the \'etale site, in
other words what happens when given a morphism of schemes. A precise formal
discussion can be found in
Topologies, Section \ref{topologies-section-etale}.
In this and the next sections we discuss this material briefly specifically
in the setting of small \'etale sites.

\medskip\noindent
Let $f : X \to Y$ be a morphism of schemes. We obtain a functor
\begin{equation}
\label{equation-functorial}
u : Y_\etale \longrightarrow X_\etale, \quad
V/Y \longmapsto X \times_Y V/X.
\end{equation}
This functor has the following important properties
\begin{enumerate}
\item $u(\text{final object}) = \text{final object}$,
\item $u$ preserves fibre products,
\item if $\{V_j \to V\}$ is a covering in $Y_\etale$, then
$\{u(V_j) \to u(V)\}$ is a covering in $X_\etale$.
\end{enumerate}
Each of these is easy to check (omitted). As a consequence we obtain what
is called a {\it morphism of sites}
$$
f_{small} : X_\etale \longrightarrow Y_\etale,
$$
see
Sites, Definition \ref{sites-definition-morphism-sites}
and
Sites, Proposition \ref{sites-proposition-get-morphism}.
It is not necessary to know about the abstract notion in detail
in order to work with \'etale sheaves and \'etale cohomology.
It usually suffices to know that there are functors
$f_{small, *}$ (pushforward) and $f_{small}^{-1}$ (pullback)
on \'etale sheaves, and to know some of their simple properties.
We will discuss these properties in the next sections, but we will
sometimes refer to the more abstract material for proofs since
that is often the natural setting to prove them.


\section{Direct images}
\label{section-direct-image}

\noindent
Let us define the pushforward of a presheaf.

\begin{definition}
\label{definition-direct-image-presheaf}
Let $f: X\to Y$ be a morphism of schemes.
Let $\mathcal{F} $ a presheaf of sets on $X_\etale$.
The {\it direct image}, or {\it pushforward} of $\mathcal{F}$
(under $f$) is
$$
f_*\mathcal{F} : Y_\etale^{opp} \longrightarrow \textit{Sets}, \quad
(V/Y) \longmapsto \mathcal{F}(X \times_Y V/X).
$$
We sometimes write $f_* = f_{small, *}$ to distinguish from other
direct image functors (such as usual Zariski pushforward or $f_{big, *}$).
\end{definition}

\noindent
This is a well-defined \'etale presheaf since the base change of an \'etale
morphism is again \'etale. A more categorical way of saying this is that
$f_*\mathcal{F}$ is the composition of functors $\mathcal{F} \circ u$
where $u$ is as in Equation (\ref{equation-functorial}). This makes it
clear that the construction is functorial in the presheaf
$\mathcal{F}$ and hence we obtain a functor
$$
f_* = f_{small, *} :
\textit{PSh}(X_\etale)
\longrightarrow
\textit{PSh}(Y_\etale)
$$
Note that if $\mathcal{F}$ is a presheaf of abelian groups, then
$f_*\mathcal{F}$ is also a presheaf of abelian groups and we obtain
$$
f_* = f_{small, *} :
\textit{PAb}(X_\etale)
\longrightarrow
\textit{PAb}(Y_\etale)
$$
as before (i.e., defined by exactly the same rule).

\begin{remark}
\label{remark-direct-image-sheaf}
We claim that the direct image of a sheaf is a sheaf.
Namely, if $\{V_j \to V\}$ is an \'etale covering in $Y_\etale$
then $\{X \times_Y V_j \to X \times_Y V\}$ is an \'etale covering in
$X_\etale$. Hence the sheaf condition for $\mathcal{F}$ with respect
to $\{X \times_Y V_i \to X \times_Y V\}$
is equivalent to the sheaf condition for $f_*\mathcal{F}$ with respect to
$\{V_i \to V\}$. Thus if $\mathcal{F}$ is a sheaf, so is
$f_*\mathcal{F}$.
\end{remark}

\begin{definition}
\label{definition-direct-image-sheaf}
Let $f: X\to Y$ be a morphism of schemes.
Let $\mathcal{F} $ a sheaf of sets on $X_\etale$.
The {\it direct image}, or {\it pushforward} of $\mathcal{F}$
(under $f$) is
$$
f_*\mathcal{F} : Y_\etale^{opp} \longrightarrow \textit{Sets}, \quad
(V/Y) \longmapsto \mathcal{F}(X \times_Y V/X)
$$
which is a sheaf by
Remark \ref{remark-direct-image-sheaf}.
We sometimes write $f_* = f_{small, *}$ to distinguish from other
direct image functors (such as usual Zariski pushforward or $f_{big, *}$).
\end{definition}

\noindent
The exact same discussion as above applies and we obtain functors
$$
f_* = f_{small, *} :
\Sh(X_\etale)
\longrightarrow
\Sh(Y_\etale)
$$
and
$$
f_* = f_{small, *} :
\textit{Ab}(X_\etale)
\longrightarrow
\textit{Ab}(Y_\etale)
$$
called {\it direct image} again.

\medskip\noindent
The functor $f_*$ on abelian sheaves is left exact. (See
Homology, Section \ref{homology-section-functors}
for what it means for a functor between abelian categories to be left exact.)
Namely, if
$0 \to \mathcal{F}_1 \to \mathcal{F}_2 \to \mathcal{F}_3$
is exact on $X_\etale$, then for every
$U/X \in \Ob(X_\etale)$
the sequence of abelian groups
$0 \to \mathcal{F}_1(U) \to \mathcal{F}_2(U) \to \mathcal{F}_3(U)$
is exact. Hence for every $V/Y \in \Ob(Y_\etale)$
the sequence of abelian groups
$0 \to f_*\mathcal{F}_1(V) \to f_*\mathcal{F}_2(V) \to f_*\mathcal{F}_3(V)$
is exact, because this is the previous sequence with $U = X \times_Y V$.

\begin{definition}
\label{definition-higher-direct-images}
Let $f: X \to Y$ be a morphism of schemes.
The right derived functors $\{R^pf_*\}_{p \geq 1}$ of
$f_* : \textit{Ab}(X_\etale) \to \textit{Ab}(Y_\etale)$
are called {\it higher direct images}.
\end{definition}

\noindent
The higher direct images and their derived category variants are
discussed in more detail in (insert future reference here).



\section{Inverse image}
\label{section-inverse-image}

\noindent
In this section we briefly discuss pullback of sheaves on the small
\'etale sites. The precise construction of this is in
Topologies, Section \ref{topologies-section-etale}.

\begin{definition}
\label{definition-inverse-image}
Let $f: X\to Y$ be a morphism of schemes. The {\it inverse image}, or
{\it pullback}\footnote{We use the notation $f^{-1}$ for pullbacks of
sheaves of sets or sheaves of abelian groups, and we reserve $f^*$ for
pullbacks of sheaves of modules via a morphism of ringed sites/topoi.}
functors are the functors
$$
f^{-1} = f_{small}^{-1} :
\Sh(Y_\etale)
\longrightarrow
\Sh(X_\etale)
$$
and
$$
f^{-1} = f_{small}^{-1} :
\textit{Ab}(Y_\etale)
\longrightarrow
\textit{Ab}(X_\etale)
$$
which are left adjoint to $f_* = f_{small, *}$. Thus
$f^{-1}$ is characterized by the fact that
$$
\Hom_{{\Sh(X_\etale)}} (f^{-1}\mathcal{G}, \mathcal{F})
=
\Hom_{\Sh(Y_\etale)} (\mathcal{G}, f_*\mathcal{F})
$$
functorially, for any $\mathcal{F} \in \Sh(X_\etale)$ and
$\mathcal{G} \in \Sh(Y_\etale)$. We similarly have
$$
\Hom_{{\textit{Ab}(X_\etale)}} (f^{-1}\mathcal{G}, \mathcal{F})
=
\Hom_{\textit{Ab}(Y_\etale)} (\mathcal{G}, f_*\mathcal{F})
$$
for $\mathcal{F} \in \textit{Ab}(X_\etale)$ and
$\mathcal{G} \in \textit{Ab}(Y_\etale)$.
\end{definition}

\noindent
It is not trivial that such an adjoint exists.
On the other hand, it exists in a fairly general setting, see
Remark \ref{remark-functoriality-general}
below. The general machinery shows that $f^{-1}\mathcal{G}$
is the sheaf associated to the presheaf
\begin{equation}
\label{equation-pullback}
U/X
\longmapsto
\colim_{U \to X \times_Y V} \mathcal{G}(V/Y)
\end{equation}
where the colimit is over the category of pairs
$(V/Y, \varphi : U/X \to X \times_Y V/X)$.
To see this apply
Sites, Proposition \ref{sites-proposition-get-morphism}
to the functor $u$ of Equation (\ref{equation-functorial})
and use the description of $u_s = (u_p\ )^\#$ in
Sites, Sections \ref{sites-section-continuous-functors} and
\ref{sites-section-functoriality-PSh}.
We will occasionally use this formula for the pullback
in order to prove some of its basic properties.

\begin{lemma}
\label{lemma-stalk-pullback}
Let $f : X \to Y$ be a morphism of schemes.
\begin{enumerate}
\item The functor
$f^{-1} : \textit{Ab}(Y_\etale) \to \textit{Ab}(X_\etale)$
is exact.
\item The functor
$f^{-1} : \Sh(Y_\etale) \to \Sh(X_\etale)$
is exact, i.e., it commutes with finite limits and colimits, see
Categories, Definition \ref{categories-definition-exact}.
\item Let $\overline{x} \to X$ be a geometric point.
Let $\mathcal{G}$ be a sheaf on $Y_\etale$.
Then there is a canonical identification
$$
(f^{-1}\mathcal{G})_{\overline{x}} = \mathcal{G}_{\overline{y}}.
$$
where $\overline{y} = f \circ \overline{x}$.
\item For any $V \to Y$ \'etale we have $f^{-1}h_V = h_{X \times_Y V}$.
\end{enumerate}
\end{lemma}

\begin{proof}
The exactness of $f^{-1}$ on sheaves of sets is a consequence of
Sites, Proposition \ref{sites-proposition-get-morphism}
applied to our functor $u$ of Equation (\ref{equation-functorial}).
In fact the exactness of pullback is part of the definition of
a morphism of topoi (or sites if you like). Thus we see (2) holds.
It implies part (1) since given an abelian sheaf $\mathcal{G}$ on
$Y_\etale$
the underlying sheaf of sets of $f^{-1}\mathcal{F}$ is the same
as $f^{-1}$ of the underlying sheaf of sets of $\mathcal{F}$, see
Sites, Section \ref{sites-section-sheaves-algebraic-structures}.
See also
Modules on Sites, Lemma \ref{sites-modules-lemma-flat-pullback-exact}.
In the literature (1) and (2) are sometimes deduced from (3) via
Theorem \ref{theorem-exactness-stalks}.

\medskip\noindent
Part (3) is a general fact about stalks of pullbacks, see
Sites, Lemma \ref{sites-lemma-point-morphism-sites}.
We will also prove (3) directly as follows. Note that by
Lemma \ref{lemma-stalk-exact}
taking stalks commutes with sheafification.
Now recall that $f^{-1}\mathcal{G}$ is the sheaf
associated to the presheaf
$$
U \longrightarrow \colim_{U \to X \times_Y V} \mathcal{G}(V),
$$
see Equation (\ref{equation-pullback}).
Thus we have
\begin{align*}
(f^{-1}\mathcal{G})_{\overline{x}}
& = \colim_{(U, \overline{u})} f^{-1}\mathcal{G}(U) \\
& = \colim_{(U, \overline{u})}
\colim_{a : U \to X \times_Y V} \mathcal{G}(V) \\
& = \colim_{(V, \overline{v})} \mathcal{G}(V) \\
& = \mathcal{G}_{\overline{y}}
\end{align*}
in the third equality the pair $(U, \overline{u})$ and the map
$a : U \to X \times_Y V$ corresponds to the pair $(V, a \circ \overline{u})$.

\medskip\noindent
Part (4) can be proved in a similar manner by identifying the colimits
which define $f^{-1}h_V$. Or you can use
Yoneda's lemma (Categories, Lemma \ref{categories-lemma-yoneda})
and the functorial equalities
$$
\Mor_{\Sh(X_\etale)}(f^{-1}h_V, \mathcal{F}) =
\Mor_{\Sh(Y_\etale)}(h_V, f_*\mathcal{F}) =
f_*\mathcal{F}(V) = \mathcal{F}(X \times_Y V)
$$
combined with the fact that representable presheaves are sheaves. See also
Sites, Lemma \ref{sites-lemma-pullback-representable-sheaf}
for a completely general result.
\end{proof}

\noindent
The pair of functors $(f_*, f^{-1})$ define a morphism of small \'etale topoi
$$
f_{small} :
\Sh(X_\etale)
\longrightarrow
\Sh(Y_\etale)
$$
Many generalities on cohomology of sheaves hold for topoi and
morphisms of topoi. We will try to point out when results are
general and when they are specific to the \'etale topos.

\begin{remark}
\label{remark-functoriality-general}
More generally, let $\mathcal{C}_1, \mathcal{C}_2$ be sites, and
assume they have final objects and fibre products. Let
$u: \mathcal{C}_2 \to \mathcal{C}_1$ be a functor satisfying:
\begin{enumerate}
\item if $\{V_i \to V\}$ is a covering of $\mathcal{C}_2$, then
$\{u(V_i) \to u(V)\}$ is a covering of $\mathcal{C}_1$ (we
say that $u$ is {\it continuous}), and
\item $u$ commutes with finite limits (i.e., $u$ is left exact, i.e.,
$u$ preserves fibre products and final objects).
\end{enumerate}
Then one can define
$f_*: \Sh(\mathcal{C}_1) \to \Sh(\mathcal{C}_2)$
by $ f_* \mathcal{F}(V) = \mathcal{F}(u(V))$.
Moreover, there exists an exact functor $f^{-1}$ which
is left adjoint to $f_*$, see
Sites, Definition \ref{sites-definition-morphism-sites} and
Proposition \ref{sites-proposition-get-morphism}.
Warning: It is not enough to require simply that $u$ is continuous
and commutes with fibre products in order to get a morphism of topoi.
\end{remark}




\section{Functoriality of big topoi}
\label{section-functoriality-big-topoi}

\noindent
Given a morphism of schemes $f : X \to Y$ there are a whole host of
morphisms of topoi associated to $f$, see
Topologies, Section \ref{topologies-section-change-topologies}
for a list. Perhaps the most used ones are the morphisms of topoi
$$
f_{big} = f_{big, \tau} :
\Sh((\Sch/X)_\tau)
\longrightarrow
\Sh((\Sch/Y)_\tau)
$$
where $\tau \in \{Zariski, \etale, smooth, syntomic, fppf\}$.
These each correspond to a continuous functor
$$
(\Sch/Y)_\tau \longrightarrow (\Sch/X)_\tau, \quad
V/Y \longmapsto X \times_Y V/X
$$
which preserves final objects, fibre products and covering, and hence
defines a morphism of sites
$$
f_{big} : (\Sch/X)_\tau \longrightarrow (\Sch/Y)_\tau.
$$
See
Topologies, Sections \ref{topologies-section-zariski},
\ref{topologies-section-etale},
\ref{topologies-section-smooth},
\ref{topologies-section-syntomic}, and
\ref{topologies-section-fppf}.
In particular, pushforward along $f_{big}$ is given by the rule
$$
(f_{big, *}\mathcal{F})(V/Y) = \mathcal{F}(X \times_Y V/X)
$$
It turns out that these morphisms of topoi have an inverse
image functor $f_{big}^{-1}$ which is very easy to describe.
Namely, we have
$$
(f_{big}^{-1}\mathcal{G})(U/X) = \mathcal{G}(U/Y)
$$
where the structure morphism of $U/Y$ is the composition of the
structure morphism $U \to X$ with $f$, see
Topologies, Lemmas \ref{topologies-lemma-morphism-big},
\ref{topologies-lemma-morphism-big-etale},
\ref{topologies-lemma-morphism-big-smooth},
\ref{topologies-lemma-morphism-big-syntomic}, and
\ref{topologies-lemma-morphism-big-fppf}.







\section{Functoriality and sheaves of modules}
\label{section-morphisms-modules}

\noindent
In this section we are going to reformulate some of the material explained in
Descent, Sections \ref{descent-section-quasi-coherent-sheaves},
\ref{descent-section-quasi-coherent-cohomology}, and
\ref{descent-section-quasi-coherent-sheaves-bis}
in the setting of \'etale topologies. Let $f : X \to Y$ be a morphism of
schemes. We have seen above, see
Sections \ref{section-functoriality}, \ref{section-direct-image}, and
\ref{section-inverse-image}
that this induces a morphism $f_{small}$ of small \'etale sites. In
Descent, Remark \ref{descent-remark-change-topologies-ringed}
we have seen that $f$ also induces a natural map
$$
f_{small}^\sharp :
\mathcal{O}_{Y_\etale}
\longrightarrow
f_{small, *}\mathcal{O}_{X_\etale}
$$
of sheaves of rings on
$Y_\etale$ such that $(f_{small}, f_{small}^\sharp)$
is a morphism of ringed sites. See
Modules on Sites, Definition \ref{sites-modules-definition-ringed-site}
for the definition of a morphism of ringed sites.
Let us just recall here that $f_{small}^\sharp$ is defined by the
compatible system of maps
$$
\text{pr}_V^\sharp : \mathcal{O}(V) \longrightarrow \mathcal{O}(X \times_Y V)
$$
for $V$ varying over the objects of $Y_\etale$.

\medskip\noindent
It is clear that this construction is compatible with compositions of
morphisms of schemes. More precisely, if $f : X \to Y$ and $g : Y \to Z$
are morphisms of schemes, then we have
$$
(g_{small}, g_{small}^\sharp)
\circ
(f_{small}, f_{small}^\sharp)
=
((g \circ f)_{small}, (g \circ f)_{small}^\sharp)
$$
as morphisms of ringed topoi. Moreover, by
Modules on Sites, Definition \ref{sites-modules-definition-pushforward}
we see that given a morphism $f : X \to Y$ of schemes
we get well defined pullback and direct image functors
\begin{align*}
f_{small}^* :
\textit{Mod}(\mathcal{O}_{Y_\etale})
\longrightarrow
\textit{Mod}(\mathcal{O}_{X_\etale}), \\
f_{small, *} :
\textit{Mod}(\mathcal{O}_{X_\etale})
\longrightarrow
\textit{Mod}(\mathcal{O}_{Y_\etale})
\end{align*}
which are adjoint in the usual way. If $g : Y \to Z$ is another morphism
of schemes, then we have
$(g \circ f)_{small}^* = f_{small}^* \circ g_{small}^*$
and $(g \circ f)_{small, *} = g_{small, *} \circ f_{small, *}$
because of what we said about compositions.

\medskip\noindent
There is quite a bit of difference between the category
of all $\mathcal{O}_X$ modules on $X$ and the category between all
$\mathcal{O}_{X_\etale}$-modules on $X_\etale$. But the
results of
Descent, Sections \ref{descent-section-quasi-coherent-sheaves},
\ref{descent-section-quasi-coherent-cohomology}, and
\ref{descent-section-quasi-coherent-sheaves-bis}
tell us that there is not much difference between considering quasi-coherent
modules on $S$ and quasi-coherent modules on $S_\etale$.
(We have already seen this in
Theorem \ref{theorem-quasi-coherent}
for example.)
In particular, if $f : X \to Y$ is any morphism of schemes, then
the pullback functors $f_{small}^*$ and $f^*$ match for
quasi-coherent sheaves, see
Descent,
Proposition \ref{descent-proposition-equivalence-quasi-coherent-functorial}.
Moreover, the same is true for pushforward provided $f$ is
quasi-compact and quasi-separated, see
Descent, Lemma \ref{descent-lemma-higher-direct-images-small-etale}.

\medskip\noindent
A few words about functoriality of the structure sheaf on big sites.
Let $f : X \to Y$ be a morphism of schemes. Choose any of the
topologies $\tau \in \{Zariski,\linebreak[0]
\etale,\linebreak[0] smooth,\linebreak[0] syntomic,\linebreak[0]
fppf\}$. Then the morphism
$f_{big} : (\Sch/X)_\tau \to (\Sch/Y)_\tau$
becomes a morphism of ringed sites by a map
$$
f_{big}^\sharp : \mathcal{O}_Y \longrightarrow f_{big, *}\mathcal{O}_X
$$
see Descent, Remark \ref{descent-remark-change-topologies-ringed}.
In fact it is given by the same construction as in the case of small
sites explained above.












\section{Comparing topologies}
\label{section-compare-topologies}

\noindent
In this section we start studying what happens when you compare
sheaves with respect to different topologies.

\begin{lemma}
\label{lemma-where-sections-are-equal}
Let $S$ be a scheme. Let $\mathcal{F}$ be a sheaf of sets on $S_\etale$.
Let $s, t \in \mathcal{F}(S)$. Then there exists an open $W \subset S$
characterized by the following property: A morphism $f : T \to S$
factors through $W$ if and only if $s|_T = t|_T$ (restriction is
pullback by $f_{small}$).
\end{lemma}

\begin{proof}
Consider the presheaf which assigns to $U \in \Ob(S_\etale)$ the empty set
if $s|_U \not = t|_U$ and a singleton else. It is clear that this is
a subsheaf of the final object of $\Sh(S_\etale)$. By
Lemma \ref{lemma-support-subsheaf-final}
we find an open $W \subset S$ representing this presheaf.
For a geometric point $\overline{x}$ of $S$ we see that $\overline{x} \in W$
if and only if the stalks of $s$ and $t$ at $\overline{x}$ agree.
By the description of stalks of pullbacks in
Lemma \ref{lemma-stalk-pullback}
we see that $W$ has the desired property.
\end{proof}

\begin{lemma}
\label{lemma-describe-pullback}
Let $S$ be a scheme. Let $\tau \in \{Zariski, \etale\}$. Consider the morphism
$$
\pi_S : (\Sch/S)_\tau \longrightarrow S_\tau
$$
of Topologies, Lemma \ref{topologies-lemma-at-the-bottom} or
\ref{topologies-lemma-at-the-bottom-etale}. Let $\mathcal{F}$ be a sheaf on
$S_\tau$. Then $\pi_S^{-1}\mathcal{F}$ is given by the rule
$$
(\pi_S^{-1}\mathcal{F})(T) = \Gamma(T_\tau, f_{small}^{-1}\mathcal{F})
$$
where $f : T \to S$. Moreover, $\pi_S^{-1}\mathcal{F}$ satisfies the
sheaf condition with respect to fpqc coverings.
\end{lemma}

\begin{proof}
Observe that we have a morphism $i_f : \Sh(T_\tau) \to \Sh(\Sch/S)_\tau)$
such that $\pi_S \circ i_f = f_{small}$ as morphisms
$T_\tau \to S_\tau$, see
Topologies, Lemmas \ref{topologies-lemma-put-in-T},
\ref{topologies-lemma-morphism-big-small},
\ref{topologies-lemma-put-in-T-etale}, and
\ref{topologies-lemma-morphism-big-small-etale}.
Since pullback is transitive we see that
$i_f^{-1} \pi_S^{-1}\mathcal{F} = f_{small}^{-1}\mathcal{F}$ as desired.

\medskip\noindent
Let $\{g_i : T_i \to T\}_{i \in I}$ be an fpqc covering.
The final statement means the following: Given a sheaf $\mathcal{G}$
on $T_\tau$ and given sections
$s_i \in \Gamma(T_i, g_{i, small}^{-1}\mathcal{G})$ whose pullbacks
to $T_i \times_T T_j$ agree, there is a unique section $s$ of $\mathcal{G}$
over $T$ whose pullback to $T_i$ agrees with $s_i$.

\medskip\noindent
Let $V \to T$ be an object of $T_\tau$ and let $t \in \mathcal{G}(V)$.
For every $i$ there is a largest open $W_i \subset T_i \times_T V$
such that the pullbacks of $s_i$ and $t$ agree as sections of the pullback
of $\mathcal{G}$ to $W_i \subset T_i \times_T V$, see
Lemma \ref{lemma-where-sections-are-equal}.
Because $s_i$ and $s_j$ agree over $T_i \times_T T_j$ we find
that $W_i$ and $W_j$ pullback to the same open over
$T_i \times_T T_j \times_T V$. By
Descent, Lemma \ref{descent-lemma-open-fpqc-covering}
we find an open $W \subset V$ whose inverse image to $T_i \times_T V$
recovers $W_i$. 

\medskip\noindent
By construction of $g_{i, small}^{-1}\mathcal{G}$ there exists
a $\tau$-covering $\{T_{ij} \to T_i\}_{j \in J_i}$, for each $j$ an
open immersion or \'etale morphism $V_{ij} \to T$, a section
$t_{ij} \in \mathcal{G}(V_{ij})$, and commutative diagrams
$$
\xymatrix{
T_{ij} \ar[r] \ar[d] & V_{ij} \ar[d] \\
T_i \ar[r] & T
}
$$
such that $s_i|_{T_{ij}}$ is the pullback of $t_{ij}$. In other words,
after replacing the covering $\{T_i \to T\}$ by $\{T_{ij} \to T\}$
we may assume there are factorizations $T_i \to V_i \to T$ with
$V_i \in \Ob(T_\tau)$ and sections $t_i \in \mathcal{G}(V_i)$
pulling back to $s_i$ over $T_i$.
By the result of the previous paragraph we find opens $W_i \subset V_i$
such that $t_i|_{W_i}$ ``agrees with'' every $s_j$ over $T_j \times_T W_i$.
Note that $T_i \to V_i$ factors through $W_i$.
Hence $\{W_i \to T\}$ is a $\tau$-covering and the lemma is proven.
\end{proof}

\begin{lemma}
\label{lemma-sections-upstairs}
Let $S$ be a scheme. Let $f : T \to S$ be a morphism such that
\begin{enumerate}
\item $f$ is flat and quasi-compact, and
\item the geometric fibres of $f$ are connected.
\end{enumerate}
Let $\mathcal{F}$ be a sheaf on $S_\etale$.
Then $\Gamma(S, \mathcal{F}) = \Gamma(T, f^{-1}_{small}\mathcal{F})$.
\end{lemma}

\begin{proof}
There is a canonical map
$\Gamma(S, \mathcal{F}) \to \Gamma(T, f_{small}^{-1}\mathcal{F})$.
Since $f$ is surjective (because its fibres are connected) we see that
this map is injective.

\medskip\noindent
To show that the map is surjective, let
$\alpha \in \Gamma(T, f_{small}^{-1}\mathcal{F})$.
Since $\{T \to S\}$ is an fpqc covering we can use
Lemma \ref{lemma-describe-pullback} to see that suffices to prove that
$\alpha$ pulls back to the same section over $T \times_S T$ by the
two projections. Let $\overline{s} \to S$ be a geometric point.
It suffices to show the agreement holds over $(T \times_S T)_{\overline{s}}$
as every geometric point of $T \times_S T$ is contained in one of
these geometric fibres. In other words, we are trying to show that
$\alpha|_{T_{\overline{s}}}$ pulls back to the same section over
$$
(T \times_S T)_{\overline{s}} =
T_{\overline{s}} \times_{\overline{s}} T_{\overline{s}}
$$
by the two projections to $T_{\overline{s}}$.
However, since $\mathcal{F}|_{T_{\overline{s}}}$ is the
pullback of $\mathcal{F}|_{\overline{s}}$ it is a constant sheaf
with value $\mathcal{F}_{\overline{s}}$. Since $T_{\overline{s}}$
is connected by assumption, any section of a constant sheaf is constant.
Hence $\alpha|_{T_{\overline{s}}}$ corresponds to an element
of $\mathcal{F}_{\overline{s}}$. Thus the two pullbacks to
$(T \times_S T)_{\overline{s}}$ both
correspond to this same element and we conclude.
\end{proof}

\noindent
Here is a version of Lemma \ref{lemma-sections-upstairs}
where we do not assume that the morphism is flat.

\begin{lemma}
\label{lemma-sections-upstairs-submersive}
Let $S$ be a scheme. Let $f : X \to S$ be a morphism such that
\begin{enumerate}
\item $f$ is submersive, and
\item the geometric fibres of $f$ are connected.
\end{enumerate}
Let $\mathcal{F}$ be a sheaf on $S_\etale$.
Then $\Gamma(S, \mathcal{F}) = \Gamma(X, f^{-1}_{small}\mathcal{F})$.
\end{lemma}

\begin{proof}
There is a canonical map
$\Gamma(S, \mathcal{F}) \to \Gamma(X, f_{small}^{-1}\mathcal{F})$.
Since $f$ is surjective (because its fibres are connected) we see that
this map is injective.

\medskip\noindent
To show that the map is surjective, let
$\tau \in \Gamma(X, f_{small}^{-1}\mathcal{F})$.
It suffices to find
an \'etale covering $\{U_i \to S\}$ and
sections $\sigma_i \in \mathcal{F}(U_i)$
such that $\sigma_i$ pulls back to $\tau|_{X \times_S U_i}$.
Namely, the injectivity shown above guarantees
that $\sigma_i$ and $\sigma_j$ restrict to the same
section of $\mathcal{F}$ over $U_i \times_S U_j$.
Thus we obtain a unique section $\sigma \in \mathcal{F}(S)$
which restricts to $\sigma_i$ over $U_i$.
Then the pullback of $\sigma$ to $X$ is $\tau$
because this is true locally.

\medskip\noindent
Let $\overline{x}$ be a geometric point of $X$ with image $\overline{s}$
in $S$. Consider the image of $\tau$ in the stalk
$$
(f_{small}^{-1}\mathcal{F})_{\overline{x}} = \mathcal{F}_{\overline{s}}
$$
See Lemma \ref{lemma-stalk-pullback}.
We can find an \'etale neighbourhood $U \to S$ of $\overline{s}$
and a section $\sigma \in \mathcal{F}(U)$ mapping to this image
in the stalk. Thus after replacing $S$ by $U$ and $X$ by $X \times_S U$
we may assume there exits a section $\sigma$ of $\mathcal{F}$ over $S$
whose image in $(f_{small}^{-1}\mathcal{F})_{\overline{x}}$ is the
same as $\tau$.

\medskip\noindent
By Lemma \ref{lemma-where-sections-are-equal}
there exists a maximal open $W \subset X$ such that
$f_{small}^{-1}\sigma$ and $\tau$ agree over $W$
and the formation of $W$ commutes with further pullback.
Observe that the pullback of $\mathcal{F}$ to the
geometric fibre $X_{\overline{s}}$ is the pullback
of $\mathcal{F}_{\overline{s}}$ viewed as a sheaf on
$\overline{s}$ by $X_{\overline{s}} \to \overline{s}$.
Hence we see that $\tau$ and $\sigma$ give sections
of the constant sheaf with value $\mathcal{F}_{\overline{s}}$
on $X_{\overline{s}}$ which agree in one point. Since
$X_{\overline{s}}$ is connected by assumption, we conclude
that $W$ contains $X_s$. The same argument for different
geometric fibres shows that $W$ contains every fibre it meets.
Since $f$ is submersive, we conclude that $W$ is the inverse
image of an open neighbourhood of $s$ in $S$.
This finishes the proof.
\end{proof}

\begin{lemma}
\label{lemma-sections-base-field-extension}
Let $K/k$ be an extension of fields with $k$ separably
algebraically closed. Let $S$ be a scheme over $k$. Denote
$p : S_K = S \times_{\Spec(k)} \Spec(K) \to S$ the projection.
Let $\mathcal{F}$ be a sheaf on $S_\etale$.
Then $\Gamma(S, \mathcal{F}) = \Gamma(S_K, p^{-1}_{small}\mathcal{F})$.
\end{lemma}

\begin{proof}
Follows from Lemma \ref{lemma-sections-upstairs}. Namely, it is clear
that $p$ is flat and quasi-compact as the base change of
$\Spec(K) \to \Spec(k)$. On the other hand, if $\overline{s} : \Spec(L) \to S$
is a geometric point, then the fibre of $p$ over $\overline{s}$
is the spectrum of $K \otimes_k L$ which is irreducible hence connected by
Algebra, Lemma \ref{algebra-lemma-separably-closed-irreducible}.
\end{proof}









\section{Recovering morphisms}
\label{section-morphisms}

\noindent
In this section we prove that the rule which associates to a scheme
its locally ringed small \'etale topos is fully faithful in a suitable
sense, see
Theorem \ref{theorem-fully-faithful}.

\begin{lemma}
\label{lemma-morphism-locally-ringed}
Let $f : X \to Y$ be a morphism of schemes.
The morphism of ringed sites $(f_{small}, f_{small}^\sharp)$
associated to $f$ is a morphism of locally ringed sites, see
Modules on Sites,
Definition \ref{sites-modules-definition-morphism-locally-ringed-topoi}.
\end{lemma}

\begin{proof}
Note that the assertion makes sense since we have seen that
$(X_\etale, \mathcal{O}_{X_\etale})$ and
$(Y_\etale, \mathcal{O}_{Y_\etale})$
are locally ringed sites, see
Lemma \ref{lemma-etale-site-locally-ringed}.
Moreover, we know that $X_\etale$ has enough points, see
Theorem \ref{theorem-exactness-stalks}
and
Remarks \ref{remarks-enough-points}.
Hence it suffices to prove that $(f_{small}, f_{small}^\sharp)$
satisfies condition (3) of
Modules on Sites,
Lemma \ref{sites-modules-lemma-locally-ringed-morphism}.
To see this take a point $p$ of $X_\etale$. By
Lemma \ref{lemma-points-small-etale-site}
$p$ corresponds to a geometric point $\overline{x}$ of $X$.
By
Lemma \ref{lemma-stalk-pullback}
the point $q = f_{small} \circ p$ corresponds to the
geometric point $\overline{y} = f \circ \overline{x}$ of $Y$.
Hence the assertion we have to prove is that the induced map
of stalks
$$
(\mathcal{O}_Y)_{\overline{y}} \longrightarrow (\mathcal{O}_X)_{\overline{x}}
$$
is a local ring map. Suppose that $a \in (\mathcal{O}_Y)_{\overline{y}}$
is an element of the left hand side which maps to an element of the maximal
ideal of the right hand side. Suppose that $a$ is the equivalence class
of a triple $(V, \overline{v}, a)$ with $V \to Y$ \'etale,
$\overline{v} : \overline{x} \to V$ over $Y$, and $a \in \mathcal{O}(V)$.
It maps to the equivalence class of
$(X \times_Y V, \overline{x} \times \overline{v}, \text{pr}_V^\sharp(a))$
in the local ring $(\mathcal{O}_X)_{\overline{x}}$. But it is clear that
being in the maximal ideal means that pulling back $\text{pr}_V^\sharp(a)$
to an element of $\kappa(\overline{x})$ gives zero. Hence also pulling back
$a$ to $\kappa(\overline{x})$ is zero. Which means that $a$ lies in the
maximal ideal of $(\mathcal{O}_Y)_{\overline{y}}$.
\end{proof}

\begin{lemma}
\label{lemma-2-morphism}
Let $X$, $Y$ be schemes. Let $f : X \to Y$ be a morphism of schemes.
Let $t$ be a $2$-morphism from $(f_{small}, f_{small}^\sharp)$ to itself, see
Modules on Sites,
Definition \ref{sites-modules-definition-2-morphism-ringed-topoi}.
Then $t = \text{id}$.
\end{lemma}

\begin{proof}
This means that $t : f^{-1}_{small} \to f^{-1}_{small}$
is a transformation of functors such that the diagram
$$
\xymatrix{
f_{small}^{-1}\mathcal{O}_Y
\ar[rd]_{f_{small}^\sharp}  & &
f_{small}^{-1}\mathcal{O}_Y \ar[ll]^t \ar[ld]^{f_{small}^\sharp} \\
& \mathcal{O}_X
}
$$
is commutative. Suppose $V \to Y$ is \'etale with $V$ affine. By
Morphisms, Lemma \ref{morphisms-lemma-quasi-affine-finite-type-over-S}
we may choose an immersion $i : V \to \mathbf{A}^n_Y$ over $Y$.
In terms of sheaves this means that $i$ induces an injection
$h_i : h_V \to \prod_{j = 1, \ldots, n} \mathcal{O}_Y$ of sheaves.
The base change $i'$ of $i$ to $X$ is an immersion
(Schemes, Lemma \ref{schemes-lemma-base-change-immersion}).
Hence $i' : X \times_Y V \to \mathbf{A}^n_X$ is an immersion, which
in turn means that
$h_{i'} : h_{X \times_Y V} \to \prod_{j = 1, \ldots, n} \mathcal{O}_X$
is an injection of sheaves.
Via the identification $f_{small}^{-1}h_V = h_{X \times_Y V}$ of
Lemma \ref{lemma-stalk-pullback}
the map $h_{i'}$ is equal to
$$
\xymatrix{
f_{small}^{-1}h_V \ar[r]^-{f^{-1}h_i} &
\prod_{j = 1, \ldots, n} f_{small}^{-1}\mathcal{O}_Y
\ar[r]^{\prod f^\sharp} &
\prod_{j = 1, \ldots, n} \mathcal{O}_X
}
$$
(verification omitted). This means that the map
$t : f_{small}^{-1}h_V \to f_{small}^{-1}h_V$
fits into the commutative diagram
$$
\xymatrix{
f_{small}^{-1}h_V \ar[r]^-{f^{-1}h_i} \ar[d]^t &
\prod_{j = 1, \ldots, n} f_{small}^{-1}\mathcal{O}_Y
\ar[r]^-{\prod f^\sharp} \ar[d]^{\prod t} &
\prod_{j = 1, \ldots, n} \mathcal{O}_X \ar[d]^{\text{id}}\\
f_{small}^{-1}h_V \ar[r]^-{f^{-1}h_i} &
\prod_{j = 1, \ldots, n} f_{small}^{-1}\mathcal{O}_Y
\ar[r]^-{\prod f^\sharp} &
\prod_{j = 1, \ldots, n} \mathcal{O}_X
}
$$
The commutativity of the right square holds by our assumption on $t$
explained above.
Since the composition of the horizontal arrows is injective
by the discussion above we conclude that the left vertical arrow
is the identity map as well. Any sheaf of sets on
$Y_\etale$ admits a surjection from a (huge) coproduct of sheaves
of the form $h_V$ with $V$ affine (combine
Topologies, Lemma \ref{topologies-lemma-alternative}
with
Sites, Lemma \ref{sites-lemma-sheaf-coequalizer-representable}).
Thus we conclude that $t : f_{small}^{-1} \to f_{small}^{-1}$
is the identity transformation as desired.
\end{proof}

\begin{lemma}
\label{lemma-faithful}
Let $X$, $Y$ be schemes.
Any two morphisms $a, b : X \to Y$ of schemes
for which there exists a $2$-isomorphism
$(a_{small}, a_{small}^\sharp) \cong (b_{small}, b_{small}^\sharp)$
in the $2$-category of ringed topoi are equal.
\end{lemma}

\begin{proof}
Let us argue this carefully since it is a bit confusing.
Let $t : a_{small}^{-1} \to b_{small}^{-1}$ be the $2$-isomorphism.
Consider any open $V \subset Y$. Note that $h_V$ is a subsheaf of the
final sheaf $*$. Thus both $a_{small}^{-1}h_V = h_{a^{-1}(V)}$
and $b_{small}^{-1}h_V = h_{b^{-1}(V)}$ are subsheaves of the final sheaf.
Thus the isomorphism
$$
t : a_{small}^{-1}h_V = h_{a^{-1}(V)} \to b_{small}^{-1}h_V = h_{b^{-1}(V)}
$$
has to be the identity, and $a^{-1}(V) = b^{-1}(V)$.
It follows that $a$ and $b$ are equal on underlying topological spaces.
Next, take a section $f \in \mathcal{O}_Y(V)$. This determines and
is determined by a map of sheaves of sets
$f : h_V \to \mathcal{O}_Y$.
Pull this back and apply $t$ to get a commutative diagram
$$
\xymatrix{
h_{b^{-1}(V)} \ar@{=}[r] &
b_{small}^{-1}h_V \ar[d]^{b_{small}^{-1}(f)} & &
a_{small}^{-1}h_V \ar[d]^{a_{small}^{-1}(f)} \ar[ll]^t &
h_{a^{-1}(V)} \ar@{=}[l]
\\
& b_{small}^{-1}\mathcal{O}_Y
\ar[rd]_{b^\sharp}  & &
a_{small}^{-1}\mathcal{O}_Y \ar[ll]^t \ar[ld]^{a^\sharp} \\
& & \mathcal{O}_X
}
$$
where the triangle is commutative by definition of a $2$-isomorphism in
Modules on Sites, Section \ref{sites-modules-section-2-category}.
Above we have seen that the composition of the top horizontal
arrows comes from the identity $a^{-1}(V) = b^{-1}(V)$.
Thus the commutativity of the diagram tells us that
$a_{small}^\sharp(f) = b_{small}^\sharp(f)$ in
$\mathcal{O}_X(a^{-1}(V)) = \mathcal{O}_X(b^{-1}(V))$.
Since this holds for every open $V$ and every $f \in \mathcal{O}_Y(V)$
we conclude that $a = b$ as morphisms of schemes.
\end{proof}

\begin{lemma}
\label{lemma-morphism-ringed-etale-topoi-affines}
Let $X$, $Y$ be affine schemes.
Let
$$
(g, g^\#) :
(\Sh(X_\etale), \mathcal{O}_X)
\longrightarrow
(\Sh(Y_\etale), \mathcal{O}_Y)
$$
be a morphism of locally ringed topoi. Then there exists a
unique morphism of schemes $f : X \to Y$ such that
$(g, g^\#)$ is $2$-isomorphic to $(f_{small}, f_{small}^\sharp)$,
see
Modules on Sites,
Definition \ref{sites-modules-definition-2-morphism-ringed-topoi}.
\end{lemma}

\begin{proof}
In this proof we write $\mathcal{O}_X$ for the structure sheaf
of the small \'etale site $X_\etale$, and similarly for
$\mathcal{O}_Y$. Say $Y = \Spec(B)$ and $X = \Spec(A)$. Since
$B = \Gamma(Y_\etale, \mathcal{O}_Y)$,
$A = \Gamma(X_\etale, \mathcal{O}_X)$
we see that $g^\sharp$ induces a ring map $\varphi : B \to A$.
Let $f = \Spec(\varphi) : X \to Y$ be the corresponding morphism
of affine schemes. We will show this $f$ does the job.

\medskip\noindent
Let $V \to Y$ be an affine scheme \'etale over $Y$. Thus we may write
$V = \Spec(C)$ with $C$ an \'etale $B$-algebra. We can write
$$
C = B[x_1, \ldots, x_n]/(P_1, \ldots, P_n)
$$
with $P_i$ polynomials such that $\Delta = \det(\partial P_i/ \partial x_j)$
is invertible in $C$, see for example
Algebra, Lemma \ref{algebra-lemma-etale-standard-smooth}.
If $T$ is a scheme over $Y$, then a $T$-valued point of $V$ is given by
$n$ sections of $\Gamma(T, \mathcal{O}_T)$ which satisfy the polynomial
equations $P_1 = 0, \ldots, P_n = 0$. In other words, the sheaf $h_V$
on $Y_\etale$ is the equalizer of the two maps
$$
\xymatrix{
\prod\nolimits_{i = 1, \ldots, n} \mathcal{O}_Y
\ar@<1ex>[r]^a \ar@<-1ex>[r]_b &
\prod\nolimits_{j = 1, \ldots, n} \mathcal{O}_Y
}
$$
where $b(h_1, \ldots, h_n) = 0$ and
$a(h_1, \ldots, h_n) =
(P_1(h_1, \ldots, h_n), \ldots, P_n(h_1, \ldots, h_n))$.
Since $g^{-1}$ is exact we conclude that the top row of the
following solid commutative diagram is an equalizer diagram as well:
$$
\xymatrix{
g^{-1}h_V \ar[r] \ar@{..>}[d] &
\prod\nolimits_{i = 1, \ldots, n} g^{-1}\mathcal{O}_Y
\ar@<1ex>[r]^{g^{-1}a} \ar@<-1ex>[r]_{g^{-1}b} \ar[d]^{\prod g^\sharp} &
\prod\nolimits_{j = 1, \ldots, n} g^{-1}\mathcal{O}_Y \ar[d]^{\prod g^\sharp}\\
h_{X \times_Y V} \ar[r] &
\prod\nolimits_{i = 1, \ldots, n} \mathcal{O}_X
\ar@<1ex>[r]^{a'} \ar@<-1ex>[r]_{b'} &
\prod\nolimits_{j = 1, \ldots, n} \mathcal{O}_X  \\
}
$$
Here $b'$ is the zero map and $a'$ is the map defined by the
images $P'_i = \varphi(P_i) \in A[x_1, \ldots, x_n]$ via the same
rule
$a'(h_1, \ldots, h_n) =
(P'_1(h_1, \ldots, h_n), \ldots, P'_n(h_1, \ldots, h_n))$.
that $a$ was defined by. The commutativity of the diagram follows from
the fact that $\varphi = g^\sharp$ on global sections. The lower
row is an equalizer diagram also, by exactly the same arguments as
before since $X \times_Y V$ is the affine scheme
$\Spec(A \otimes_B C)$ and
$A \otimes_B C = A[x_1, \ldots, x_n]/(P'_1, \ldots, P'_n)$.
Thus we obtain a unique dotted arrow
$g^{-1}h_V \to h_{X \times_Y V}$ fitting into the diagram

\medskip\noindent
We claim that the map of sheaves $g^{-1}h_V \to h_{X \times_Y V}$
is an isomorphism. Since the small \'etale site of $X$ has enough points
(Theorem \ref{theorem-exactness-stalks})
it suffices to prove this on stalks. Hence let $\overline{x}$ be a
geometric point of $X$, and denote $p$ the associate point of the
small \'etale topos of $X$. Set $q = g \circ p$. This is a point of
the small \'etale topos of $Y$. By
Lemma \ref{lemma-points-small-etale-site}
we see that $q$ corresponds to a geometric point $\overline{y}$ of
$Y$. Consider the map of stalks
$$
(g^\sharp)_p :
(\mathcal{O}_Y)_{\overline{y}} =
\mathcal{O}_{Y, q} =
(g^{-1}\mathcal{O}_Y)_p
\longrightarrow
\mathcal{O}_{X, p} =
(\mathcal{O}_X)_{\overline{x}}
$$
Since $(g, g^\sharp)$ is a morphism of {\it locally} ringed topoi
$(g^\sharp)_p$ is a local ring homomorphism of strictly henselian
local rings. Applying localization to the big commutative diagram above and
Algebra, Lemma \ref{algebra-lemma-strictly-henselian-solutions}
we conclude that $(g^{-1}h_V)_p \to (h_{X \times_Y V})_p$ is an isomorphism
as desired.

\medskip\noindent
We claim that the isomorphisms $g^{-1}h_V \to h_{X \times_Y V}$ are
functorial. Namely, suppose that $V_1 \to V_2$ is a morphism of affine
schemes \'etale over $Y$. Write
$V_i = \Spec(C_i)$ with
$$
C_i = B[x_{i, 1}, \ldots, x_{i, n_i}]/(P_{i, 1}, \ldots, P_{i, n_i})
$$
The morphism $V_1 \to V_2$ is given by a $B$-algebra map $C_2 \to C_1$
which in turn is given by some polynomials
$Q_j \in B[x_{1, 1}, \ldots, x_{1, n_1}]$ for $j = 1, \ldots, n_2$.
Then it is an easy matter to show that the diagram of sheaves
$$
\xymatrix{
h_{V_1} \ar[d] \ar[r] & \prod_{i = 1, \ldots, n_1} \mathcal{O}_Y
\ar[d]^{Q_1, \ldots, Q_{n_2}}\\
h_{V_2} \ar[r] & \prod_{i = 1, \ldots, n_2} \mathcal{O}_Y
}
$$
is commutative, and pulling back to $X_\etale$ we obtain the
solid commutative diagram
$$
\xymatrix{
g^{-1}h_{V_1} \ar@{..>}[dd] \ar[rrd] \ar[r] &
\prod_{i = 1, \ldots, n_1} g^{-1}\mathcal{O}_Y
\ar[dd]^{g^\sharp}
\ar[rrd]^{Q_1, \ldots, Q_{n_2}} \\
& & g^{-1}h_{V_2} \ar@{..>}[dd] \ar[r] &
\prod_{i = 1, \ldots, n_2} g^{-1}\mathcal{O}_Y
\ar[dd]^{g^\sharp} \\
h_{X \times_Y V_1} \ar[r] \ar[rrd] &
\prod\nolimits_{i = 1, \ldots, n_1} \mathcal{O}_X
\ar[rrd]^{Q'_1, \ldots, Q'_{n_2}} \\
& & h_{X \times_Y V_2} \ar[r] &
\prod\nolimits_{i = 1, \ldots, n_2} \mathcal{O}_X
}
$$
where $Q'_j \in A[x_{1, 1}, \ldots, x_{1, n_1}]$ is the image of
$Q_j$ via $\varphi$. Since the dotted arrows exist, make the
two squares commute, and the horizontal arrows are injective
we see that the whole diagram commutes. This proves functoriality
(and also that the construction of $g^{-1}h_V \to h_{X \times_Y V}$
is independent of the choice of the presentation, although we
strictly speaking do not need to show this).

\medskip\noindent
At this point we are able to show that $f_{small, *} \cong g_*$.
Namely, let $\mathcal{F}$ be a sheaf on $X_\etale$. For every
$V \in \Ob(X_\etale)$ affine we have
\begin{align*}
(g_*\mathcal{F})(V)
& =
\Mor_{\Sh(Y_\etale)}(h_V, g_*\mathcal{F}) \\
& =
\Mor_{\Sh(X_\etale)}(g^{-1}h_V, \mathcal{F}) \\
& =
\Mor_{\Sh(X_\etale)}(h_{X \times_Y V}, \mathcal{F}) \\
& =
\mathcal{F}(X \times_Y V) \\
& =
f_{small, *}\mathcal{F}(V)
\end{align*}
where in the third equality we use the isomorphism
$g^{-1}h_V \cong h_{X \times_Y V}$ constructed above. These isomorphisms
are clearly functorial in $\mathcal{F}$ and functorial in $V$
as the isomorphisms $g^{-1}h_V \cong h_{X \times_Y V}$ are functorial.
Now any sheaf on $Y_\etale$ is determined by the restriction
to the subcategory of affine schemes
(Topologies, Lemma \ref{topologies-lemma-alternative}),
and hence we obtain an isomorphism of functors $f_{small, *} \cong g_*$
as desired.

\medskip\noindent
Finally, we have to check that, via the isomorphism
$f_{small, *} \cong g_*$ above, the maps $f_{small}^\sharp$ and
$g^\sharp$ agree. By construction this is already the case for the
global sections of $\mathcal{O}_Y$, i.e., for the elements of $B$.
We only need to check the result on
sections over an affine $V$ \'etale over $Y$ (by
Topologies, Lemma \ref{topologies-lemma-alternative}
again). Writing
$V = \Spec(C)$, $C = B[x_i]/(P_j)$ as before it suffices
to check that the coordinate functions $x_i$ are mapped to
the same sections of $\mathcal{O}_X$ over $X \times_Y V$.
And this is exactly what it means that the diagram
$$
\xymatrix{
g^{-1}h_V \ar[r] \ar@{..>}[d] &
\prod\nolimits_{i = 1, \ldots, n} g^{-1}\mathcal{O}_Y
\ar[d]^{\prod g^\sharp} \\
h_{X \times_Y V} \ar[r] &
\prod\nolimits_{i = 1, \ldots, n} \mathcal{O}_X
}
$$
commutes. Thus the lemma is proved.
\end{proof}

\noindent
Here is a version for general schemes.

\begin{theorem}
\label{theorem-fully-faithful}
Let $X$, $Y$ be schemes. Let
$$
(g, g^\#) :
(\Sh(X_\etale), \mathcal{O}_X)
\longrightarrow
(\Sh(Y_\etale), \mathcal{O}_Y)
$$
be a morphism of locally ringed topoi. Then there exists a
unique morphism of schemes $f : X \to Y$ such that
$(g, g^\#)$ is isomorphic to $(f_{small}, f_{small}^\sharp)$.
In other words, the construction
$$
\Sch \longrightarrow \textit{Locally ringed topoi},
\quad
X \longrightarrow (X_\etale, \mathcal{O}_X)
$$
is fully faithful (morphisms up to $2$-isomorphisms on the right hand side).
\end{theorem}

\begin{proof}
You can prove this theorem by carefully adjusting the arguments of
the proof of
Lemma \ref{lemma-morphism-ringed-etale-topoi-affines}
to the global setting. However, we want to indicate how we
can glue the result of that lemma to get a global morphism
due to the rigidity provided by the result of
Lemma \ref{lemma-2-morphism}.
Unfortunately, this is a bit messy.

\medskip\noindent
Let us prove existence when $Y$ is affine. In this case choose an
affine open covering $X = \bigcup U_i$. For each $i$ the inclusion
morphism $j_i : U_i \to X$ induces a morphism of locally ringed topoi
$(j_{i, small}, j_{i, small}^\sharp) :
(\Sh(U_{i, \etale}), \mathcal{O}_{U_i})
\to
(\Sh(X_\etale), \mathcal{O}_X)$
by
Lemma \ref{lemma-morphism-locally-ringed}.
We can compose this with $(g, g^\sharp)$ to obtain a morphism
of locally ringed topoi
$$
(g, g^\sharp) \circ (j_{i, small}, j_{i, small}^\sharp) :
(\Sh(U_{i, \etale}), \mathcal{O}_{U_i})
\to
(\Sh(Y_\etale), \mathcal{O}_Y)
$$
see
Modules on Sites,
Lemma \ref{sites-modules-lemma-composition-morphisms-locally-ringed-topoi}.
By
Lemma \ref{lemma-morphism-ringed-etale-topoi-affines}
there exists a unique morphism of schemes $f_i : U_i \to Y$
and a $2$-isomorphism
$$
t_i :
(f_{i, small}, f_{i, small}^\sharp)
\longrightarrow
(g, g^\sharp) \circ (j_{i, small}, j_{i, small}^\sharp).
$$
Set $U_{i, i'} = U_i \cap U_{i'}$, and denote $j_{i, i'} : U_{i, i'} \to U_i$
the inclusion morphism. Since we have
$j_i \circ j_{i, i'} = j_{i'} \circ j_{i', i}$
we see that
\begin{align*}
(g, g^\sharp) \circ
(j_{i, small}, j_{i, small}^\sharp) \circ
(j_{i, i', small}, j_{i, i', small}^\sharp)
= \\
(g, g^\sharp) \circ
(j_{i', small}, j_{i', small}^\sharp) \circ
(j_{i', i, small}, j_{i', i, small}^\sharp)
\end{align*}
Hence by uniqueness (see
Lemma \ref{lemma-faithful})
we conclude that
$f_i \circ j_{i, i'} = f_{i'} \circ j_{i', i}$, in other words the
morphisms of schemes $f_i = f \circ j_i$ are the restrictions of a
global morphism of schemes $f : X \to Y$. Consider the diagram
of $2$-isomorphisms (where we drop the components ${}^\sharp$ to ease the
notation)
$$
\xymatrix{
g \circ j_{i, small} \circ j_{i, i', small}
\ar[rr]^{t_i \star \text{id}_{j_{i, i', small}}}
\ar@{=}[d] & &
f_{small} \circ j_{i, small} \circ j_{i, i', small} \ar@{=}[d] \\
g \circ j_{i', small} \circ j_{i', i, small}
\ar[rr]^{t_{i'} \star \text{id}_{j_{i', i, small}}} & &
f_{small} \circ j_{i', small} \circ j_{i', i, small}
}
$$
The notation $\star$ indicates horizontal composition, see
Categories, Definition \ref{categories-definition-2-category}
in general and
Sites, Section \ref{sites-section-2-category}
for our particular case. By the result of
Lemma \ref{lemma-2-morphism}
this diagram commutes. Hence for any sheaf $\mathcal{G}$
on $Y_\etale$ the isomorphisms
$t_i : f_{small}^{-1}\mathcal{G}|_{U_i} \to g^{-1}\mathcal{G}|_{U_i}$
agree over $U_{i, i'}$ and we obtain a global isomorphism
$t : f_{small}^{-1}\mathcal{G} \to g^{-1}\mathcal{G}$.
It is clear that this isomorphism is functorial in $\mathcal{G}$
and is compatible with the maps $f_{small}^\sharp$ and $g^\sharp$
(because it is compatible with these maps locally).
This proves the theorem in case $Y$ is affine.

\medskip\noindent
In the general case, let $V \subset Y$ be an affine open.
Then $h_V$ is a subsheaf of the final sheaf $*$ on $Y_\etale$.
As $g$ is exact we see that $g^{-1}h_V$ is a subsheaf of the final
sheaf on $X_\etale$. Hence by
Lemma \ref{lemma-support-subsheaf-final}
there exists an open subscheme $W \subset X$ such that $g^{-1}h_V = h_W$. By
Modules on Sites,
Lemma \ref{sites-modules-lemma-localize-morphism-locally-ringed-topoi}
there exists a commutative diagram of morphisms of locally ringed
topoi
$$
\xymatrix{
(\Sh(W_\etale), \mathcal{O}_W) \ar[r] \ar[d]_{g'} &
(\Sh(X_\etale), \mathcal{O}_X) \ar[d]^g \\
(\Sh(V_\etale), \mathcal{O}_V) \ar[r] &
(\Sh(Y_\etale), \mathcal{O}_Y)
}
$$
where the horizontal arrows are the localization morphisms
(induced by the inclusion morphisms $V \to Y$ and $W \to X$)
and where $g'$ is induced from $g$. By the result of the preceding
paragraph we obtain a morphism of schemes $f' : W \to V$ and
a $2$-isomorphism
$t : (f'_{small}, (f'_{small})^\sharp) \to (g', (g')^\sharp)$.
Exactly as before these morphisms $f'$ (for varying affine opens $V \subset Y$)
agree on overlaps by uniqueness, so we get a morphism $f : X \to Y$.
Moreover, the $2$-isomorphisms $t$ are compatible on overlaps by
Lemma \ref{lemma-2-morphism}
again and we obtain a global $2$-isomorphism
$(f_{small}, (f_{small})^\sharp) \to (g, (g)^\sharp)$.
as desired. Some details omitted.
\end{proof}










\section{Push and pull}
\label{section-monomorphisms}

\noindent
Let $f : X \to Y$ be a morphism of schemes.
Here is a list of conditions we will consider in the following:
\begin{enumerate}
\item[(A)] For every \'etale morphism $U \to X$ and $u \in U$ there exist
an \'etale morphism $V \to Y$ and a disjoint union decomposition
$X \times_Y V = W \amalg W'$ and a morphism $h : W \to U$ over $X$
with $u$ in the image of $h$.
\item[(B)] For every $V \to Y$ \'etale, and every \'etale covering
$\{U_i \to X \times_Y V\}$ there exists an \'etale covering
$\{V_j \to V\}$ such that for each $j$ we have
$X \times_Y V_j = \coprod W_{ij}$ where $W_{ij} \to X \times_Y V$
factors through $U_i \to X \times_Y V$ for some $i$.
\item[(C)] For every $U \to X$ \'etale, there exists a $V \to Y$ \'etale
and a surjective morphism $X \times_Y V \to U$ over $X$.
\end{enumerate}
It turns out that each of these properties has meaning in terms of
the behaviour of the functor $f_{small, *}$. We will work this
out in the next few sections.



\section{Property (A)}
\label{section-A}

\noindent
Please see Section \ref{section-monomorphisms} for the definition of property
(A).

\begin{lemma}
\label{lemma-property-A-implies}
Let $f : X \to Y$ be a morphism of schemes.
Assume (A).
\begin{enumerate}
\item
$f_{small, *} :
\textit{Ab}(X_\etale)
\to
\textit{Ab}(Y_\etale)$
reflects injections and surjections,
\item $f_{small}^{-1}f_{small, *}\mathcal{F} \to \mathcal{F}$
is surjective for any abelian sheaf $\mathcal{F}$ on $X_\etale$,
\item
$f_{small, *} :
\textit{Ab}(X_\etale)
\to
\textit{Ab}(Y_\etale)$
is faithful.
\end{enumerate}
\end{lemma}

\begin{proof}
Let $\mathcal{F}$ be an abelian sheaf on $X_\etale$.
Let $U$ be an object of $X_\etale$. By assumption we can find a
covering $\{W_i \to U\}$ in $X_\etale$ such that each $W_i$ is
an open and closed subscheme of $X \times_Y V_i$ for some object
$V_i$ of $Y_\etale$. The sheaf condition shows that
$$
\mathcal{F}(U) \subset \prod \mathcal{F}(W_i)
$$
and that $\mathcal{F}(W_i)$ is a direct summand of
$\mathcal{F}(X \times_Y V_i) = f_{small, *}\mathcal{F}(V_i)$.
Hence it is clear that $f_{small, *}$ reflects injections.

\medskip\noindent
Next, suppose that $a : \mathcal{G} \to \mathcal{F}$ is a map of
abelian sheaves such that $f_{small, *}a$ is surjective. Let
$s \in \mathcal{F}(U)$ with $U$ as above. With $W_i$, $V_i$ as
above we see that it suffices to show that $s|_{W_i}$ is \'etale locally
the image of a section of $\mathcal{G}$ under $a$. Since $\mathcal{F}(W_i)$
is a direct summand of $\mathcal{F}(X \times_Y V_i)$
it suffices to show that for any $V \in \Ob(Y_\etale)$
any element $s \in \mathcal{F}(X \times_Y V)$
is \'etale locally on $X \times_Y V$ the image of a section of
$\mathcal{G}$ under $a$. Since
$\mathcal{F}(X \times_Y V) = f_{small, *}\mathcal{F}(V)$
we see by assumption that there exists a covering $\{V_j \to V\}$ such that
$s$ is the image of
$s_j \in f_{small, *}\mathcal{G}(V_j) = \mathcal{G}(X \times_Y V_j)$.
This proves $f_{small, *}$ reflects surjections.

\medskip\noindent
Parts (2), (3) follow formally from part (1), see
Modules on Sites, Lemma \ref{sites-modules-lemma-reflect-surjections}.
\end{proof}

\begin{lemma}
\label{lemma-locally-quasi-finite-A}
Let $f : X \to Y$ be a separated locally quasi-finite morphism of schemes.
Then property (A) above holds.
\end{lemma}

\begin{proof}
Let $U \to X$ be an \'etale morphism and $u \in U$.
The geometric statement (A) reduces directly to the case where $U$ and $Y$
are affine schemes. Denote $x \in X$ and $y \in Y$ the
images of $u$. Since $X \to Y$ is locally quasi-finite, and $U \to X$ is
locally quasi-finite (see
Morphisms, Lemma \ref{morphisms-lemma-etale-locally-quasi-finite})
we see that $U \to Y$ is locally quasi-finite (see
Morphisms, Lemma \ref{morphisms-lemma-composition-quasi-finite}).
Moreover both $X \to Y$ and $U \to Y$ are separated. Thus
More on Morphisms, Lemma
\ref{more-morphisms-lemma-etale-splits-off-quasi-finite-part-technical-variant}
applies to both morphisms. This means we may pick an \'etale neighbourhood
$(V, v) \to (Y, y)$ such that
$$
X \times_Y V = W \amalg R, \quad
U \times_Y V = W' \amalg R'
$$
and points $w \in W$, $w' \in W'$ such that
\begin{enumerate}
\item $W$, $R$ are open and closed in $X \times_Y V$,
\item $W'$, $R'$ are open and closed in $U \times_Y V$,
\item $W \to V$ and $W' \to V$ are finite,
\item $w$, $w'$ map to $v$,
\item $\kappa(v) \subset \kappa(w)$ and $\kappa(v) \subset \kappa(w')$
are purely inseparable, and
\item no other point of $W$ or $W'$ maps to $v$.
\end{enumerate}
Here is a commutative diagram
$$
\xymatrix{
U \ar[d] & U \times_Y V \ar[l] \ar[d] & W' \amalg R' \ar[d] \ar[l] \\
X \ar[d] & X \times_Y V \ar[l] \ar[d] & W \amalg R \ar[l] \\
Y & V \ar[l]
}
$$
After shrinking $V$ we may assume that $W'$ maps into $W$:
just remove the image the inverse image of $R$ in $W'$; this is
a closed set (as $W' \to V$ is finite) not containing $v$.
Then $W' \to W$ is finite because both $W \to V$ and $W' \to V$ are finite.
Hence $W' \to W$ is finite \'etale, and there is exactly one point in the
fibre over $w$ with $\kappa(w) = \kappa(w')$. Hence $W' \to W$ is an
isomorphism in an open neighbourhood $W^\circ$ of $w$, see
\'Etale Morphisms, Lemma \ref{etale-lemma-finite-etale-one-point}.
Since $W \to V$ is finite the image of $W \setminus W^\circ$ is a closed
subset $T$ of $V$ not containing $v$. Thus after replacing $V$ by
$V \setminus T$ we may assume that $W' \to W$ is an isomorphism.
Now the decomposition $X \times_Y V = W \amalg R$ and the morphism
$W \to U$ are as desired and we win.
\end{proof}

\begin{lemma}
\label{lemma-integral-A}
Let $f : X \to Y$ be an integral morphism of schemes.
Then property (A) holds.
\end{lemma}

\begin{proof}
Let $U \to X$ be \'etale, and let $u \in U$ be a point.
We have to find $V \to Y$ \'etale, a disjoint union decomposition
$X \times_Y V = W \amalg W'$ and an $X$-morphism $W \to U$
with $u$ in the image. We may shrink $U$ and $Y$ and assume
$U$ and $Y$ are affine. In this case also $X$ is affine, since
an integral morphism is affine by definition. Write $Y = \Spec(A)$,
$X = \Spec(B)$ and $U = \Spec(C)$. Then $A \to B$ is an
integral ring map, and $B \to C$ is an \'etale ring map. By
Algebra, Lemma \ref{algebra-lemma-etale}
we can find a finite $A$-subalgebra $B' \subset B$ and an \'etale ring
map $B' \to C'$ such that $C = B \otimes_{B'} C'$. Thus the question
reduces to the \'etale morphism
$U' = \Spec(C') \to X' = \Spec(B')$
over the finite morphism $X' \to Y$. In this case the result follows from
Lemma \ref{lemma-locally-quasi-finite-A}.
\end{proof}

\begin{lemma}
\label{lemma-when-push-pull-surjective}
Let $f : X \to Y$ be a morphism of schemes. Denote
$f_{small} :
\Sh(X_\etale)
\to
\Sh(Y_\etale)$
the associated morphism of small \'etale topoi. Assume at least one
of the following
\begin{enumerate}
\item $f$ is integral, or
\item $f$ is separated and locally quasi-finite.
\end{enumerate}
Then the functor
$f_{small, *} :
\textit{Ab}(X_\etale)
\to
\textit{Ab}(Y_\etale)$
has the following properties
\begin{enumerate}
\item the map
$f_{small}^{-1}f_{small, *}\mathcal{F} \to \mathcal{F}$
is always surjective,
\item $f_{small, *}$ is faithful, and
\item $f_{small, *}$ reflects injections and surjections.
\end{enumerate}
\end{lemma}

\begin{proof}
Combine
Lemmas \ref{lemma-locally-quasi-finite-A},
\ref{lemma-integral-A}, and
\ref{lemma-property-A-implies}.
\end{proof}



\section{Property (B)}
\label{section-B}

\noindent
Please see Section \ref{section-monomorphisms} for the definition of property
(B).

\begin{lemma}
\label{lemma-property-B-implies}
Let $f : X \to Y$ be a morphism of schemes. Assume (B) holds.
Then the functor
$f_{small, *} :
\Sh(X_\etale)
\to
\Sh(Y_\etale)$
transforms surjections into surjections.
\end{lemma}

\begin{proof}
This follows from
Sites, Lemma \ref{sites-lemma-weaker}.
\end{proof}

\begin{lemma}
\label{lemma-simplify-B}
Let $f : X \to Y$ be a morphism of schemes. Suppose
\begin{enumerate}
\item $V \to Y$ is an \'etale morphism of schemes,
\item $\{U_i \to X \times_Y V\}$ is an \'etale covering, and
\item $v \in V$ is a point.
\end{enumerate}
Assume that for any such data there exists an \'etale neighbourhood
$(V', v') \to (V, v)$, a disjoint union decomposition
$X \times_Y V' = \coprod W'_i$, and morphisms $W'_i \to U_i$
over $X \times_Y V$. Then property (B) holds.
\end{lemma}

\begin{proof}
Omitted.
\end{proof}

\begin{lemma}
\label{lemma-finite-B}
Let $f : X \to Y$ be a finite morphism of schemes.
Then property (B) holds.
\end{lemma}

\begin{proof}
Consider $V \to Y$ \'etale, $\{U_i \to X \times_Y V\}$ an \'etale covering, and
$v \in V$. We have to find a $V' \to V$ and decomposition and maps as in
Lemma \ref{lemma-simplify-B}.
We may shrink $V$ and $Y$, hence we may assume that $V$ and $Y$ are affine.
Since $X$ is finite over $Y$, this also implies that $X$ is affine.
During the proof we may (finitely often) replace $(V, v)$ by an
\'etale neighbourhood $(V', v')$ and correspondingly the covering
$\{U_i \to X \times_Y V\}$ by $\{V' \times_V U_i \to X \times_Y V'\}$.

\medskip\noindent
Since $X \times_Y V \to V$ is finite there exist finitely
many (pairwise distinct) points $x_1, \ldots, x_n \in X \times_Y V$
mapping to $v$. We may apply
More on Morphisms, Lemma
\ref{more-morphisms-lemma-etale-splits-off-quasi-finite-part-technical-variant}
to $X \times_Y V \to V$ and the points $x_1, \ldots, x_n$ lying over
$v$ and find an \'etale neighbourhood $(V', v') \to (V, v)$
such that
$$
X \times_Y V' = R \amalg \coprod T_a
$$
with $T_a \to V'$ finite with exactly one point $p_a$ lying over $v'$
and moreover $\kappa(v') \subset \kappa(p_a)$ purely inseparable, and
such that $R \to V'$ has empty fibre over $v'$.
Because $X \to Y$ is finite, also $R \to V'$ is finite. Hence after
shrinking $V'$ we may assume that $R = \emptyset$. Thus we may
assume that $X \times_Y V = X_1 \amalg \ldots \amalg X_n$ with
exactly one point $x_l \in X_l$ lying over $v$ with moreover
$\kappa(v) \subset \kappa(x_l)$ purely inseparable. Note that this
property is preserved under refinement of the \'etale neighbourhood
$(V, v)$.

\medskip\noindent
For each $l$ choose an $i_l$ and a point $u_l \in U_{i_l}$ mapping to $x_l$.
Now we apply property (A) for the finite morphism
$X \times_Y V \to V$ and the \'etale
morphisms $U_{i_l} \to X \times_Y V$ and the points $u_l$.
This is permissible by
Lemma \ref{lemma-integral-A}
This gives produces an \'etale neighbourhood $(V', v') \to (V, v)$
and decompositions
$$
X \times_Y V' = W_l \amalg R_l
$$
and $X$-morphisms $a_l : W_l \to U_{i_l}$ whose image contains $u_{i_l}$.
Here is a picture:
$$
\xymatrix{
& & & U_{i_l} \ar[d] & \\
W_l \ar[rrru] \ar[r] & W_l \amalg R_l \ar@{=}[r] &
X \times_Y V' \ar[r] \ar[d] &
X \times_Y V \ar[r] \ar[d] & X \ar[d] \\
& & V' \ar[r] & V \ar[r] & Y
}
$$
After replacing $(V, v)$ by $(V', v')$ we conclude that each
$x_l$ is contained in an open and closed neighbourhood $W_l$ such that
the inclusion morphism $W_l \to X \times_Y V$ factors through
$U_i \to X \times_Y V$ for some $i$. Replacing $W_l$ by $W_l \cap X_l$
we see that these open and closed sets are disjoint and moreover
that $\{x_1, \ldots, x_n\} \subset W_1 \cup \ldots \cup W_n$.
Since $X \times_Y V \to V$ is finite we may shrink $V$ and assume that
$X \times_Y V = W_1 \amalg \ldots \amalg W_n$ as desired.
\end{proof}

\begin{lemma}
\label{lemma-integral-B}
Let $f : X \to Y$ be an integral morphism of schemes.
Then property (B) holds.
\end{lemma}

\begin{proof}
Consider $V \to Y$ \'etale, $\{U_i \to X \times_Y V\}$ an \'etale covering, and
$v \in V$. We have to find a $V' \to V$ and decomposition and maps as in
Lemma \ref{lemma-simplify-B}.
We may shrink $V$ and $Y$, hence we may assume that $V$ and $Y$ are affine.
Since $X$ is integral over $Y$, this also implies that $X$ and
$X \times_Y V$ are affine. We may refine the covering
$\{U_i \to X \times_Y V\}$, and hence we may assume that
$\{U_i \to X \times_Y V\}_{i = 1, \ldots, n}$ is a standard \'etale covering.
Write $Y = \Spec(A)$, $X = \Spec(B)$,
$V = \Spec(C)$, and $U_i = \Spec(B_i)$.
Then $A \to B$ is an integral ring map, and $B \otimes_A C \to B_i$ are
\'etale ring maps. By
Algebra, Lemma \ref{algebra-lemma-etale}
we can find a finite $A$-subalgebra $B' \subset B$ and an \'etale ring
map $B' \otimes_A C \to B'_i$ for $i = 1, \ldots, n$
such that $B_i = B \otimes_{B'} B'_i$. Thus the question
reduces to the \'etale covering
$\{\Spec(B'_i) \to X' \times_Y V\}_{i = 1, \ldots, n}$
with $X' = \Spec(B')$ finite over $Y$.
In this case the result follows from
Lemma \ref{lemma-finite-B}.
\end{proof}

\begin{lemma}
\label{lemma-what-integral}
Let $f : X \to Y$ be a morphism of schemes.
Assume $f$ is integral (for example finite).
Then
\begin{enumerate}
\item $f_{small, *}$ transforms surjections into surjections (on sheaves
of sets and on abelian sheaves),
\item $f_{small}^{-1}f_{small, *}\mathcal{F} \to \mathcal{F}$
is surjective for any abelian sheaf $\mathcal{F}$ on $X_\etale$,
\item
$f_{small, *} :
\textit{Ab}(X_\etale)
\to
\textit{Ab}(Y_\etale)$
is faithful and reflects injections and surjections, and
\item
$f_{small, *} :
\textit{Ab}(X_\etale)
\to
\textit{Ab}(Y_\etale)$
is exact.
\end{enumerate}
\end{lemma}

\begin{proof}
Parts (2), (3) we have seen in
Lemma \ref{lemma-when-push-pull-surjective}.
Part (1) follows from
Lemmas \ref{lemma-integral-B} and \ref{lemma-property-B-implies}.
Part (4) is a consequence of part (1), see
Modules on Sites, Lemma \ref{sites-modules-lemma-exactness}.
\end{proof}








\section{Property (C)}
\label{section-C}

\noindent
Please see Section \ref{section-monomorphisms} for the definition of property
(C).

\begin{lemma}
\label{lemma-property-C-implies}
Let $f : X \to Y$ be a morphism of schemes. Assume (C) holds. Then the functor
$f_{small, *} :
\Sh(X_\etale)
\to
\Sh(Y_\etale)$
reflects injections and surjections.
\end{lemma}

\begin{proof}
Follows from
Sites, Lemma \ref{sites-lemma-cover-from-below}.
We omit the verification that property (C) implies that the functor
$Y_\etale \to X_\etale$, $V \mapsto X \times_Y V$
satisfies the assumption of
Sites, Lemma \ref{sites-lemma-cover-from-below}.
\end{proof}

\begin{remark}
\label{remark-property-C-strong}
Property (C) holds if $f : X \to Y$ is an open immersion. Namely, if
$U \in \Ob(X_\etale)$, then we can view $U$ also as an object
of $Y_\etale$ and $U \times_Y X = U$. Hence property (C)
does not imply that $f_{small, *}$ is exact as this is not
the case for open immersions (in general).
\end{remark}

\begin{lemma}
\label{lemma-property-C-closed-implies}
Let $f : X \to Y$ be a morphism of schemes. Assume that
for any $V \to Y$ \'etale we have that
\begin{enumerate}
\item $X \times_Y V \to V$ has property (C), and
\item $X \times_Y V \to V$ is closed.
\end{enumerate}
Then the functor
$Y_\etale \to X_\etale$, $V \mapsto X \times_Y V$
is almost cocontinuous, see
Sites, Definition \ref{sites-definition-almost-cocontinuous}.
\end{lemma}

\begin{proof}
Let $V \to Y$ be an object of $Y_\etale$ and let
$\{U_i \to X \times_Y V\}_{i \in I}$ be a covering of $X_\etale$.
By assumption (1) for each $i$ we can find an \'etale morphism
$h_i : V_i \to V$ and a surjective morphism $X \times_Y V_i \to U_i$
over $X \times_Y V$. Note that $\bigcup h_i(V_i) \subset V$ is an
open set containing the closed set $Z = \Im(X \times_Y V \to V)$.
Let $h_0 : V_0 = V \setminus Z \to V$ be the open immersion.
It is clear that $\{V_i \to V\}_{i \in I \cup \{0\}}$ is an
\'etale covering such that for each $i \in I \cup \{0\}$ we have
either $V_i \times_Y X = \emptyset$ (namely if $i = 0$), or
$V_i \times_Y X \to V \times_Y X$ factors through $U_i \to X \times_Y V$
(if $i \not = 0$). Hence the functor $Y_\etale \to X_\etale$
is almost cocontinuous.
\end{proof}

\begin{lemma}
\label{lemma-integral-homeo-onto-image-C}
Let $f : X \to Y$ be an integral morphism of schemes which defines
a homeomorphism of $X$ with a closed subset of $Y$.
Then property (C) holds.
\end{lemma}

\begin{proof}
Let $g : U \to X$ be an \'etale morphism. We need to find an object
$V \to Y$ of $Y_\etale$ and a surjective morphism $X \times_Y V \to U$
over $X$. Suppose that for every $u \in U$ we can find an object
$V_u \to Y$ of $Y_\etale$ and a morphism $h_u : X \times_Y V_u \to U$
over $X$ with $u \in \Im(h_u)$. Then we can take $V = \coprod V_u$
and $h = \coprod h_u$ and we win. Hence given a point
$u \in U$ we find a pair $(V_u, h_u)$ as above. To do this we may
shrink $U$ and assume that $U$ is affine. In this case
$g : U \to X$ is locally quasi-finite. Let
$g^{-1}(g(\{u\})) = \{u, u_2, \ldots, u_n\}$. Since there are no
specializations $u_i \leadsto u$ we may replace $U$ by an affine neighbourhood
so that $g^{-1}(g(\{u\})) = \{u\}$.

\medskip\noindent
The image $g(U) \subset X$ is open,
hence $f(g(U))$ is locally closed in $Y$. Choose an open $V \subset Y$ such
that $f(g(U)) = f(X) \cap V$. It follows that $g$ factors through
$X \times_Y V$ and that the resulting $\{U \to X \times_Y V\}$ is an \'etale
covering. Since $f$ has property (B) , see
Lemma \ref{lemma-integral-B},
we see that there exists an \'etale covering $\{V_j \to V\}$ such that
$X \times_Y V_j \to X \times_Y V$ factor through $U$.
This implies that $V' = \coprod V_j$ is \'etale over $Y$ and that there is a
morphism $h : X \times_Y V' \to U$ whose image
surjects onto $g(U)$. Since $u$ is the only point in its fibre it must
be in the image of $h$ and we win.
\end{proof}

\noindent
We urge the reader to think of the following lemma as a
way station\footnote{A way station is a place where people stop to eat
and rest when they are on a long journey.} on the journey towards the
ultimate truth regarding $f_{small, *}$ for integral universally injective
morphisms.

\begin{lemma}
\label{lemma-integral-universally-injective}
Let $f : X \to Y$ be a morphism of schemes. Assume that $f$ is
universally injective and integral (for example a closed immersion).
Then
\begin{enumerate}
\item
$f_{small, *} :
\Sh(X_\etale)
\to
\Sh(Y_\etale)$
reflects injections and surjections,
\item
$f_{small, *} :
\Sh(X_\etale)
\to
\Sh(Y_\etale)$
commutes with pushouts and coequalizers (and more generally
finite connected colimits),
\item $f_{small, *}$ transforms surjections into surjections (on sheaves
of sets and on abelian sheaves),
\item the map
$f_{small}^{-1}f_{small, *}\mathcal{F} \to \mathcal{F}$
is surjective for any sheaf (of sets or of abelian groups)
$\mathcal{F}$ on $X_\etale$,
\item the functor $f_{small, *}$ is faithful (on sheaves of sets and
on abelian sheaves),
\item
$f_{small, *} :
\textit{Ab}(X_\etale)
\to
\textit{Ab}(Y_\etale)$
is exact, and
\item the functor
$Y_\etale \to X_\etale$, $V \mapsto X \times_Y V$ is
almost cocontinuous.
\end{enumerate}
\end{lemma}

\begin{proof}
By
Lemmas \ref{lemma-integral-A},
\ref{lemma-integral-B} and
\ref{lemma-integral-homeo-onto-image-C}
we know that the morphism $f$ has properties (A), (B), and (C).
Moreover, by
Lemma \ref{lemma-property-C-closed-implies}
we know that the functor $Y_\etale \to X_\etale$ is
almost cocontinuous. Now we have
\begin{enumerate}
\item property (C) implies (1) by
Lemma \ref{lemma-property-C-implies},
\item almost continuous implies (2) by
Sites, Lemma \ref{sites-lemma-morphism-of-sites-almost-cocontinuous},
\item property (B) implies (3) by
Lemma \ref{lemma-property-B-implies}.
\end{enumerate}
Properties (4), (5), and (6) follow formally from the first three, see
Sites, Lemma \ref{sites-lemma-exactness-properties}
and
Modules on Sites, Lemma \ref{sites-modules-lemma-exactness}.
Property (7) we saw above.
\end{proof}




\section{Topological invariance of the small \'etale site}
\label{section-topological-invariance}

\noindent
In the following theorem we show that the small \'etale site is a topological
invariant in the following sense: If $f : X \to Y$ is a morphism of schemes
which is a universal homeomorphism, then $X_\etale \cong Y_\etale$
as sites. This improves the result of
\'Etale Morphisms, Theorem \ref{etale-theorem-remarkable-equivalence}.
We first prove the result for morphisms and then we state the result
for categories.

\begin{theorem}
\label{theorem-etale-topological}
Let $X$ and $Y$ be two schemes over a base scheme $S$. Let
$S' \to S$ be a universal homeomorphism.
Denote $X'$ (resp.\ $Y'$) the base change to $S'$.
If $X$ is \'etale over $S$, then the map
$$
\Mor_S(Y, X) \longrightarrow \Mor_{S'}(Y', X')
$$
is bijective.
\end{theorem}

\begin{proof}
After base changing via $Y \to S$, we may assume that $Y = S$.
Thus we may and do assume both $X$ and $Y$ are \'etale over $S$.
In other words, the theorem states that the base change functor
is a fully faithful functor from the category of schemes \'etale
over $S$ to the category of schemes \'etale over $S'$.

\medskip\noindent
Consider the forgetful functor
\begin{equation}
\label{equation-descent-etale-forget}
\begin{matrix}
\text{descent data }(X', \varphi')\text{ relative to }S'/S \\
\text{ with }X'\text{ \'etale over }S'
\end{matrix}
\longrightarrow
\text{schemes }X'\text{ \'etale over }S'
\end{equation}
We claim this functor is an equivalence. On the other hand, the
functor
\begin{equation}
\label{equation-descent-etale}
\text{schemes }X\text{ \'etale over }S \longrightarrow
\begin{matrix}
\text{descent data }(X', \varphi')\text{ relative to }S'/S \\
\text{ with }X'\text{ \'etale over }S'
\end{matrix}
\end{equation}
is fully faithful by \'Etale Morphisms, Lemma
\ref{etale-lemma-fully-faithful-cases}.
Thus the claim implies the theorem.

\medskip\noindent
Proof of the claim.
Recall that a universal homeomorphism is the same thing as an
integral, universally injective, surjective morphism, see
Morphisms, Lemma \ref{morphisms-lemma-universal-homeomorphism}.
In particular, the diagonal $\Delta : S' \to S' \times_S S'$ is a thickening
by Morphisms, Lemma \ref{morphisms-lemma-universally-injective}.
Thus by \'Etale Morphisms, Theorem
\ref{etale-theorem-etale-topological}
we see that given $X' \to S'$ \'etale there is a unique isomorphism
$$
\varphi' : X' \times_S S' \to S' \times_S X'
$$
of schemes \'etale over $S' \times_S S'$ which pulls back under
$\Delta$ to $\text{id} : X' \to X'$ over $S'$.
Since $S' \to S' \times_S S' \times_S S'$
is a thickening as well (it is bijective and a closed immersion)
we conclude that $(X', \varphi')$ is a descent datum relative to $S'/S$.
The canonical nature of the construction of $\varphi'$ shows
that it is compatible with morphisms between schemes \'etale over $S'$.
In other words, we obtain a quasi-inverse
$X' \mapsto (X', \varphi')$ of the functor
(\ref{equation-descent-etale-forget}). This proves the claim and
finishes the proof of the theorem.
\end{proof}

\begin{theorem}
\label{theorem-topological-invariance}
\begin{reference}
\cite[IV Theorem 18.1.2]{EGA}
\end{reference}
Let $f : X \to Y$ be a morphism of schemes.
Assume $f$ is integral, universally injective and surjective
(i.e., $f$ is a universal homeomorphism, see
Morphisms, Lemma \ref{morphisms-lemma-universal-homeomorphism}).
The functor
$$
V \longmapsto V_X = X \times_Y V
$$
defines an equivalence of categories
$$
\{
\text{schemes }V\text{ \'etale over }Y
\}
\leftrightarrow
\{
\text{schemes }U\text{ \'etale over }X
\}
$$
\end{theorem}

\noindent
We give two proofs. The first uses effectivity of descent
for quasi-compact, separated, \'etale morphisms relative
to surjective integral morphisms. The second uses the material
on properties (A), (B), and (C) discussed earlier in the chapter.

\begin{proof}[First proof]
By Theorem \ref{theorem-etale-topological}
we see that the functor is fully faithful.
It remains to show that the functor is essentially surjective.
Let $U \to X$ be an \'etale morphism of schemes.

\medskip\noindent
Suppose that the result holds if $U$ and $Y$ are affine.
In that case, we choose an affine open covering
$U = \bigcup U_i$ such that each $U_i$ maps
into an affine open of $Y$. By assumption (affine case) we can
find \'etale morphisms $V_i \to Y$ such that $X \times_Y V_i \cong U_i$
as schemes over $X$. Let $V_{i, i'} \subset V_i$
be the open subscheme whose underlying topological space
corresponds to $U_i \cap U_{i'}$. Because we have isomorphisms
$$
X \times_Y V_{i, i'} \cong U_i \cap U_{i'} \cong X \times_Y V_{i', i}
$$
as schemes over $X$ we see by fully faithfulness that we
obtain isomorphisms
$\theta_{i, i'} : V_{i, i'} \to V_{i', i}$ of schemes over $Y$.
We omit the verification that these isomorphisms satisfy the
cocycle condition of Schemes, Section \ref{schemes-section-glueing-schemes}.
Applying Schemes, Lemma \ref{schemes-lemma-glue-schemes}
we obtain a scheme $V \to Y$ by
glueing the schemes $V_i$ along the identifications $\theta_{i, i'}$.
It is clear that $V \to Y$ is \'etale and $X \times_Y V \cong U$
by construction.

\medskip\noindent
Thus it suffices to show the lemma in case $U$ and $Y$ are affine.
Recall that in the proof of Theorem \ref{theorem-etale-topological}
we showed that $U$ comes with a unique descent datum
$(U, \varphi)$ relative to $X/Y$. By
\'Etale Morphisms, Proposition \ref{etale-proposition-effective}
(which applies because $U \to X$ is quasi-compact and separated
as well as \'etale by our reduction to the affine case)
there exists an \'etale morphism $V \to Y$ such that
$X \times_Y V \cong U$ and the proof is complete.
\end{proof}

\begin{proof}[Second proof]
By Theorem \ref{theorem-etale-topological}
we see that the functor is fully faithful.
It remains to show that the functor is essentially surjective.
Let $U \to X$ be an \'etale morphism of schemes.

\medskip\noindent
Suppose that the result holds if $U$ and $Y$ are affine.
In that case, we choose an affine open covering
$U = \bigcup U_i$ such that each $U_i$ maps
into an affine open of $Y$. By assumption (affine case) we can
find \'etale morphisms $V_i \to Y$ such that $X \times_Y V_i \cong U_i$
as schemes over $X$. Let $V_{i, i'} \subset V_i$
be the open subscheme whose underlying topological space
corresponds to $U_i \cap U_{i'}$. Because we have isomorphisms
$$
X \times_Y V_{i, i'} \cong U_i \cap U_{i'} \cong X \times_Y V_{i', i}
$$
as schemes over $X$ we see by fully faithfulness that we
obtain isomorphisms
$\theta_{i, i'} : V_{i, i'} \to V_{i', i}$ of schemes over $Y$.
We omit the verification that these isomorphisms satisfy the
cocycle condition of Schemes, Section \ref{schemes-section-glueing-schemes}.
Applying Schemes, Lemma \ref{schemes-lemma-glue-schemes}
we obtain a scheme $V \to Y$ by
glueing the schemes $V_i$ along the identifications $\theta_{i, i'}$.
It is clear that $V \to Y$ is \'etale and $X \times_Y V \cong U$
by construction.

\medskip\noindent
Thus it suffices to prove that the functor
\begin{equation}
\label{equation-affine-etale}
\{
\text{affine schemes }V\text{ \'etale over }Y
\}
\leftrightarrow
\{
\text{affine schemes }U\text{ \'etale over }X
\}
\end{equation}
is essentially surjective when $X$ and $Y$ are affine.

\medskip\noindent
Let $U \to X$ be an affine scheme \'etale over $X$.
We have to find $V \to Y$ \'etale (and affine) such that $X \times_Y V$
is isomorphic to $U$ over $X$. Note that an \'etale morphism of affines
has universally bounded fibres, see
Morphisms,
Lemmas \ref{morphisms-lemma-etale-locally-quasi-finite} and
\ref{morphisms-lemma-locally-quasi-finite-qc-source-universally-bounded}.
Hence we can do induction on the integer $n$ bounding the degree of the fibres
of $U \to X$. See
Morphisms, Lemma \ref{morphisms-lemma-etale-universally-bounded}
for a description of this integer in the case of an \'etale morphism.
If $n = 1$, then $U \to X$ is an open immersion (see
\'Etale Morphisms, Theorem \ref{etale-theorem-etale-radicial-open}),
and the result is clear. Assume $n > 1$.

\medskip\noindent
By
Lemma \ref{lemma-integral-homeo-onto-image-C}
there exists an \'etale morphism of schemes $W \to Y$ and a
surjective morphism $W_X \to U$ over $X$.
As $U$ is quasi-compact we may replace $W$ by a disjoint union of
finitely many affine opens of $W$, hence we may assume that $W$
is affine as well. Here is a diagram
$$
\xymatrix{
U \ar[d] & U \times_Y W \ar[l] \ar[d] & W_X \amalg R \ar@{=}[l]\\
X \ar[d] & W_X \ar[l] \ar[d] \\
Y & W \ar[l]
}
$$
The disjoint union decomposition arises because by construction the
\'etale morphism of affine schemes $U \times_Y W \to W_X$ has a section.
OK, and now we see that the morphism $R \to X \times_Y W$ is an \'etale
morphism of affine schemes whose fibres have degree universally bounded
by $n - 1$. Hence by induction assumption there exists a scheme
$V' \to W$ \'etale such that $R \cong W_X \times_W V'$.
Taking $V'' = W \amalg V'$ we find a scheme $V''$ \'etale over $W$ whose
base change to $W_X$ is isomorphic to $U \times_Y W$
over $X \times_Y W$.

\medskip\noindent
At this point we can use descent to find $V$ over $Y$ whose base
change to $X$ is isomorphic to $U$ over $X$. Namely, by the fully
faithfulness of the functor (\ref{equation-affine-etale})
corresponding to the universal homeomorphism
$X \times_Y (W \times_Y W) \to (W \times_Y W)$
there exists a unique isomorphism $\varphi : V'' \times_Y W \to W \times_Y V''$
whose base change to $X \times_Y (W \times_Y W)$ is the canonical
descent datum for $U \times_Y W$ over $X \times_Y W$. In particular
$\varphi$ satisfies the cocycle condition. Hence by
Descent, Lemma \ref{descent-lemma-affine}
we see that $\varphi$ is effective (recall that all schemes above are affine).
Thus we obtain $V \to Y$ and an isomorphism $V'' \cong W \times_Y V$
such that the canonical descent datum on $W \times_Y V/W/Y$ agrees
with $\varphi$. Note that $V \to Y$ is \'etale, by
Descent, Lemma \ref{descent-lemma-descending-property-etale}.
Moreover, there is an isomorphism $V_X \cong U$ which comes from
descending the isomorphism
$$
V_X \times_X W_X =
X \times_Y V \times_Y W =
(X \times_Y W) \times_W (W \times_Y V) \cong
W_X  \times_W V'' \cong U \times_Y W
$$
which we have by construction. Some details omitted.
\end{proof}

\begin{remark}
\label{remark-affine-inside-equivalence}
In the situation of
Theorem \ref{theorem-topological-invariance}
it is also true that $V \mapsto V_X$ induces an equivalence
between those \'etale morphisms $V \to Y$ with $V$ affine and
those \'etale morphisms $U \to X$ with $U$ affine.
This follows for example from
Limits, Proposition \ref{limits-proposition-affine}.
\end{remark}

\begin{proposition}[Topological invariance of \'etale cohomology]
\label{proposition-topological-invariance}
Let $X_0 \to X$ be a universal homeomorphism of schemes
(for example the closed immersion defined by a nilpotent sheaf of ideals).
Then
\begin{enumerate}
\item the \'etale sites $X_\etale$ and $(X_0)_\etale$ are isomorphic,
\item the \'etale topoi $\Sh(X_\etale)$ and $\Sh((X_0)_\etale)$
are equivalent, and
\item $H^q_\etale(X, \mathcal{F}) = H^q_\etale(X_0, \mathcal{F}|_{X_0})$
for all $q$ and
for any abelian sheaf $\mathcal{F}$ on $X_\etale$.
\end{enumerate}
\end{proposition}

\begin{proof}
The equivalence of categories $X_\etale \to (X_0)_\etale$ is
given by Theorem \ref{theorem-topological-invariance}. We omit
the proof that under this equivalence the \'etale coverings correspond.
Hence (1) holds. Parts (2) and (3) follow formally from (1).
\end{proof}






\section{Closed immersions and pushforward}
\label{section-closed-immersions}

\noindent
Before stating and proving
Proposition \ref{proposition-closed-immersion-pushforward}
in its correct generality we briefly state and prove it for
closed immersions. Namely, some of the preceding arguments
are quite a bit easier to follow in the case of a closed immersion and
so we repeat them here in their simplified form.

\medskip\noindent
In the rest of this section $i : Z \to X$ is a closed immersion.
The functor
$$
\Sch/X \longrightarrow \Sch/Z, \quad
U \longmapsto U_Z = Z \times_X U
$$
will be denoted $U \mapsto U_Z$ as indicated. Since being a closed immersion
is preserved under arbitrary base change the scheme $U_Z$ is a closed subscheme
of $U$.

\begin{lemma}
\label{lemma-closed-immersion-almost-full}
Let $i : Z \to X$ be a closed immersion of schemes.
Let $U, U'$ be schemes \'etale over $X$. Let $h : U_Z \to U'_Z$
be a morphism over $Z$. Then there exists a diagram
$$
\xymatrix{
U & W \ar[l]_a \ar[r]^b & U'
}
$$
in $X_\etale$ such that $a_Z : W_Z \to U_Z$
is an isomorphism and $h = b_Z \circ (a_Z)^{-1}$.
\end{lemma}

\begin{proof}
Consider the scheme $M = U \times_X U'$. The graph $\Gamma_h \subset M_Z$
of $h$ is open. This is true for example as $\Gamma_h$ is the image of a
section of the \'etale morphism $\text{pr}_{1, Z} : M_Z \to U_Z$, see
\'Etale Morphisms, Proposition \ref{etale-proposition-properties-sections}.
Hence there exists an open subscheme $W \subset M$ whose intersection with
the closed subset $M_Z$ is $\Gamma_h$. Set $a = \text{pr}_1|_W$
and $b = \text{pr}_2|_W$.
\end{proof}

\begin{lemma}
\label{lemma-closed-immersion-almost-essentially-surjective}
Let $i : Z \to X$ be a closed immersion of schemes.
Let $V \to Z$ be an \'etale morphism of schemes.
There exist \'etale morphisms $U_i \to X$ and morphisms
$U_{i, Z} \to V$ such that $\{U_{i, Z} \to V\}$
is a Zariski covering of $V$.
\end{lemma}

\begin{proof}
Since we only have to find a Zariski covering of $V$ consisting of schemes
of the form $U_Z$ with $U$ \'etale over $X$, we may Zariski localize on $X$
and $V$. Hence we may assume $X$ and $V$ affine. In the affine case this is
Algebra, Lemma \ref{algebra-lemma-lift-etale}.
\end{proof}

\noindent
If $\overline{x} : \Spec(k) \to X$ is a geometric point of $X$, then
either $\overline{x}$ factors (uniquely) through the closed subscheme $Z$, or
$Z_{\overline{x}} = \emptyset$. If $\overline{x}$ factors through $Z$
we say that $\overline{x}$ is a geometric point of $Z$ (because it is) and
we use the notation ``$\overline{x} \in Z$'' to indicate this.

\begin{lemma}
\label{lemma-stalk-pushforward-closed-immersion}
Let $i : Z \to X$ be a closed immersion of schemes.
Let $\mathcal{G}$ be a sheaf of sets on $Z_\etale$.
Let $\overline{x}$ be a geometric point of $X$.
Then
$$
(i_{small, *}\mathcal{G})_{\overline{x}} =
\left\{
\begin{matrix}
* & \text{if} & \overline{x} \not \in Z \\
\mathcal{G}_{\overline{x}} & \text{if} & \overline{x} \in Z
\end{matrix}
\right.
$$
where $*$ denotes a singleton set.
\end{lemma}

\begin{proof}
Let $U = X \setminus Z$.
Note that $i_{small, *}\mathcal{G}|_{U_\etale} = *$ is the final
object in the category of \'etale sheaves on $U$, i.e., the sheaf
which associates a singleton set to each scheme \'etale over $U$.
This explains the value of $(i_{small, *}\mathcal{G})_{\overline{x}}$
if $\overline{x} \not \in Z$.

\medskip\noindent
Next, suppose that $\overline{x} \in Z$. Note that
$$
(i_{small, *}\mathcal{G})_{\overline{x}}
=
\colim_{(U, \overline{u})} \mathcal{G}(U_Z)
$$
and on the other hand
$$
\mathcal{G}_{\overline{x}}
=
\colim_{(V, \overline{v})} \mathcal{G}(V).
$$
Let $\mathcal{C}_1 = \{(U, \overline{u})\}^{opp}$ be the opposite of the
category of \'etale neighbourhoods of $\overline{x}$ in $X$, and let
$\mathcal{C}_2 = \{(V, \overline{v})\}^{opp}$ be the opposite of the
category of \'etale neighbourhoods of $\overline{x}$ in $Z$. The canonical map
$$
\mathcal{G}_{\overline{x}}
\longrightarrow
(i_{small, *}\mathcal{G})_{\overline{x}}
$$
corresponds to the functor $F : \mathcal{C}_1 \to \mathcal{C}_2$,
$F(U, \overline{u}) = (U_Z, \overline{x})$. Now
Lemmas \ref{lemma-closed-immersion-almost-essentially-surjective} and
\ref{lemma-closed-immersion-almost-full}
imply that $\mathcal{C}_1$ is cofinal in $\mathcal{C}_2$, see
Categories, Definition \ref{categories-definition-cofinal}.
Hence it follows that the displayed arrow is an isomorphism, see
Categories, Lemma \ref{categories-lemma-cofinal}.
\end{proof}

\begin{proposition}
\label{proposition-closed-immersion-pushforward}
Let $i : Z \to X$ be a closed immersion of schemes.
\begin{enumerate}
\item The functor
$$
i_{small, *} :
\Sh(Z_\etale)
\longrightarrow
\Sh(X_\etale)
$$
is fully faithful and its essential image is those sheaves of sets
$\mathcal{F}$ on $X_\etale$ whose restriction to $X \setminus Z$ is
isomorphic to $*$, and
\item the functor
$$
i_{small, *} :
\textit{Ab}(Z_\etale)
\longrightarrow
\textit{Ab}(X_\etale)
$$
is fully faithful and its essential image is those abelian sheaves on
$X_\etale$ whose support is contained in $Z$.
\end{enumerate}
In both cases $i_{small}^{-1}$ is a left inverse to the functor
$i_{small, *}$.
\end{proposition}

\begin{proof}
Let's discuss the case of sheaves of sets.
For any sheaf $\mathcal{G}$ on $Z$ the morphism
$i_{small}^{-1}i_{small, *}\mathcal{G} \to \mathcal{G}$
is an isomorphism by
Lemma \ref{lemma-stalk-pushforward-closed-immersion}
(and
Theorem \ref{theorem-exactness-stalks}).
This implies formally that $i_{small, *}$ is fully faithful, see
Sites, Lemma \ref{sites-lemma-exactness-properties}.
It is clear that $i_{small, *}\mathcal{G}|_{U_\etale} \cong *$
where $U = X \setminus Z$. Conversely, suppose that $\mathcal{F}$
is a sheaf of sets on $X$ such that $\mathcal{F}|_{U_\etale} \cong *$.
Consider the adjunction mapping
$$
\mathcal{F} \longrightarrow i_{small, *}i_{small}^{-1}\mathcal{F}
$$
Combining
Lemmas \ref{lemma-stalk-pushforward-closed-immersion} and
\ref{lemma-stalk-pullback}
we see that it is an isomorphism. This finishes the proof of (1).
The proof of (2) is identical.
\end{proof}





\section{Integral universally injective morphisms}
\label{section-integral-universally-injective}

\noindent
Here is the general version of
Proposition \ref{proposition-closed-immersion-pushforward}.

\begin{proposition}
\label{proposition-integral-universally-injective-pushforward}
Let $f : X \to Y$ be a morphism of schemes which is integral
and universally injective.
\begin{enumerate}
\item The functor
$$
f_{small, *} :
\Sh(X_\etale)
\longrightarrow
\Sh(Y_\etale)
$$
is fully faithful and its essential image is those sheaves of sets
$\mathcal{F}$ on $Y_\etale$ whose restriction to $Y \setminus f(X)$ is
isomorphic to $*$, and
\item the functor
$$
f_{small, *} :
\textit{Ab}(X_\etale)
\longrightarrow
\textit{Ab}(Y_\etale)
$$
is fully faithful and its essential image is those abelian sheaves on
$Y_\etale$ whose support is contained in $f(X)$.
\end{enumerate}
In both cases $f_{small}^{-1}$ is a left inverse to the functor
$f_{small, *}$.
\end{proposition}

\begin{proof}
We may factor $f$ as
$$
\xymatrix{
X \ar[r]^h & Z \ar[r]^i & Y
}
$$
where $h$ is integral, universally injective and surjective
and $i : Z \to Y$ is a closed immersion.
Apply
Proposition \ref{proposition-closed-immersion-pushforward}
to $i$ and apply
Theorem \ref{theorem-topological-invariance}
to $h$.
\end{proof}









\section{Big sites and pushforward}
\label{section-big}

\noindent
In this section we prove some technical results on $f_{big, *}$ for
certain types of morphisms of schemes.

\begin{lemma}
\label{lemma-monomorphism-big-push-pull}
Let $\tau \in \{Zariski, \etale, smooth, syntomic, fppf\}$.
Let $f : X \to Y$ be a monomorphism of schemes.
Then the canonical map
$f_{big}^{-1}f_{big, *}\mathcal{F} \to \mathcal{F}$
is an isomorphism for any sheaf $\mathcal{F}$ on
$(\Sch/X)_\tau$.
\end{lemma}

\begin{proof}
In this case the functor $(\Sch/X)_\tau \to (\Sch/Y)_\tau$
is continuous, cocontinuous and fully faithful. Hence the result follows from
Sites, Lemma \ref{sites-lemma-back-and-forth}.
\end{proof}

\begin{remark}
\label{remark-push-pull-shriek}
In the situation of
Lemma \ref{lemma-monomorphism-big-push-pull}
it is true that the canonical map
$\mathcal{F} \to f_{big}^{-1}f_{big!}\mathcal{F}$
is an isomorphism for any sheaf of sets $\mathcal{F}$ on
$(\Sch/X)_\tau$. The proof is the same. This also
holds for sheaves of abelian groups. However, note
that the functor $f_{big!}$ for sheaves of abelian groups is defined in
Modules on Sites, Section \ref{sites-modules-section-exactness-lower-shriek}
and is in general different from $f_{big!}$ on sheaves of sets.
The result for sheaves of abelian groups follows from
Modules on Sites, Lemma \ref{sites-modules-lemma-back-and-forth}.
\end{remark}

\begin{lemma}
\label{lemma-closed-immersion-cover-from-below}
Let $f : X \to Y$ be a closed immersion of schemes.
Let $U \to X$ be a syntomic (resp.\ smooth, resp.\ \'etale) morphism.
Then there exist syntomic (resp.\ smooth, resp.\ \'etale) morphisms
$V_i \to Y$ and morphisms $V_i \times_Y X \to U$ such that
$\{V_i \times_Y X \to U\}$ is a Zariski covering of $U$.
\end{lemma}

\begin{proof}
Let us prove the lemma when $\tau = syntomic$.
The question is local on $U$. Thus we may assume that $U$ is
an affine scheme mapping into an affine of $Y$.
Hence we reduce to proving the following case:
$Y = \Spec(A)$, $X = \Spec(A/I)$, and
$U = \Spec(\overline{B})$, where
$A/I \to \overline{B}$ be a syntomic ring map.
By Algebra, Lemma \ref{algebra-lemma-lift-syntomic}
we can find elements $\overline{g}_i \in \overline{B}$
such that
$\overline{B}_{\overline{g}_i} = A_i/IA_i$ for certain syntomic ring maps
$A \to A_i$.
This proves the lemma in the syntomic case.
The proof of the smooth case is the same except it uses
Algebra, Lemma \ref{algebra-lemma-lift-smooth}.
In the \'etale case use
Algebra, Lemma \ref{algebra-lemma-lift-etale}.
\end{proof}

\begin{lemma}
\label{lemma-prepare-closed-immersion-almost-cocontinuous}
Let $f : X \to Y$ be a closed immersion of schemes.
Let $\{U_i \to X\}$ be a syntomic (resp.\ smooth, resp.\ \'etale) covering.
There exists a syntomic (resp.\ smooth, resp.\ \'etale) covering $\{V_j \to Y\}$
such that for each $j$, either $V_j \times_Y X = \emptyset$, or the
morphism $V_j \times_Y X \to X$ factors through $U_i$ for some $i$.
\end{lemma}

\begin{proof}
For each $i$ we can choose syntomic (resp.\ smooth, resp.\ \'etale) morphisms
$g_{ij} : V_{ij} \to Y$ and morphisms $V_{ij} \times_Y X \to U_i$ over $X$,
such that $\{V_{ij} \times_Y X \to U_i\}$ are Zariski coverings, see
Lemma \ref{lemma-closed-immersion-cover-from-below}.
This in particular implies that
$\bigcup_{ij} g_{ij}(V_{ij})$ contains the closed subset $f(X)$.
Hence the family of syntomic (resp.\ smooth, resp.\ \'etale) maps $g_{ij}$
together with the open immersion $Y \setminus f(X) \to Y$ forms the desired
syntomic (resp.\ smooth, resp.\ \'etale) covering of $Y$.
\end{proof}

\begin{lemma}
\label{lemma-closed-immersion-almost-cocontinuous}
Let $f : X \to Y$ be a closed immersion of schemes.
Let $\tau \in \{syntomic, smooth, \etale\}$.
The functor $V \mapsto X \times_Y V$ defines an almost
cocontinuous functor (see
Sites, Definition \ref{sites-definition-almost-cocontinuous})
$(\Sch/Y)_\tau \to (\Sch/X)_\tau$ between
big $\tau$ sites.
\end{lemma}

\begin{proof}
We have to show the following: given a morphism $V \to Y$
and any syntomic (resp.\ smooth, resp.\ \'etale)
covering $\{U_i \to X \times_Y V\}$, there exists a
smooth (resp.\ smooth, resp.\ \'etale) covering $\{V_j \to V\}$
such that for each $j$, either $X \times_Y V_j$ is empty, or
$X \times_Y V_j \to Z \times_Y V$ factors through one of
the $U_i$. This follows on applying
Lemma \ref{lemma-prepare-closed-immersion-almost-cocontinuous}
above to the closed immersion $X \times_Y V \to V$.
\end{proof}

\begin{lemma}
\label{lemma-closed-immersion-pushforward-exact}
Let $f : X \to Y$ be a closed immersion of schemes.
Let $\tau \in \{syntomic, smooth, \etale\}$.
\begin{enumerate}
\item The pushforward
$f_{big, *} :
\Sh((\Sch/X)_\tau)
\to
\Sh((\Sch/Y)_\tau)$
commutes with coequalizers and pushouts.
\item The pushforward
$f_{big, *} :
\textit{Ab}((\Sch/X)_\tau)
\to
\textit{Ab}((\Sch/Y)_\tau)$
is exact.
\end{enumerate}
\end{lemma}

\begin{proof}
This follows from
Sites, Lemma \ref{sites-lemma-morphism-of-sites-almost-cocontinuous},
Modules on Sites,
Lemma \ref{sites-modules-lemma-morphism-ringed-sites-almost-cocontinuous},
and
Lemma \ref{lemma-closed-immersion-almost-cocontinuous}
above.
\end{proof}

\begin{remark}
\label{remark-fppf-closed-immersion-not-closed}
In Lemma \ref{lemma-closed-immersion-pushforward-exact} the case $\tau = fppf$
is missing. The reason is that given a ring $A$, an ideal $I$ and a
faithfully flat, finitely presented ring map $A/I \to \overline{B}$, there
is no reason to think that one can find {\it any} flat finitely presented ring
map $A \to B$ with $B/IB \not = 0$ such that $A/I \to B/IB$ factors through
$\overline{B}$. Hence the proof of
Lemma \ref{lemma-closed-immersion-almost-cocontinuous}
does not work for the fppf topology.
In fact it is likely false that
$f_{big, *} : \textit{Ab}((\Sch/X)_{fppf})
\to \textit{Ab}((\Sch/Y)_{fppf})$
is exact when $f$ is a closed immersion.
If you know an example, please email
\href{mailto:stacks.project@gmail.com}{stacks.project@gmail.com}.
\end{remark}












\section{Exactness of big lower shriek}
\label{section-exactness-lower-shriek}

\noindent
This is just the following technical result. Note that the functor $f_{big!}$
has nothing whatsoever to do with cohomology with compact support in
general.

\begin{lemma}
\label{lemma-exactness-lower-shriek}
Let $\tau \in \{Zariski, \etale, smooth, syntomic, fppf\}$.
Let $f : X \to Y$ be a morphism of schemes. Let
$$
f_{big} :
\Sh((\Sch/X)_\tau)
\longrightarrow
\Sh((\Sch/Y)_\tau)
$$
be the corresponding morphism of topoi as in
Topologies, Lemma
\ref{topologies-lemma-morphism-big},
\ref{topologies-lemma-morphism-big-etale},
\ref{topologies-lemma-morphism-big-smooth},
\ref{topologies-lemma-morphism-big-syntomic}, or
\ref{topologies-lemma-morphism-big-fppf}.
\begin{enumerate}
\item The functor
$f_{big}^{-1} : \textit{Ab}((\Sch/Y)_\tau) \to \textit{Ab}((\Sch/X)_\tau)$
has a left adjoint
$$
f_{big!} : \textit{Ab}((\Sch/X)_\tau) \to \textit{Ab}((\Sch/Y)_\tau)
$$
which is exact.
\item The functor
$f_{big}^* :
\textit{Mod}((\Sch/Y)_\tau, \mathcal{O})
\to
\textit{Mod}((\Sch/X)_\tau, \mathcal{O})$
has a left adjoint
$$
f_{big!} :
\textit{Mod}((\Sch/X)_\tau, \mathcal{O})
\to
\textit{Mod}((\Sch/Y)_\tau, \mathcal{O})
$$
which is exact.
\end{enumerate}
Moreover, the two functors $f_{big!}$ agree on underlying sheaves
of abelian groups.
\end{lemma}

\begin{proof}
Recall that $f_{big}$ is the morphism of topoi associated to the
continuous and cocontinuous functor
$u : (\Sch/X)_\tau \to (\Sch/Y)_\tau$, $U/X \mapsto U/Y$.
Moreover, we have $f_{big}^{-1}\mathcal{O} = \mathcal{O}$.
Hence the existence of $f_{big!}$ follows from
Modules on Sites, Lemma \ref{sites-modules-lemma-g-shriek-adjoint},
respectively
Modules on Sites, Lemma \ref{sites-modules-lemma-lower-shriek-modules}.
Note that if $U$ is an object of $(\Sch/X)_\tau$ then the functor
$u$ induces an equivalence of categories
$$
u' :
(\Sch/X)_\tau/U
\longrightarrow
(\Sch/Y)_\tau/U
$$
because both sides of the arrow are equal to $(\Sch/U)_\tau$.
Hence the agreement of $f_{big!}$ on underlying abelian sheaves
follows from the discussion in
Modules on Sites, Remark \ref{sites-modules-remark-when-shriek-equal}.
The exactness of $f_{big!}$ follows from
Modules on Sites, Lemma \ref{sites-modules-lemma-exactness-lower-shriek}
as the functor $u$ above which commutes with fibre products and equalizers.
\end{proof}

\noindent
Next, we prove a technical lemma that will be useful later when comparing
sheaves of modules on different sites associated to algebraic stacks.

\begin{lemma}
\label{lemma-compare-structure-sheaves}
Let $X$ be a scheme. Let
$\tau \in \{Zariski, \etale, smooth, syntomic, fppf\}$.
Let $\mathcal{C}_1 \subset \mathcal{C}_2 \subset (\Sch/X)_\tau$ be full
subcategories with the following properties:
\begin{enumerate}
\item For an object $U/X$ of $\mathcal{C}_t$,
\begin{enumerate}
\item if $\{U_i \to U\}$ is a covering of $(\Sch/X)_\tau$, then
$U_i/X$ is an object of $\mathcal{C}_t$,
\item $U \times \mathbf{A}^1/X$ is an object of $\mathcal{C}_t$.
\end{enumerate}
\item $X/X$ is an object of $\mathcal{C}_t$.
\end{enumerate}
We endow $\mathcal{C}_t$ with the structure of a site whose coverings are
exactly those coverings $\{U_i \to U\}$ of $(\Sch/X)_\tau$ with
$U \in \Ob(\mathcal{C}_t)$. Then
\begin{enumerate}
\item[(a)] The functor $\mathcal{C}_1 \to \mathcal{C}_2$
is fully faithful, continuous, and cocontinuous.
\end{enumerate}
Denote $g : \Sh(\mathcal{C}_1) \to \Sh(\mathcal{C}_2)$ the corresponding
morphism of topoi. Denote $\mathcal{O}_t$ the restriction of $\mathcal{O}$
to $\mathcal{C}_t$. Denote $g_!$ the functor of
Modules on Sites, Definition \ref{sites-modules-definition-g-shriek}.
\begin{enumerate}
\item[(b)] The canonical map $g_!\mathcal{O}_1 \to \mathcal{O}_2$
is an isomorphism.
\end{enumerate}
\end{lemma}

\begin{proof}
Assertion (a) is immediate from the definitions.
In this proof all schemes are schemes over $X$ and all morphisms of
schemes are morphisms of schemes over $X$. Note that $g^{-1}$ is
given by restriction, so that for an object $U$ of $\mathcal{C}_1$
we have $\mathcal{O}_1(U) = \mathcal{O}_2(U) = \mathcal{O}(U)$.
Recall that $g_!\mathcal{O}_1$ is the sheaf associated to the presheaf
$g_{p!}\mathcal{O}_1$ which associates to $V$ in $\mathcal{C}_2$ the group
$$
\colim_{V \to U} \mathcal{O}(U)
$$
where $U$ runs over the objects of $\mathcal{C}_1$ and the colimit is
taken in the category of abelian groups. Below we will use frequently
that if
$$
V \to U \to U'
$$
are morphisms with $U, U' \in \Ob(\mathcal{C}_1)$
and if $f' \in \mathcal{O}(U')$ restricts to $f \in \mathcal{O}(U)$,
then $(V \to U, f)$ and $(V \to U', f')$ define the same element of the
colimit. Also, $g_!\mathcal{O}_1 \to \mathcal{O}_2$ maps the element
$(V \to U, f)$ simply to the pullback of $f$ to $V$.

\medskip\noindent
Surjectivity. Let $V$ be a scheme and let $h \in \mathcal{O}(V)$.
Then we obtain a morphism $V \to X \times \mathbf{A}^1$ induced by $h$
and the structure morphism $V \to X$. Writing
$\mathbf{A}^1 = \Spec(\mathbf{Z}[x])$ we see the element
$x \in \mathcal{O}(X \times \mathbf{A}^1)$ pulls
back to $h$. Since $X \times \mathbf{A}^1$ is an object of $\mathcal{C}_1$
by assumptions (1)(b) and (2) we obtain the desired surjectivity.

\medskip\noindent
Injectivity. Let $V$ be a scheme. Let
$s = \sum_{i = 1, \ldots, n} (V \to U_i, f_i)$ be an element of the colimit
displayed above. For any $i$ we can use the morphism
$f_i : U_i \to X \times \mathbf{A}^1$
to see that $(V \to U_i, f_i)$ defines the same element of the colimit as
$(f_i : V \to X \times \mathbf{A}^1, x)$. Then we can consider
$$
f_1 \times \ldots \times f_n : V \to X \times \mathbf{A}^n
$$
and we see that $s$ is equivalent in the colimit to
$$
\sum\nolimits_{i = 1, \ldots, n}
(f_1 \times \ldots \times f_n : V \to X \times \mathbf{A}^n, x_i) =
(f_1 \times \ldots \times f_n : V \to X \times \mathbf{A}^n,
x_1 + \ldots + x_n)
$$
Now, if $x_1 + \ldots + x_n$ restricts to zero on $V$, then we see
that $f_1 \times \ldots \times f_n$ factors through
$X \times \mathbf{A}^{n - 1} = V(x_1 + \ldots + x_n)$. Hence we see
that $s$ is equivalent to zero in the colimit.
\end{proof}











%9.29.09
\section{\'Etale cohomology}
\label{section-etale-cohomology}

\noindent
In the following sections we prove some basic results on \'etale cohomology.
Here is an example of something we know for cohomology of topological
spaces which also holds for \'etale cohomology.

\begin{lemma}[Mayer-Vietoris for \'etale cohomology]
\label{lemma-mayer-vietoris}
Let $X$ be a scheme. Suppose that $X = U \cup V$ is a
union of two opens. For any abelian sheaf $\mathcal{F}$ on $X_\etale$
there exists a long exact cohomology sequence
$$
\begin{matrix}
0 \to
H^0_\etale(X, \mathcal{F}) \to
H^0_\etale(U, \mathcal{F}) \oplus H^0_\etale(V, \mathcal{F}) \to
H^0_\etale(U \cap V, \mathcal{F}) \phantom{\to \ldots} \\
\phantom{0} \to H^1_\etale(X, \mathcal{F}) \to
H^1_\etale(U, \mathcal{F}) \oplus H^1_\etale(V, \mathcal{F}) \to
H^1_\etale(U \cap V, \mathcal{F}) \to \ldots
\end{matrix}
$$
This long exact sequence is functorial in $\mathcal{F}$.
\end{lemma}

\begin{proof}
Observe that if $\mathcal{I}$ is an injective abelian sheaf, then
$$
0 \to \mathcal{I}(X) \to \mathcal{I}(U) \oplus \mathcal{I}(V) \to
\mathcal{I}(U \cap V) \to 0
$$
is exact. This is true in the first and middle spots as $\mathcal{I}$
is a sheaf. It is true on the right, because
$\mathcal{I}(U) \to \mathcal{I}(U \cap V)$ is surjective by
Cohomology on Sites, Lemma
\ref{sites-cohomology-lemma-restriction-along-monomorphism-surjective}.
Another way to prove it would be to show that the cokernel
of the map $\mathcal{I}(U) \oplus \mathcal{I}(V) \to
\mathcal{I}(U \cap V)$ is the first {\v C}ech cohomology group
of $\mathcal{I}$ with respect to the covering
$X = U \cup V$ which vanishes by
Lemmas \ref{lemma-hom-injective} and \ref{lemma-forget-injectives}.
Thus, if $\mathcal{F} \to \mathcal{I}^\bullet$ is an injective
resolution, then
$$
0 \to \mathcal{I}^\bullet(X) \to
\mathcal{I}^\bullet(U) \oplus \mathcal{I}^\bullet(V) \to
\mathcal{I}^\bullet(U \cap V) \to 0
$$
is a short exact sequence of complexes and the associated long
exact cohomology sequence is the sequence of the statement of the lemma.
\end{proof}

\begin{lemma}[Relative Mayer-Vietoris]
\label{lemma-relative-mayer-vietoris}
Let $f : X \to Y$ be a morphism of schemes. Suppose that $X = U \cup V$
is a union of two open subschemes. Denote
$a = f|_U : U \to Y$, $b = f|_V : V \to Y$, and
$c = f|_{U \cap V} : U \cap V \to Y$.
For every abelian sheaf $\mathcal{F}$ on $X_\etale$
there exists a long exact sequence
$$
0 \to
f_*\mathcal{F} \to
a_*(\mathcal{F}|_U) \oplus b_*(\mathcal{F}|_V) \to
c_*(\mathcal{F}|_{U \cap V}) \to
R^1f_*\mathcal{F} \to \ldots
$$
on $Y_\etale$.
This long exact sequence is functorial in $\mathcal{F}$.
\end{lemma}

\begin{proof}
Let $\mathcal{F} \to \mathcal{I}^\bullet$ be an injective resolution
of $\mathcal{F}$ on $X_\etale$. We claim that we
get a short exact sequence of complexes
$$
0 \to
f_*\mathcal{I}^\bullet \to
a_*\mathcal{I}^\bullet|_U \oplus b_*\mathcal{I}^\bullet|_V \to
c_*\mathcal{I}^\bullet|_{U \cap V} \to
0.
$$
Namely, for any $W$ in $Y_\etale$, and for any $n \geq 0$ the
corresponding sequence of groups of sections over $W$
$$
0 \to
\mathcal{I}^n(W \times_Y X) \to
\mathcal{I}^n(W \times_Y U)
\oplus \mathcal{I}^n(W \times_Y V) \to
\mathcal{I}^n(W \times_Y (U \cap V)) \to
0
$$
was shown to be short exact in the proof of Lemma \ref{lemma-mayer-vietoris}.
The lemma follows by taking cohomology sheaves and using the fact that
$\mathcal{I}^\bullet|_U$ is an injective resolution of $\mathcal{F}|_U$
and similarly for $\mathcal{I}^\bullet|_V$, $\mathcal{I}^\bullet|_{U \cap V}$.
\end{proof}







\section{Colimits}
\label{section-colimit}

\noindent
We recall that if $(\mathcal{F}_i, \varphi_{ii'})$ is a diagram of sheaves on
a site $\mathcal{C}$ its colimit (in the category of sheaves) is the
sheafification of the presheaf $U \mapsto \colim_i \mathcal{F}_i(U)$. See
Sites, Lemma \ref{sites-lemma-colimit-sheaves}.
If the system is directed, $U$ is a quasi-compact object of
$\mathcal{C}$ which has a cofinal system of coverings by quasi-compact
objects, then $\mathcal{F}(U) = \colim \mathcal{F}_i(U)$, see
Sites, Lemma \ref{sites-lemma-directed-colimits-sections}.
See Cohomology on Sites, Lemma
\ref{sites-cohomology-lemma-colim-works-over-collection}
for a result dealing with higher cohomology groups of colimits
of abelian sheaves.

\medskip\noindent
In Cohomology on Sites, Lemma \ref{sites-cohomology-lemma-colimit}
we generalize this result to a system of sheaves
on an inverse system of sites. Here is the corresponding
notion in the case of a system of \'etale sheaves living
on an inverse system of schemes.

\begin{definition}
\label{definition-inverse-system-sheaves}
Let $I$ be a preordered set. Let $(X_i, f_{i'i})$ be an inverse
system of schemes over $I$.
A {\it system $(\mathcal{F}_i, \varphi_{i'i})$ of sheaves
on $(X_i, f_{i'i})$} is given by
\begin{enumerate}
\item a sheaf $\mathcal{F}_i$ on $(X_i)_\etale$ for all $i \in I$,
\item for $i' \geq i$ a map
$\varphi_{i'i} : f_{i'i}^{-1}\mathcal{F}_i \to \mathcal{F}_{i'}$
of sheaves on $(X_{i'})_\etale$
\end{enumerate}
such that $\varphi_{i''i} = \varphi_{i''i'} \circ f_{i'' i'}^{-1}\varphi_{i'i}$
whenever $i'' \geq i' \geq i$.
\end{definition}

\noindent
In the situation of Definition \ref{definition-inverse-system-sheaves},
assume $I$ is a directed set and the transition morphisms $f_{i'i}$ affine.
Let $X = \lim X_i$ be the limit in the category of schemes, see
Limits, Section \ref{limits-section-limits}.
Denote $f_i : X \to X_i$ the projection morphisms and consider the maps
$$
f_i^{-1}\mathcal{F}_i = f_{i'}^{-1}f_{i'i}^{-1}\mathcal{F}_i
\xrightarrow{f_{i'}^{-1}\varphi_{i'i}}
f_{i'}^{-1}\mathcal{F}_{i'}
$$
This turns $f_i^{-1}\mathcal{F}_i$ into a system of sheaves on $X_\etale$
over $I$ (it is a good exercise to check this).
We often want to know whether there is an isomorphism
$$
H^q_\etale(X, \colim f_i^{-1}\mathcal{F}_i) =
\colim H^q_\etale(X_i, \mathcal{F}_i)
$$
It will turn out this is true if $X_i$ is quasi-compact and quasi-separated
for all $i$, see Theorem \ref{theorem-colimit}.

\begin{lemma}
\label{lemma-colimit-affine-sites}
Let $I$ be a directed set. Let $(X_i, f_{i'i})$ be an inverse
system of schemes over $I$ with affine transition morphisms.
Let $X = \lim_{i \in I} X_i$. With
notation as in Topologies, Lemma \ref{topologies-lemma-alternative} we have
$$
X_{affine, \etale} = \colim (X_i)_{affine, \etale}
$$
as sites in the sense of
Sites, Lemma \ref{sites-lemma-colimit-sites}.
\end{lemma}

\begin{proof}
Let us first prove this when $X$ and $X_i$ are quasi-compact and
quasi-separated for all $i$ (as this is true in all cases of
interest). In this case any object of
$X_{affine, \etale}$, resp.\ $(X_i)_{affine, \etale}$ is
of finite presentation over $X$. Moreover, the
category of schemes of finite presentation over $X$ is the
colimit of the categories of schemes of finite presentation
over $X_i$, see
Limits, Lemma \ref{limits-lemma-descend-finite-presentation}.
The same holds for the subcategories of affine objects \'etale
over $X$ by Limits, Lemmas
\ref{limits-lemma-limit-affine} and \ref{limits-lemma-descend-etale}.
Finally, if $\{U^j \to U\}$ is a covering of $X_{affine, \etale}$
and if $U_i^j \to U_i$ is morphism of affine schemes \'etale over
$X_i$ whose base change to $X$ is $U^j \to U$, then we see that
the base change of $\{U^j_i \to U_i\}$ to some $X_{i'}$ is
a covering for $i'$ large enough, see
Limits, Lemma \ref{limits-lemma-descend-surjective}.

\medskip\noindent
In the general case, let $U$ be an object of $X_{affine, \etale}$.
Then $U \to X$ is \'etale and separated (as $U$ is separated)
but in general not quasi-compact. Still, $U \to X$ is locally
of finite presentation and hence by
Limits, Lemma \ref{limits-lemma-descend-finite-presentation-variant}
there exists an $i$, a quasi-compact and quasi-separated scheme $U_i$, and
a morphism $U_i \to X_i$ which is locally of finite presentation
whose base change to $X$ is $U \to X$. Then $U = \lim_{i' \geq i} U_{i'}$
where $U_{i'} = U_i \times_{X_i} X_{i'}$.
After increasing $i$ we may assume $U_i$ is affine, see
Limits, Lemma \ref{limits-lemma-limit-affine}.
To check that $U_i \to X_i$ is \'etale for $i$ sufficiently large,
choose a finite affine open covering $U_i = U_{i, 1} \cup \ldots \cup U_{i, m}$
such that $U_{i, j} \to U_i \to X_i$ maps into an affine open
$W_{i, j} \subset X_i$. Then we can apply
Limits, Lemma \ref{limits-lemma-descend-etale}
to see that $U_{i, j} \to W_{i, j}$ is \'etale
after possibly increasing $i$.
In this way we see that the functor
$\colim (X_i)_{affine, \etale} \to X_{affine, \etale}$
is essentially surjective. Fully faithfulness follows
directly from the already used
Limits, Lemma \ref{limits-lemma-descend-finite-presentation-variant}.
The statement on coverings is proved in exactly the same
manner as done in the first paragraph of the proof.
\end{proof}

\noindent
Using the above we get the following general result on colimits and cohomology.

\begin{theorem}
\label{theorem-colimit}
Let $X = \lim_{i \in I} X_i$ be a limit of a directed system of schemes
with affine transition morphisms $f_{i'i} : X_{i'} \to X_i$. We assume
that $X_i$ is quasi-compact and quasi-separated for all $i \in I$.
Let $(\mathcal{F}_i, \varphi_{i'i})$ be a system of abelian sheaves
on $(X_i, f_{i'i})$. Denote $f_i : X \to X_i$ the projection and set
$\mathcal{F} = \colim f_i^{-1}\mathcal{F}_i$. Then
$$
\colim_{i\in I} H_\etale^p(X_i, \mathcal{F}_i) = H_\etale^p(X, \mathcal{F}).
$$
for all $p \geq 0$.
\end{theorem}

\begin{proof}
By Topologies, Lemma \ref{topologies-lemma-alternative}
we can compute the cohomology
of $\mathcal{F}$ on $X_{affine, \etale}$.
Thus the result by a combination of
Lemma \ref{lemma-colimit-affine-sites}
and
Cohomology on Sites, Lemma \ref{sites-cohomology-lemma-colimit}.
\end{proof}

\noindent
The following two results are special cases of the theorem above.

\begin{lemma}
\label{lemma-colimit}
Let $X$ be a quasi-compact and quasi-separated scheme. Let $I$
be a directed set. Let $(\mathcal{F}_i, \varphi_{ij})$ be a system
of abelian sheaves on $X_\etale$ over $I$. Then
$$
\colim_{i\in I} H_\etale^p(X, \mathcal{F}_i) = H_\etale^p(X,
\colim_{i\in I} \mathcal{F}_i).
$$
\end{lemma}

\begin{proof}
This is a special case of Theorem \ref{theorem-colimit}.
We also sketch a direct proof.
We prove it for all $X$ at the same time, by induction on $p$.
\begin{enumerate}
\item
For any quasi-compact and quasi-separated scheme $X$ and any \'etale covering
$\mathcal{U}$ of $X$, show that there exists a refinement
$\mathcal{V} = \{V_j \to X\}_{j\in J}$ with $J$ finite and each $V_j$
quasi-compact and quasi-separated such that all
$V_{j_0} \times_X \ldots \times_X V_{j_p}$ are also
quasi-compact and quasi-separated.
\item
Using the previous step and the definition of colimits in the category of
sheaves, show that the theorem holds for $p = 0$ and all $X$.
\item
Using the locality of cohomology
(Lemma \ref{lemma-locality-cohomology}),
the {\v C}ech-to-cohomology spectral sequence
(Theorem \ref{theorem-cech-ss}) and the fact that the induction
hypothesis applies to all
$V_{j_0} \times_X \ldots \times_X V_{j_p}$
in the above situation, prove the induction step $p \to p + 1$.
\end{enumerate}
\end{proof}

\begin{lemma}
\label{lemma-directed-colimit-cohomology}
Let $A$ be a ring, $(I, \leq)$ a directed set and $(B_i, \varphi_{ij})$ a
system of $A$-algebras. Set $B = \colim_{i\in I} B_i$. Let $X \to \Spec(A)$
be a quasi-compact and quasi-separated morphism of schemes. Let
$\mathcal{F}$ an abelian sheaf on $X_\etale$.
Denote $Y_i = X \times_{\Spec(A)} \Spec(B_i)$,
$Y = X \times_{\Spec(A)} \Spec(B)$,
$\mathcal{G}_i = (Y_i \to X)^{-1}\mathcal{F}$ and
$\mathcal{G} = (Y \to X)^{-1}\mathcal{F}$. Then
$$
H_\etale^p(Y, \mathcal{G}) =
\colim_{i\in I} H_\etale^p (Y_i, \mathcal{G}_i).
$$
\end{lemma}

\begin{proof}
This is a special case of Theorem \ref{theorem-colimit}.
We also outline a direct proof as follows.
\begin{enumerate}
\item Given $V \to Y$ \'etale with $V$ quasi-compact and
quasi-separated, there exist $i\in I$ and $V_i \to Y_i$ such that
$V = V_i \times_{Y_i} Y$.
If all the schemes considered were affine, this would correspond to the
following algebra statement: if $B = \colim B_i$ and $B \to C$ is \'etale,
then there exist $i \in I$ and $B_i \to C_i$ \'etale such that
$C \cong B \otimes_{B_i} C_i$.
This is proved in Algebra, Lemma \ref{algebra-lemma-etale}.
\item In the situation of (1) show that
$\mathcal{G}(V) = \colim_{i' \geq i} \mathcal{G}_{i'}(V_{i'})$
where $V_{i'}$ is the base change of $V_i$ to $Y_{i'}$.
\item By (1), we see that for every \'etale covering
$\mathcal{V} = \{V_j \to Y\}_{j\in J}$ with $J$ finite and the
$V_j$s quasi-compact and quasi-separated, there exists $i \in I$ and
an \'etale covering $\mathcal{V}_i = \{V_{ij} \to Y_i\}_{j \in J}$
such that $\mathcal{V} \cong \mathcal{V}_i \times_{Y_i} Y$.
\item Show that (2) and (3) imply
$$
\check H^*(\mathcal{V}, \mathcal{G})=
\colim_{i\in I} \check H^*(\mathcal{V}_i, \mathcal{G}_i).
$$
\item Cleverly use the {\v C}ech-to-cohomology spectral sequence
(Theorem \ref{theorem-cech-ss}).
\end{enumerate}
\end{proof}

\begin{lemma}
\label{lemma-higher-direct-images}
Let $f: X\to Y$ be a morphism of schemes and $\mathcal{F}\in
\textit{Ab}(X_\etale)$. Then $R^pf_*\mathcal{F}$ is the sheaf
associated to the presheaf
$$
(V \to Y) \longmapsto H_\etale^p(X \times_Y V, \mathcal{F}|_{X \times_Y V}).
$$
More generally, for $K \in D(X_\etale)$ we have that $R^pf_*K$ is the
sheaf associated to the presheaf
$$
(V \to Y) \longmapsto H_\etale^p(X \times_Y V, K|_{X \times_Y V}).
$$
\end{lemma}

\begin{proof}
This lemma is valid for topological spaces, and the proof in this case is the
same. See Cohomology on Sites, Lemma
\ref{sites-cohomology-lemma-higher-direct-images}
for the case of a sheaf and see
Cohomology on Sites, Lemma
\ref{sites-cohomology-lemma-sheafification-cohomology}
for the case of a complex of abelian sheaves.
\end{proof}

\begin{lemma}
\label{lemma-relative-colimit}
Let $S$ be a scheme. Let $X = \lim_{i \in I} X_i$ be a limit of a
directed system of schemes over $S$ with affine transition morphisms
$f_{i'i} : X_{i'} \to X_i$. We assume the structure morphisms
$g_i : X_i \to S$ and $g : X \to S$ are quasi-compact and quasi-separated.
Let $(\mathcal{F}_i, \varphi_{i'i})$ be a system of abelian sheaves
on $(X_i, f_{i'i})$. Denote $f_i : X \to X_i$ the projection and set
$\mathcal{F} = \colim f_i^{-1}\mathcal{F}_i$. Then
$$
\colim_{i\in I} R^p g_{i, *} \mathcal{F}_i = R^p g_* \mathcal{F}
$$
for all $p \geq 0$.
\end{lemma}

\begin{proof}
Recall (Lemma \ref{lemma-higher-direct-images})
that $R^p g_{i, *} \mathcal{F}_i$ is the sheaf associated to the
presheaf $U \mapsto H^p_\etale(U \times_S X_i, \mathcal{F}_i)$
and similarly for $R^pg_*\mathcal{F}$. Moreover, the colimit of a
system of sheaves is the sheafification of the colimit on the level
of presheaves. Note that every object of $S_\etale$ has a covering
by quasi-compact and quasi-separated objects (e.g., affine schemes).
Moreover, if $U$ is a quasi-compact and quasi-separated object,
then we have
$$
\colim H^p_\etale(U \times_S X_i, \mathcal{F}_i) =
H^p_\etale(U \times_S X, \mathcal{F})
$$
by Theorem \ref{theorem-colimit}. Thus the lemma follows.
\end{proof}

\begin{lemma}
\label{lemma-relative-colimit-general}
Let $I$ be a directed set. Let $g_i : X_i \to S_i$ be an inverse system of
morphisms of schemes over $I$. Assume $g_i$ is quasi-compact and
quasi-separated and for $i' \geq i$ the transition morphisms
$f_{i'i} : X_{i'} \to X_i$ and $h_{i'i} : S_{i'} \to S_i$ are affine.
Let $g : X \to S$ be the limit of the morphisms $g_i$, see
Limits, Section \ref{limits-section-limits}.
Denote $f_i : X \to X_i$ and $h_i : S \to S_i$ the projections.
Let $(\mathcal{F}_i, \varphi_{i'i})$ be a system of sheaves
on $(X_i, f_{i'i})$. Set $\mathcal{F} = \colim f_i^{-1}\mathcal{F}_i$. Then
$$
R^p g_* \mathcal{F} =
\colim_{i \in I} h_i^{-1}R^p g_{i, *} \mathcal{F}_i
$$
for all $p \geq 0$.
\end{lemma}

\begin{proof}
How is the map of the lemma constructed?
For $i' \geq i$ we have a commutative diagram
$$
\xymatrix{
X \ar[r]_{f_{i'}} \ar[d]_g &
X_{i'} \ar[r]_{f_{i'i}} \ar[d]_{g_{i'}} &
X_i \ar[d]^{g_i} \\
S \ar[r]^{h_{i'}} &
S_{i'} \ar[r]^{h_{i'i}} &
S_i
}
$$
If we combine the base change map
$h_{i'i}^{-1}Rg_{i, *}\mathcal{F}_i \to Rg_{i', *}f_{i'i}^{-1}\mathcal{F}_i$
(Cohomology on Sites, Lemma
\ref{sites-cohomology-lemma-base-change-map-flat-case} or
Remark \ref{sites-cohomology-remark-base-change})
with the map $Rg_{i', *}\varphi_{i'i}$, then we obtain
$\psi_{i'i} : h_{i' i}^{-1} R^p g_{i, *} \mathcal{F}_i \to
R^pg_{i', *} \mathcal{F}_{i'}$. Similarly, using the left square
in the diagram we obtain maps
$\psi_i : h_i^{-1}R^pg_{i, *}\mathcal{F}_i \to R^pg_*\mathcal{F}$.
The maps $h_{i'}^{-1}\psi_{i'i}$ and $\psi_i$ are the maps used in
the statement of the lemma. For this to make sense, we have to check that
$\psi_{i''i} = \psi_{i''i'} \circ h_{i''i'}^{-1}\psi_{i'i}$ and
$\psi_{i'} \circ h_{i'}^{-1}\psi_{i'i} = \psi_i$; this follows
from Cohomology on Sites, Remark
\ref{sites-cohomology-remark-compose-base-change-horizont